\documentclass[11pt,openany]{book}

% -----------------------------------------------------------------------------
% Transseries for Mere Mortals
% A self-contained tutorial written for a reader who knows calculus, elementary
% differential equations, and ordinary power series.
% -----------------------------------------------------------------------------

\usepackage[T1]{fontenc}
\usepackage[utf8]{inputenc}
\usepackage{microtype}
\usepackage{geometry}
\geometry{
  paperwidth=7.25in,
  paperheight=10.25in,
  inner=0.86in,
  outer=0.72in,
  top=0.78in,
  bottom=0.84in,
  headsep=0.22in,
  footskip=0.38in
}
\usepackage{amsmath,amssymb,amsthm,mathtools}
\usepackage{newtxtext,newtxmath}
\usepackage{bm}
\usepackage{array,booktabs,longtable,tabularx,multirow}
\usepackage{enumitem}
\usepackage{xcolor}
\usepackage{tikz}
\usetikzlibrary{arrows.meta,calc,decorations.pathreplacing,positioning,patterns,shapes.geometric}
\usepackage[most]{tcolorbox}
\usepackage{listings}
\usepackage{fancyhdr}
\usepackage{titlesec}
\usepackage{imakeidx}
\usepackage{bookmark}
\usepackage{hyperref}
\usepackage[capitalise,nameinlink,noabbrev]{cleveref}
\usepackage{etoolbox}
\usepackage{needspace}

\definecolor{DeepBlue}{HTML}{173A5E}
\definecolor{MediumBlue}{HTML}{2F6690}
\definecolor{PaleBlue}{HTML}{EAF2F8}
\definecolor{DeepTeal}{HTML}{1D5C63}
\definecolor{PaleTeal}{HTML}{E9F5F3}
\definecolor{DeepGold}{HTML}{7A5B13}
\definecolor{PaleGold}{HTML}{FBF5E6}
\definecolor{DeepRed}{HTML}{7A2E2E}
\definecolor{PaleRed}{HTML}{FBECEC}
\definecolor{SoftGray}{HTML}{F4F5F6}
\definecolor{DarkGray}{HTML}{3D4650}

\hypersetup{
  colorlinks=true,
  linkcolor=DeepBlue,
  citecolor=DeepTeal,
  urlcolor=MediumBlue,
  pdftitle={Transseries for Mere Mortals},
  pdfauthor={OpenAI},
  pdfsubject={A step-by-step tutorial on logarithmic-exponential and resurgent transseries},
  pdfkeywords={transseries, asymptotics, Hahn series, resurgence, Borel summation, Stokes phenomenon}
}

\setcounter{tocdepth}{2}
\setcounter{secnumdepth}{3}
\makeindex[intoc,title=Index]

\pagestyle{fancy}
\fancyhf{}
\fancyhead[LE]{\small\scshape\nouppercase{\leftmark}}
\fancyhead[RO]{\small\scshape\nouppercase{\rightmark}}
\fancyfoot[LE,RO]{\thepage}
\renewcommand{\headrulewidth}{0.4pt}
\renewcommand{\footrulewidth}{0pt}

\titleformat{\chapter}[display]
  {\normalfont\bfseries\color{DeepBlue}}
  {\filleft\Large\chaptername\ \thechapter}
  {0.5ex}
  {\titlerule\vspace{1.0ex}\filright\Huge}
  [\vspace{0.6ex}\titlerule]
\titleformat{\section}{\Large\bfseries\color{DeepBlue}}{\thesection}{0.7em}{}
\titleformat{\subsection}{\large\bfseries\color{DeepTeal}}{\thesubsection}{0.7em}{}
\titleformat{\subsubsection}{\normalsize\bfseries\color{DarkGray}}{\thesubsubsection}{0.7em}{}

\setlist[itemize]{leftmargin=1.4em,itemsep=0.25ex,topsep=0.5ex}
\setlist[enumerate]{leftmargin=1.8em,itemsep=0.35ex,topsep=0.5ex}
\setlength{\parindent}{1.15em}
\setlength{\parskip}{0.15ex}
\setlength{\headheight}{14pt}
\raggedbottom
\allowdisplaybreaks[2]

\theoremstyle{definition}
\newtheorem{definition}{Definition}[chapter]
\newtheorem{example}{Example}[chapter]
\newtheorem{exercise}{Exercise}[chapter]
\newtheorem{algorithm}{Algorithm}[chapter]
\theoremstyle{plain}
\newtheorem{theorem}{Theorem}[chapter]
\newtheorem{proposition}{Proposition}[chapter]
\newtheorem{lemma}{Lemma}[chapter]
\newtheorem{corollary}{Corollary}[chapter]
\theoremstyle{remark}
\newtheorem{remark}{Remark}[chapter]

\newtcolorbox{intuitionbox}[1][]{
  enhanced,breakable,colback=PaleBlue,colframe=MediumBlue,
  boxrule=0.6pt,arc=1.2mm,left=1.5mm,right=1.5mm,top=1.1mm,bottom=1.1mm,
  title={Intuition},fonttitle=\bfseries,#1
}
\newtcolorbox{checkpointbox}[1][]{
  enhanced,breakable,colback=PaleTeal,colframe=DeepTeal,
  boxrule=0.6pt,arc=1.2mm,left=1.5mm,right=1.5mm,top=1.1mm,bottom=1.1mm,
  title={Checkpoint},fonttitle=\bfseries,#1
}
\newtcolorbox{warningbox}[1][]{
  enhanced,breakable,colback=PaleRed,colframe=DeepRed,
  boxrule=0.6pt,arc=1.2mm,left=1.5mm,right=1.5mm,top=1.1mm,bottom=1.1mm,
  title={A common trap},fonttitle=\bfseries,#1
}
\newtcolorbox{summarybox}[1][]{
  enhanced,breakable,colback=PaleGold,colframe=DeepGold,
  boxrule=0.6pt,arc=1.2mm,left=1.5mm,right=1.5mm,top=1.1mm,bottom=1.1mm,
  title={Chapter summary},fonttitle=\bfseries,#1
}
\newtcolorbox{workedbox}[1][]{
  enhanced,breakable,colback=SoftGray,colframe=DarkGray,
  boxrule=0.5pt,arc=1.2mm,left=1.5mm,right=1.5mm,top=1.1mm,bottom=1.1mm,
  title={Worked calculation},fonttitle=\bfseries,#1
}

\lstdefinestyle{pseudo}{
  basicstyle=\ttfamily\small,
  columns=fullflexible,
  frame=single,
  rulecolor=\color{DarkGray},
  backgroundcolor=\color{SoftGray},
  showstringspaces=false,
  keepspaces=true,
  xleftmargin=0.5em,
  xrightmargin=0.5em,
  aboveskip=0.8em,
  belowskip=0.8em
}

% Mathematical notation.
\newcommand{\R}{\mathbb{R}}
\newcommand{\C}{\mathbb{C}}
\newcommand{\N}{\mathbb{N}}
\newcommand{\Z}{\mathbb{Z}}
\newcommand{\Q}{\mathbb{Q}}
\newcommand{\T}{\mathbb{T}}
\newcommand{\M}{\mathfrak{M}}
\newcommand{\G}{\mathfrak{G}}
\newcommand{\m}{\mathfrak{m}}
\newcommand{\n}{\mathfrak{n}}
\newcommand{\one}{\mathbf{1}}
\newcommand{\eps}{\varepsilon}
\newcommand{\dd}{\,\mathrm{d}}
\newcommand{\e}{\mathrm{e}}
\newcommand{\ii}{\mathrm{i}}
\newcommand{\bigO}{\mathcal{O}}
\newcommand{\smallo}{\mathrm{o}}
\newcommand{\asym}{\sim}
\newcommand{\precdom}{\prec}
\newcommand{\preceqdom}{\preccurlyeq}
\newcommand{\succdom}{\succ}
\newcommand{\succeqdom}{\succcurlyeq}
\newcommand{\equivdom}{\asymp}
\newcommand{\formal}{\widehat{\vphantom{f}\smash{\phantom{f}}}}
\newcommand{\trans}[1]{\widehat{#1}}
\DeclareMathOperator{\supp}{supp}
\DeclareMathOperator{\lm}{lm}
\DeclareMathOperator{\lc}{lc}
\DeclareMathOperator{\magop}{mag}
\DeclareMathOperator{\ord}{ord}
\DeclareMathOperator{\Ei}{Ei}
\DeclareMathOperator{\Ai}{Ai}
\DeclareMathOperator{\Bi}{Bi}
\DeclareMathOperator{\Res}{Res}
\DeclareMathOperator{\id}{id}
\DeclareMathOperator{\Log}{Log}
\newcommand{\Borel}{\mathcal{B}}
\newcommand{\Laplace}{\mathcal{L}}
\newcommand{\Stokes}{\mathfrak{S}}
\newcommand{\alien}{\Delta}
\newcommand{\median}{\mathcal{S}_{\mathrm{med}}}
\newcommand{\lateralsum}[1]{\mathcal{S}_{#1}}

\crefname{definition}{definition}{definitions}
\crefname{example}{example}{examples}
\crefname{exercise}{exercise}{exercises}
\crefname{algorithm}{algorithm}{algorithms}
\crefname{theorem}{theorem}{theorems}
\crefname{proposition}{proposition}{propositions}
\crefname{lemma}{lemma}{lemmas}
\crefname{corollary}{corollary}{corollaries}

\begin{document}

\frontmatter

\begin{titlepage}
\thispagestyle{empty}
\begin{tikzpicture}[remember picture,overlay]
  \fill[PaleBlue] (current page.north west) rectangle (current page.south east);
  \fill[DeepBlue] (current page.north west) rectangle ([xshift=1.0in]current page.south west);
  \draw[MediumBlue,line width=1.1pt]
    ([xshift=1.33in,yshift=-1.25in]current page.north west)
    -- ([xshift=-0.55in,yshift=-1.25in]current page.north east);
  \foreach \y/\s in {-3.45/0.42,-4.35/0.31,-5.2/0.23,-6.0/0.17,-6.75/0.12} {
    \draw[DeepTeal,line width=1.0pt]
      ([xshift=1.45in,yshift=\y in]current page.north west)
      .. controls +(1.1in,0.45in) and +(-1.1in,-0.45in) ..
      ([xshift=-0.7in,yshift=\y in]current page.north east);
  }
  \node[anchor=north west,text width=5.15in,align=left]
    at ([xshift=1.45in,yshift=-1.72in]current page.north west) {
      {\fontsize{33}{37}\selectfont\bfseries\color{DeepBlue} Transseries\par}
      \vspace{0.15em}
      {\fontsize{24}{29}\selectfont\bfseries\color{DeepTeal} for Mere Mortals\par}
      \vspace{0.7em}
      {\Large\color{DarkGray} A step-by-step tutorial from ordinary\\ asymptotic series to exponentials beyond all orders\par}
    };
  \node[anchor=south west,text width=4.9in]
    at ([xshift=1.45in,yshift=0.95in]current page.south west) {
      {\large\color{DarkGray} Formal logarithmic-exponential transseries, differential equations, Borel summation, Stokes phenomena, resurgence, Hardy fields, and surreal numbers}\par
      \vspace{0.8em}
      {\normalsize\color{DeepBlue} Prepared as a self-contained mathematical tutorial}\par
    };
\end{tikzpicture}
\null
\end{titlepage}

\thispagestyle{empty}
\vspace*{\fill}
\begin{center}
\begin{minipage}{0.82\textwidth}
\small
This tutorial is an original synthesis based on the source materials supplied with the request and on the literature listed in the bibliography. Its principal expository sources include Edgar's introductory papers, van der Hoeven's treatment of transseries and real differential algebra, the asymptotic differential algebra of Aschenbrenner--van den Dries--van der Hoeven, and the resurgent-transseries literature of Costin, Dorigoni, and Aniceto--Ba\c{s}ar--Schiappa.

\medskip
The word \emph{transseries} is used in several overlapping traditions. This book states explicitly which formal or analytic structure is being used at each point. It also distinguishes formal identities from claims about actual functions.

\medskip
Compiled 31 August 2026.
\end{minipage}
\end{center}
\vspace*{\fill}
\cleardoublepage

\chapter*{Preface}
\addcontentsline{toc}{chapter}{Preface}

Power series are among the first genuinely magical objects in mathematics. A finite formula can be replaced by an infinite list of coefficients; differentiation and integration become mechanical; and local behavior becomes algebra. But ordinary power series are only the first rung of a much taller ladder.

Suppose that a function has an expansion in inverse powers of a large variable,
\[
 f(x)\asym \frac{a_0}{x}+\frac{a_1}{x^2}+\frac{a_2}{x^3}+\cdots,
 \qquad x\to+\infty.
\]
The expansion may diverge. Even when it is useful, it cannot detect a term such as $\e^{-x}$, because $\e^{-x}$ is smaller than every power $x^{-N}$. Two different functions can therefore have exactly the same power-series asymptotics. If a differential equation has a free constant multiplying $\e^{-x}$, a power series literally loses one degree of freedom.

A transseries repairs this loss by allowing a hierarchy of scales:
\[
 x^{-1},x^{-2},\ldots;
 \qquad \e^{-x},\e^{-x}x^{-1},\ldots;
 \qquad \e^{-2x},\ldots;
 \qquad \log x,\e^{\sqrt{x}},\e^{-\e^x},\ldots.
\]
The result is not merely a longer series. It is a formal language with an ordering, arithmetic, exponentials, logarithms, differentiation, integration, and composition. In analytic problems it is also the natural bookkeeping device for Borel summation, Stokes jumps, and resurgence.

This book has three goals.

\begin{enumerate}
\item Build the central ideas from calculus and ordinary asymptotic expansions, with no assumption that the reader already knows ordered fields, valuations, or resurgence.
\item Show enough of the formal construction that the algebra does not look like hand-waving with infinitely many symbols.
\item Work through actual calculations: factorial divergence, Lambert's $W$ function, formal integration, algebraic equations, Riccati equations, WKB, Borel transforms, and Stokes discontinuities.
\end{enumerate}

The intended reader is a mathematically mature undergraduate who has completed a standard calculus sequence and has seen elementary differential equations and power series. Complex analysis is helpful in the chapters on Borel summation, but the necessary contour ideas are reviewed informally.

\section*{How to read the book}

Chapters~\ref{ch:scales}--\ref{ch:first-transseries} are the shortest route to the main motivation. Chapters~\ref{ch:monomials}--\ref{ch:explog} build the formal field. Chapters~\ref{ch:differentiation}--\ref{ch:fixed-points} explain its calculus. Chapters~\ref{ch:dominant-balance}--\ref{ch:singular-perturbation} apply it to equations. Chapters~\ref{ch:borel}--\ref{ch:resurgence} explain analytic reconstruction and resurgence. The final part places transseries among Hardy fields, surreal numbers, and model theory.

A reader mainly interested in applications may read Chapters~\ref{ch:scales}--\ref{ch:first-transseries}, skim Chapters~\ref{ch:monomials}--\ref{ch:explog}, and continue at Chapter~\ref{ch:dominant-balance}. A reader interested in foundations should follow the order given.

\begin{warningbox}
A formal transseries is not automatically the expansion of a function, and an asymptotic expansion is not automatically convergent. Throughout the book, a hat such as $\trans y$ often emphasizes a formal object. Symbols without hats may denote actual functions when the analytic setting has been specified.
\end{warningbox}

\chapter*{Notation and standing conventions}
\addcontentsline{toc}{chapter}{Notation and standing conventions}

Unless stated otherwise, $x\to+\infty$. We compare positive quantities by eventual size. The notation
\[
 f\precdom g
 \quad\text{means}\quad
 \frac{f}{g}\longrightarrow 0,
\]
and $f\equivdom g$ means that $f/g$ tends to a nonzero finite constant or, in a purely formal setting, that $f$ and $g$ have the same leading monomial. The familiar notation $f\asym g$ is reserved for asymptotic equivalence $f/g\to1$ when confusion is possible.

For a transseries $T$, $\supp T$ is its support, $\lm(T)$ its leading monomial, and $\lc(T)$ the coefficient of that monomial. We write $T\precdom1$ for an infinitesimal transseries and $T\succdom1$ for an infinitely large one.

The real logarithm is used on positive real arguments. In complex-analytic chapters, $\Log$ denotes a chosen branch. Formal exponentials and logarithms are algebraic operations defined by canonical decompositions; they do not require numerical convergence.

\tableofcontents

\mainmatter

\part{Why ordinary series are not enough}

\chapter{Scales, orders, and asymptotic expansions}\label{ch:scales}

\section{An expansion needs a scale}

When we write a Taylor series near $x=0$, the scale is understood:
\[
 1,\ x,\ x^2,\ x^3,\ldots.
\]
For asymptotics at infinity it is often more natural to use
\[
 1,\ x^{-1},\ x^{-2},\ x^{-3},\ldots.
\]
The important feature is not that these objects are powers. It is that each one is much smaller than the preceding one.

\begin{definition}[Asymptotic scale]
A sequence $\phi_0,\phi_1,\phi_2,\ldots$ is an \emph{asymptotic scale} as $x\to+\infty$ if
\[
 \phi_{n+1}(x)\precdom \phi_n(x)
 \qquad(n=0,1,2,\ldots).
\]
Equivalently, $\phi_{n+1}/\phi_n\to0$.
\end{definition}

Examples include
\[
 1,\frac1x,\frac1{x^2},\ldots;
 \qquad
 x,\log x,\log\log x,1,\frac1{\log\log x},\ldots;
\]
and
\[
 1,\e^{-x},\e^{-2x},\e^{-3x},\ldots.
\]
The scales can be mixed. For instance,
\[
 1,\frac1x,\frac1{x^2},\ldots,\e^{-x},\frac{\e^{-x}}x,
 \frac{\e^{-x}}{x^2},\ldots,\e^{-2x},\ldots
\]
is ordered by decreasing size, although it is not naturally indexed by a single integer without some bookkeeping.

\begin{figure}[htbp]
\centering
\begin{tikzpicture}[x=1cm,y=1cm,>=Latex]
  \draw[->,thick] (0,0) -- (11.2,0) node[right] {smaller};
  \foreach \x/\lab in {
    0.6/{$x$},
    2.0/{$1$},
    3.3/{$1/x$},
    4.5/{$1/x^2$},
    6.0/{$\e^{-\sqrt{x}}$},
    7.6/{$\e^{-x}$},
    9.2/{$\e^{-x^2}$},
    10.5/{$\e^{-\e^x}$}}
  {
    \draw (\x,0.12)--(\x,-0.12);
    \node[rotate=35,anchor=west] at (\x,-0.22) {\lab};
  }
  \node[anchor=west,text=DeepTeal] at (0.2,1.0)
    {$x\succdom1\succdom x^{-1}\succdom x^{-2}\succdom\e^{-\sqrt{x}}\succdom\e^{-x}\succdom\e^{-x^2}\succdom\e^{-\e^x}$};
\end{tikzpicture}
\caption{A tiny fragment of the hierarchy of asymptotic scales. There is no last scale.}
\label{fig:scale-hierarchy}
\end{figure}

\section{Poincar\'e asymptotic expansions}

\begin{definition}[Poincar\'e expansion]
Let $\phi_0,\phi_1,\ldots$ be an asymptotic scale. We write
\[
 f(x)\asym \sum_{n=0}^{\infty}a_n\phi_n(x)
\]
if, for every fixed $N\ge0$,
\[
 f(x)-\sum_{n=0}^{N}a_n\phi_n(x)
   =\smallo\bigl(\phi_N(x)\bigr)
 \qquad(x\to+\infty).
\]
The infinite sum is notation for the whole sequence of finite remainder statements. It need not converge for any $x$.
\end{definition}

There are two ideas to absorb.

First, $N$ is fixed before $x\to\infty$. We are not asserting that an actual infinite numerical sum equals $f(x)$. Second, a later term improves the error relative to the preceding scale, even when the coefficient sequence grows so quickly that the series diverges.

\begin{example}
For $x>0$,
\[
 \frac1{x+1}=\frac1x\frac1{1+x^{-1}}.
\]
Expanding the geometric factor gives
\[
 \frac1{x+1}\asym \frac1x-\frac1{x^2}+\frac1{x^3}-\cdots.
\]
Here the series actually converges whenever $x>1$, but convergence is stronger than asymptoticity and is not part of the definition.
\end{example}

\begin{example}
Repeated integration by parts gives
\[
 \int_x^\infty \e^{-t}t^{-1/2}\dd t
 \asym \e^{-x}x^{-1/2}
 \left(1-\frac1{2x}+\frac{3}{4x^2}-\frac{15}{8x^3}+\cdots\right).
\]
This is an expansion on the mixed scale
\(
 \e^{-x}x^{-1/2},\e^{-x}x^{-3/2},\ldots
\).
The exponential factor is part of every monomial in the scale.
\end{example}

\section{How coefficients are extracted}

If the scale is fixed, the coefficients are determined one at a time. Suppose
\[
 f\asym a_0\phi_0+a_1\phi_1+a_2\phi_2+\cdots.
\]
Then
\[
 a_0=\lim_{x\to\infty}\frac{f(x)}{\phi_0(x)},
\]
provided the normalization of $\phi_0$ makes this limit meaningful. After subtracting the first term,
\[
 a_1=\lim_{x\to\infty}
 \frac{f(x)-a_0\phi_0(x)}{\phi_1(x)},
\]
and so on. This is the asymptotic analogue of extracting Taylor coefficients, but it uses successive limits rather than derivatives at a point.

\begin{proposition}[Uniqueness relative to a scale]
If $\phi_{n+1}\precdom\phi_n$ and
\[
 f\asym\sum_{n\ge0}a_n\phi_n
 \qquad\text{and}\qquad
 f\asym\sum_{n\ge0}b_n\phi_n,
\]
then $a_n=b_n$ for every $n$.
\end{proposition}

\begin{proof}
If the coefficient lists differ, let $k$ be the first index with $a_k\ne b_k$. Subtract the two remainder statements through index $k$. The difference is
\[
 (a_k-b_k)\phi_k=\smallo(\phi_k),
\]
which is impossible because $(a_k-b_k)\phi_k/\phi_k$ is a nonzero constant.
\end{proof}

\section{The invisible-function problem}

Uniqueness of coefficients does \emph{not} mean uniqueness of the function. For every integer $N\ge0$,
\[
 \e^{-x}=\smallo(x^{-N}).
\]
Therefore
\[
 0\asym 0+\frac0x+\frac0{x^2}+\cdots
 \quad\text{and}\quad
 \e^{-x}\asym 0+\frac0x+\frac0{x^2}+\cdots.
\]
More generally, if
\[
 f(x)\asym\sum_{n\ge0}a_nx^{-n},
\]
then for every constant $C$,
\[
 f(x)+C\e^{-x}\asym\sum_{n\ge0}a_nx^{-n}.
\]
No amount of inverse-power data distinguishes these functions.

\begin{intuitionbox}
An ordinary inverse-power expansion is like a measuring instrument with a finite family of resolutions. An exponentially small quantity falls below every one of those resolutions. Transseries enlarge the instrument by adding new scales rather than pretending the invisible quantity is zero.
\end{intuitionbox}

The same phenomenon repeats at every level. Relative to powers of $\e^{-x}$, the term $\e^{-x^2}$ is beyond all orders:
\[
 \e^{-x^2}=\smallo(\e^{-Nx})
 \qquad\text{for every fixed }N.
\]
Relative to powers of $\e^{-x^2}$, the term $\e^{-\e^x}$ is beyond all orders again. There is no final universal list of scales.

\section{Logarithms create another hierarchy}

For every fixed $a>0$ and $b\in\R$,
\[
 \log x\precdom x^a,
 \qquad
 (\log x)^b\precdom x^a,
 \qquad
 x^{-a}\precdom(\log x)^{-b}\quad(b>0).
\]
Thus a more refined region of the scale hierarchy is
\[
 x\succdom x^{1/2}\succdom (\log x)^{100}
 \succdom\log x\succdom\log\log x\succdom1
 \succdom\frac1{\log\log x}\succdom\frac1{\log x}
 \succdom\frac1x.
\]
The comparison can usually be made by taking logarithms. For positive $f,g$,
\[
 f\precdom g
 \quad\Longleftrightarrow\quad
 \log f-\log g\longrightarrow-\infty.
\]
For example,
\[
 \log\left(\frac{\e^{\sqrt{x}}}{x^{1000}}\right)
 =\sqrt{x}-1000\log x\longrightarrow+\infty,
\]
so $\e^{\sqrt{x}}\succdom x^{1000}$.

\section{A practical comparison algorithm}

For expressions built from positive powers, exponentials, and logarithms, the following informal procedure often works.

\begin{algorithm}[Comparing two positive scales]
To compare $f$ and $g$ as $x\to+\infty$:
\begin{enumerate}
\item Form the quotient $f/g$.
\item Simplify its logarithm.
\item If the logarithm tends to $+\infty$, then $f\succdom g$; if it tends to $-\infty$, then $f\precdom g$.
\item If it tends to a finite number, remove the corresponding constant factor and compare the next correction.
\end{enumerate}
\end{algorithm}

\begin{workedbox}
Compare
\[
 f(x)=\e^{x+\sqrt{x}}x^{-7}
 \qquad\text{and}\qquad
 g(x)=\e^x x^{100}.
\]
Their quotient is
\[
 \frac fg=\e^{\sqrt{x}}x^{-107},
\]
and
\[
 \log(f/g)=\sqrt{x}-107\log x\longrightarrow+\infty.
\]
Therefore $f\succdom g$. The apparently large factor $x^{107}$ loses to the smaller exponential $\e^{\sqrt{x}}$.
\end{workedbox}

\section{Exercises}

\begin{exercise}
Arrange the following in decreasing order as $x\to+\infty$:
\[
 x^{-10},\quad \e^{-\sqrt{x}},\quad (\log x)^{-3},\quad
 \e^{-x},\quad x^{-1/100},\quad \e^{-x/\log x}.
\]
\end{exercise}

\begin{exercise}
Show that if $f\asym\sum_{n\ge0}a_nx^{-n}$, then
$f+\e^{-\sqrt{x}}\asym\sum_{n\ge0}a_nx^{-n}$.
Explain why the same conclusion need not hold after adding $(\log x)^{-1}$.
\end{exercise}

\begin{exercise}
Suppose $f(x)=\sqrt{x^2+x}$. Obtain the first four nonzero terms of its expansion in descending powers of $x$ as $x\to+\infty$.
\end{exercise}

\begin{summarybox}
An asymptotic expansion is defined relative to an ordered scale. Its coefficients are unique, but the represented function is not: terms smaller than every monomial in the chosen scale are invisible. Exponentials and iterated logarithms force us to work with a much richer hierarchy than ordinary powers.
\end{summarybox}


\chapter{Factorial divergence and the missing constant}\label{ch:euler}

The simplest useful doorway into transseries is a first-order differential equation whose formal power expansion diverges and loses its arbitrary constant.

\section{The Euler equation}

Consider
\begin{equation}\label{eq:euler-ode}
 y'(x)+y(x)=\frac1x,
 \qquad x\to+\infty.
\end{equation}
We first look for a formal inverse-power series
\begin{equation}\label{eq:euler-ansatz}
 \trans y_0(x)=\sum_{n=0}^{\infty}a_nx^{-n-1}.
\end{equation}
Differentiating term by term gives
\[
 \trans y_0'
 =-\sum_{n=0}^{\infty}(n+1)a_nx^{-n-2}
 =-\sum_{n=1}^{\infty}n a_{n-1}x^{-n-1}.
\]
Substitution into \cref{eq:euler-ode} yields
\[
 a_0x^{-1}+\sum_{n=1}^{\infty}(a_n-na_{n-1})x^{-n-1}=x^{-1}.
\]
Hence
\[
 a_0=1,
 \qquad
 a_n=na_{n-1}\quad(n\ge1),
\]
so
\begin{equation}\label{eq:euler-factorial-series}
 \trans y_0(x)=\frac1x+\frac1{x^2}+\frac{2!}{x^3}
 +\frac{3!}{x^4}+\frac{4!}{x^5}+\cdots
 =\sum_{n=0}^{\infty}\frac{n!}{x^{n+1}}.
\end{equation}

\section{The series diverges for every fixed $x$}

The ratio of consecutive terms is
\[
 \frac{(n+1)!x^{-n-2}}{n!x^{-n-1}}=\frac{n+1}{x}.
\]
For fixed $x$, this ratio eventually exceeds $1$, so the terms do not even tend to zero. The series cannot converge.

That does not make it useless. For a large $x$, the first several terms decrease before the factorial growth wins. At $x=20$, for example, the term sizes decrease until $n$ is close to $19$ or $20$. Truncating near the smallest term gives an excellent approximation.

\begin{definition}[Optimal truncation]
For a divergent asymptotic series whose terms first decrease and then increase, \emph{optimal truncation} means stopping at or near the term of smallest magnitude.
\end{definition}

For \cref{eq:euler-factorial-series}, the smallest term occurs near $n=x-1$. Stirling's formula gives, for $n\approx x$,
\[
 \frac{n!}{x^{n+1}}
 \approx \frac{\sqrt{2\pi x}(x/\e)^x}{x^{x+1}}
 =\sqrt{\frac{2\pi}{x}}\,\e^{-x}.
\]
The best attainable error from the inverse-power series is therefore on an exponentially small scale. The divergence is telling us where the next scale lives.

\begin{intuitionbox}
Divergence is often information, not failure. The late terms of the power series grow in precisely the way needed to announce an omitted contribution of order $\e^{-x}$.
\end{intuitionbox}

\section{The exact family of solutions}

Multiplying \cref{eq:euler-ode} by the integrating factor $\e^x$ gives
\[
 (\e^xy)'=\frac{\e^x}{x}.
\]
On any interval not crossing $x=0$,
\begin{equation}\label{eq:euler-exact-family}
 y_C(x)=\e^{-x}\left(C+\int_{x_0}^{x}\frac{\e^t}{t}\dd t\right),
\end{equation}
where $C$ is arbitrary. Equivalently, after absorbing a constant into $C$,
\[
 y_C(x)=\e^{-x}\Ei(x)+C\e^{-x}.
\]
All members of this one-parameter family have the same formal expansion \cref{eq:euler-factorial-series}, because
\[
 C\e^{-x}=\smallo(x^{-N})
\]
for every $N$.

The ordinary power series has forgotten the constant of integration. The smallest correction needed to restore it is exactly $C\e^{-x}$.

\section{The first transseries solution}

We now write
\begin{equation}\label{eq:first-transseries}
 \trans y(x;C)
 =\sum_{n=0}^{\infty}\frac{n!}{x^{n+1}}+C\e^{-x}.
\end{equation}
This is a very small transseries: one divergent power-series sector plus one exponentially small sector. Nevertheless it already does three jobs that the power series alone cannot do:

\begin{enumerate}
\item It retains the arbitrary constant of the differential equation.
\item It records the natural scale of the optimally truncated error.
\item It anticipates the ambiguity that appears when the divergent series is Borel summed across its Stokes direction.
\end{enumerate}

Because \cref{eq:euler-ode} is linear, no powers $C^2\e^{-2x},C^3\e^{-3x},\ldots$ are generated. Nonlinear equations usually produce an entire tower of such sectors.

\section{Checking the formal identity}

A formal identity is checked coefficientwise. Differentiate \cref{eq:first-transseries}:
\[
 \frac{\dd}{\dd x}\left(C\e^{-x}\right)=-C\e^{-x},
\]
so this term cancels when $y'+y$ is formed. The factorial series telescopes:
\begin{align*}
 \left(\sum_{n\ge0}\frac{n!}{x^{n+1}}\right)'
 +\sum_{n\ge0}\frac{n!}{x^{n+1}}
 &=-\sum_{n\ge0}\frac{(n+1)!}{x^{n+2}}
   +\sum_{n\ge0}\frac{n!}{x^{n+1}}\\
 &=\frac1x.
\end{align*}
Thus the transseries solves the equation formally, with no analytic summation yet assumed.

\section{Remainders by integration by parts}

The factorial coefficients can also be derived from the exact integral. Formally integrate by parts:
\[
 \int^x\frac{\e^t}{t}\dd t
 =\frac{\e^x}{x}+\int^x\frac{\e^t}{t^2}\dd t.
\]
Repeating $N$ times gives
\[
 \int^x\frac{\e^t}{t}\dd t
 =\e^x\sum_{n=0}^{N-1}\frac{n!}{x^{n+1}}
  +N!\int^x\frac{\e^t}{t^{N+1}}\dd t
  +\text{constant}.
\]
After multiplication by $\e^{-x}$,
\[
 y_C(x)=\sum_{n=0}^{N-1}\frac{n!}{x^{n+1}}
 +N!\e^{-x}\int^x\frac{\e^t}{t^{N+1}}\dd t
 +C\e^{-x}.
\]
For fixed $N$, the integral remainder is $\bigO(x^{-N-1})$, as required. When $N$ grows with $x$, the same expression leads toward exponentially improved asymptotics, but that requires more care than ordinary Poincar\'e asymptotics.

\section{A numerical feel for optimal truncation}

At $x=10$, the terms $n!/10^{n+1}$ decrease through approximately $n=9$ and then grow. The partial sum through $n=9$ is
\[
 0.1+0.01+0.002+0.0006+\cdots+\frac{9!}{10^{10}},
\]
and its least term has size about $3.6\times10^{-5}$. The exponential scale is
\[
 \e^{-10}\approx4.54\times10^{-5}.
\]
The close agreement is not accidental; it is the Stirling estimate above.

\begin{warningbox}
For a divergent asymptotic series, adding more terms after the least term generally makes the numerical approximation worse. ``More terms'' is not the same as ``more accuracy.''
\end{warningbox}

\section{Exercises}

\begin{exercise}
For the equation
\[
 y'-y=\frac1x,
\]
find a formal inverse-power solution and the exponentially small or large homogeneous correction. Which correction is small as $x\to+\infty$?
\end{exercise}

\begin{exercise}
Consider $y'+ay=x^{-1}$ with $a>0$. Derive the formal coefficients and identify the missing homogeneous term.
\end{exercise}

\begin{exercise}
Use Stirling's formula to show more carefully that the least term of
$\sum n!x^{-n-1}$ is of order $x^{-1/2}\e^{-x}$.
\end{exercise}

\begin{summarybox}
The Euler equation produces a factorially divergent series. Its least term is exponentially small, and the exact differential equation has a free term $C\e^{-x}$ invisible to all inverse powers. The smallest useful transseries already combines a perturbative sector with a beyond-all-orders sector.
\end{summarybox}


\chapter{What a transseries is---and is not}\label{ch:first-transseries}

There is no single one-line definition that covers every object called a transseries in the literature. The word names a family of closely related formal expansions. Two traditions dominate this tutorial.

\section{The logarithmic-exponential viewpoint}

In the formal-algebraic tradition, a transseries is a generalized series whose monomials can involve powers, exponentials, and logarithms, arranged in a well-controlled dominance order. Typical examples look like
\begin{align*}
 T(x)={}&\e^{x^2+\sqrt{x}}
 \left(1+\frac{3}{x}+\frac{7\log x}{x^2}+\cdots\right)\\
 &+5x^{\pi}\e^{-x}
 \left(1-\frac1{\log x}+\cdots\right)
 +\e^{-\e^x}+\cdots.
\end{align*}
The aim is to form an ordered differential field closed under a large collection of operations:
\[
 +,\ - ,\ \times,\ /,\ \exp,\ \log,\ \frac{\dd}{\dd x},
 \ \int,\ \text{and suitable composition}.
\]
This is the setting developed from generalized power series and logarithmic-exponential series, and presented accessibly by Edgar \cite{edgar-beginners,edgar-composition} and systematically by van der Hoeven \cite{vdhoeven-book}.

\section{The resurgent viewpoint}

In singular perturbation theory and mathematical physics, one often writes a transseries in a small parameter $z$ as
\begin{equation}\label{eq:resurgent-template}
 \trans F(z;\sigma)
 =\sum_{k=0}^{\infty}\sigma^k
   \e^{-kA/z}z^{k\beta}
   \sum_{n=0}^{\infty}a_{k,n}z^n.
\end{equation}
Here:

\begin{itemize}
\item $k$ labels nonperturbative sectors;
\item $A$ is an action or singulant;
\item $\sigma$ is a transseries parameter, often encoding a boundary condition;
\item each inner power series may diverge;
\item Borel singularities of one sector are related to other sectors.
\end{itemize}

The formal expression is only half the story. One studies directional Borel sums, their jumps across Stokes rays, and the web of large-order relations called \emph{resurgence}. Costin's work gives rigorous summability results for broad classes of differential equations \cite{costin-1995,costin-1998,costin-book}; modern pedagogical accounts include Dorigoni \cite{dorigoni} and Aniceto--Ba\c{s}ar--Schiappa \cite{abs-primer}.

\section{How the viewpoints fit together}

Both viewpoints organize many asymptotic scales in one formal object. Both need a rule ensuring that sums and products are meaningful. Both distinguish a leading part from successively smaller corrections. But their emphases differ.

\begin{center}
\small
\begin{tabularx}{\textwidth}{@{}>{\bfseries}p{0.22\textwidth}XX@{}}
\toprule
 & Formal logarithmic-exponential field & Resurgent transseries in analysis \\
\midrule
Typical variable & $x\to+\infty$ & $z\to0$ in the complex plane \\
Typical monomials & $x^a(\log x)^b\e^L$ & $\e^{-kA/z}z^{\beta} (\log z)^m$ \\
Central structure & Ordered Hahn field, derivation, composition & Borel plane, Stokes automorphisms, alien derivatives \\
Main question & Which formal operations are well defined? & Which functions are reconstructed, and how do sums jump? \\
Convergence required? & No & Usually replaced by summability \\
\bottomrule
\end{tabularx}
\end{center}

A concrete transseries may belong to both worlds. For example,
\[
 \sum_{k\ge0}C^k\e^{-kx}\sum_{n\ge0}a_{k,n}x^{-n}
\]
is a logarithmic-exponential formal series in $x$ and also has the sector form used in resurgence.

\begin{warningbox}
Do not infer analytic facts from formal closure alone. A formal transseries can be differentiated and multiplied without being assigned an actual numerical value. Conversely, two different summation prescriptions may assign different functions to the same divergent formal sector along a Stokes direction.
\end{warningbox}

\section{A minimal mental model}

Before constructing the full field, it is useful to picture a transseries as an ordered dictionary
\[
 \text{monomial}\longmapsto\text{coefficient}.
\]
For example,
\[
 T=3x^2-7+4x^{-1}+5\e^{-x}+2x\e^{-2x}
\]
corresponds to
\[
 x^2\mapsto3,
 \quad1\mapsto-7,
 \quad x^{-1}\mapsto4,
 \quad\e^{-x}\mapsto5,
 \quad x\e^{-2x}\mapsto2.
\]
The monomials are ordered
\[
 x^2\succdom1\succdom x^{-1}\succdom\e^{-x}\succdom x\e^{-2x}.
\]
The leading monomial is $x^2$ and the leading coefficient is $3$.

What makes transseries nontrivial is that the dictionary may be infinite, the keys form a sophisticated multiplicative group, and operations such as $\exp(T)$ generate new keys. The support conditions in Chapters~\ref{ch:hahn} and \ref{ch:grids} prevent an infinite amount of algebra from accumulating at a single monomial.

\section{A toy nonlinear transseries}

Consider the logistic-type equation
\begin{equation}\label{eq:logistic-toy}
 y'=-y+y^2.
\end{equation}
Near the zero solution as $x\to+\infty$, the linearized equation is $y'=-y$, suggesting a leading correction $C\e^{-x}$. Write
\[
 y=\sum_{k=1}^{\infty}a_k C^k\e^{-kx}.
\]
Then
\[
 y'=-\sum_{k\ge1}k a_kC^k\e^{-kx},
\]
and
\[
 y^2=\sum_{k\ge2}
 \left(\sum_{i=1}^{k-1}a_i a_{k-i}\right)
 C^k\e^{-kx}.
\]
Equating the coefficient of $C^k\e^{-kx}$ gives
\[
 (1-k)a_k=\sum_{i=1}^{k-1}a_i a_{k-i}
 \qquad(k\ge2),
\]
while $a_1$ is a normalization; take $a_1=1$. The recurrence yields
\[
 a_2=-1,
 \quad a_3=1,
 \quad a_4=-1,
 \quad\ldots,
\]
so
\begin{equation}\label{eq:logistic-geometric}
 \trans y=C\e^{-x}-C^2\e^{-2x}+C^3\e^{-3x}-\cdots.
\end{equation}
In this unusually friendly example, the transseries converges geometrically when $|C\e^{-x}|<1$ and sums to
\begin{equation}\label{eq:logistic-closed}
 y(x)=\frac{C\e^{-x}}{1+C\e^{-x}},
\end{equation}
which indeed solves \cref{eq:logistic-toy}.

This example displays the standard nonlinear mechanism: one exponentially small seed generates all multiples of its exponential scale. In harder problems each sector carries its own divergent power series.

\section{What is normally excluded}

The standard logarithmic-exponential field is enormous, but it is not the class of all imaginable formal expressions. Typical exclusions include:

\begin{itemize}
\item supports that contain infinitely increasing monomials, such as $1+x+x^2+x^3+\cdots$ when $x\succdom1$;
\item uncontrolled infinite nesting, such as an infinite tower $\exp(\exp(\exp(\cdots)))$;
\item infinite descent through logarithmic depth without a support discipline;
\item arbitrary rearrangements of formal sums.
\end{itemize}

These are not aesthetic restrictions. Without them, the coefficient of a product monomial may receive infinitely many unrelated contributions, a ``largest term'' may not exist, and recursive definitions may fail to stabilize.

\section{The promise of the formal theory}

The chapters ahead will justify the following informal slogan.

\begin{checkpointbox}
A logarithmic-exponential transseries behaves like a power series whose powers have been replaced by a well-ordered universe of asymptotic monomials. The universe is designed so that leading terms exist and the familiar algebraic operations can be extended recursively.
\end{checkpointbox}

\section{Exercises}

\begin{exercise}
Verify directly that \cref{eq:logistic-closed} solves \cref{eq:logistic-toy}. Expand it as a geometric series and identify the domain in which the numerical series converges.
\end{exercise}

\begin{exercise}
For
\[
 T=\e^{x}+x^{10}\e^{\sqrt{x}}+x^{-1}+\e^{-x^2},
\]
arrange the displayed monomials in decreasing order and identify the leading monomial.
\end{exercise}

\begin{exercise}
Explain in words why a nonlinear term $y^3$ would generate sectors $\e^{-3x},\e^{-5x},\ldots$ from a leading seed $\e^{-x}$ if symmetry prevents even powers.
\end{exercise}

\begin{summarybox}
``Transseries'' names overlapping formal and analytic frameworks. Logarithmic-exponential transseries emphasize an ordered differential field of generalized monomials; resurgent transseries emphasize divergent sectors, Borel singularities, and Stokes jumps. Both extend asymptotic bookkeeping beyond ordinary powers.
\end{summarybox}

\part{The formal logarithmic-exponential field}

\chapter{Transmonomials and dominance}\label{ch:monomials}

Ordinary Laurent series use monomials $x^n$. Transseries replace this one-dimensional list by a multiplicative group containing monomials such as
\[
 x^a,\qquad (\log x)^b,\qquad \e^{x},\qquad
 \e^{x+\sqrt{x}},\qquad \e^{-x^2}(\log x)^7,
\]
and many more. Before allowing infinite sums, we must understand the keys in our ordered dictionary.

\section{Ordered multiplicative groups}

\begin{definition}[Ordered abelian group of monomials]
A \emph{monomial group} $\M$ is an abelian group written multiplicatively, equipped with a total order $<$ that is compatible with multiplication:
\[
 \m<\n\quad\Longrightarrow\quad \m\mathfrak p<\n\mathfrak p
 \qquad(\mathfrak p\in\M).
\]
Its identity is the constant monomial $1$.
\end{definition}

For asymptotics at $+\infty$, the order is chosen to match eventual growth. We write
\[
 \m\precdom\n
\]
when $\m$ is asymptotically smaller than $\n$. Thus
\[
 \e^{-x}\precdom x^{-100}\precdom1\precdom\log x\precdom x^{1/10}\precdom\e^{\sqrt{x}}\precdom\e^x.
\]
Multiplicative compatibility says that comparisons survive multiplication by a positive monomial. For example,
\[
 \e^{-x}\precdom x^{-100}
 \quad\Longrightarrow\quad
 x^7\e^{-x}\precdom x^{-93}.
\]

Every monomial has an inverse. If $\m\succdom1$, then $\m^{-1}\precdom1$. The group language packages this simple reciprocal relation.

\section{Power monomials as the first model}

Take
\[
 \M_0=\{x^a:a\in\R\}.
\]
Multiplication corresponds to addition of exponents:
\[
 x^a x^b=x^{a+b},
 \qquad
 (x^a)^{-1}=x^{-a}.
\]
The dominance order is
\[
 x^a\precdom x^b\quad\Longleftrightarrow\quad a<b.
\]
Thus the map $a\mapsto x^a$ converts the ordered additive group $(\R,+,<)$ into an ordered multiplicative group.

When exponentials are introduced, their exponents themselves become ordered objects. If $L_1-L_2\succdom1$ and is positive, then
\[
 \e^{L_1}\succdom\e^{L_2},
\]
because
\[
 \frac{\e^{L_1}}{\e^{L_2}}=\e^{L_1-L_2}\succdom1.
\]
This recursive use of order is the central architectural idea of logarithmic-exponential transseries.

\section{Coefficient times monomial}

A term has the form $a\m$, where $a\in\R\setminus\{0\}$ and $\m\in\M$. The coefficient does not change the monomial's dominance class. For instance,
\[
 10^{100}\e^{-x}\precdom10^{-100}x^{-100}
\]
as $x\to+\infty$, despite the enormous finite constants.

\begin{definition}[Leading monomial and leading coefficient]
For a nonzero formal sum $T$ whose support has a largest monomial, the \emph{leading monomial} $\lm(T)$ is that largest monomial. The \emph{leading coefficient} $\lc(T)$ is its coefficient. The product
\[
 \lc(T)\lm(T)
\]
is the leading term.
\end{definition}

For
\[
 T=-7\e^{x}+x^{100}\e^{\sqrt{x}}+3x^{-1}+\e^{-x},
\]
we have
\[
 \lm(T)=\e^x,
 \qquad
 \lc(T)=-7.
\]
The sign of $T$ is the sign of its leading coefficient. All smaller terms are incapable of overturning it formally, just as a lower-degree term cannot change the sign of a polynomial for sufficiently large positive $x$.

\section{Dominance, comparability, and asymptotic equivalence}

Three relations are useful.

\begin{definition}
For nonzero transseries $S,T$:
\begin{align*}
 S\precdom T
 &\quad\Longleftrightarrow\quad
 \lm(S)\precdom\lm(T),\\
 S\equivdom T
 &\quad\Longleftrightarrow\quad
 \lm(S)=\lm(T),\\
 S\asym T
 &\quad\Longleftrightarrow\quad
 \lm(S-T)\precdom\lm(T)
 \ \text{and}\ \lc(S)=\lc(T).
\end{align*}
The last relation means $S/T=1+$ an infinitesimal.
\end{definition}

For example,
\[
 x^2+3x\not\equivdom 7x^2-\e^x,
\]
because the second expression is dominated by $-\e^x$. On the other hand,
\[
 x^2+3x\equivdom 7x^2-x^{-10},
\]
since these two expressions have the same leading monomial, even though their
leading coefficients differ. Finally,
\[
 x^2+3x\asym x^2-100+\e^{-x},
\]
because both have the same leading term $x^2$. Exact coefficients and exact
support matter in every one of these comparisons.

\section{Large, constant, and small parts}

A transseries $T$ can be split according to whether its monomials are greater than, equal to, or less than $1$:
\begin{equation}\label{eq:additive-decomposition-preview}
 T=L+c+S,
\end{equation}
where
\begin{itemize}
\item $L$ is \emph{purely large}: every monomial in $L$ is $\succdom1$;
\item $c\in\R$ is the constant term;
\item $S$ is \emph{small}: every monomial in $S$ is $\precdom1$.
\end{itemize}
This decomposition will define the exponential in Chapter~\ref{ch:explog}.

For example,
\[
 T=\underbrace{x^2-3\log x}_{L}
   +\underbrace{7}_{c}
   +\underbrace{x^{-1}+\e^{-x}}_{S}.
\]
The adjective ``purely'' matters. A large positive transseries such as $x+1+x^{-1}$ is not purely large because it contains constant and small monomials.

\section{Logarithmic derivatives}

A monomial often has the form
\[
 \m=x^a\e^L.
\]
Its derivative is best expressed through its logarithmic derivative:
\begin{equation}\label{eq:log-derivative-monomial}
 \frac{\m'}{\m}=\frac{a}{x}+L'.
\end{equation}
Thus
\[
 \m'=\left(\frac{a}{x}+L'\right)\m.
\]
This simple identity is the engine behind formal differentiation and integration. Examples:
\begin{align*}
 (x^a)'&=\frac{a}{x}x^a,\\
 (\e^{x^2})'&=2x\e^{x^2},\\
 (x^3\e^{-\sqrt{x}})'&=
 \left(\frac3x-\frac1{2\sqrt{x}}\right)x^3\e^{-\sqrt{x}}.
\end{align*}

\section{A dominance calculation by layers}

Compare
\[
 \m_1=x^{17}\exp\!\left(x+\sqrt{x}+\frac1x\right),
 \qquad
 \m_2=x^{-200}\exp\!\left(x+2\log x\right).
\]
Take the logarithm of their quotient:
\begin{align*}
 \log(\m_1/\m_2)
 &=217\log x+\sqrt{x}+\frac1x-2\log x\\
 &=\sqrt{x}+215\log x+\frac1x.
\end{align*}
This is positive and infinitely large, so $\m_1\succdom\m_2$. Notice the layered comparison:

\begin{enumerate}
\item The leading $x$ in the exponential cancels.
\item The next exponential contribution $\sqrt{x}$ beats every logarithmic contribution.
\item Only if that also canceled would the powers of $x$ decide the comparison.
\end{enumerate}

\begin{intuitionbox}
A transmonomial has something like a lexicographic address. Compare the highest exponential layer first; if it ties, move down to the next layer; only then compare powers and constants. The formal construction makes this intuition exact.
\end{intuitionbox}

\section{Exercises}

\begin{exercise}
Find the leading monomial and leading coefficient of
\[
 3x^5\e^{-x}+7\e^{-x/2}-2x^{100}\e^{-x}+x^{-1}.
\]
\end{exercise}

\begin{exercise}
Compare
\[
 \exp(x+\log^2x)
 \quad\text{and}\quad
 x^{1000}\e^x.
\]
Then compare
\(
 \exp(x+\sqrt{\log x})
\)
with the same second monomial.
\end{exercise}

\begin{exercise}
Compute the logarithmic derivative of
\[
 \m=x^{-3}(\log x)^7\exp\!\left(-x^2+\e^{\sqrt{x}}\right).
\]
\end{exercise}

\begin{summarybox}
Transmonomials form an ordered multiplicative group. A term is a real coefficient times a monomial; a transseries is ordered by its largest supported monomial. Dominance is recursive because exponentials are compared by comparing their exponents.
\end{summarybox}


\chapter{Hahn series: infinite sums that still behave}\label{ch:hahn}

If we allow arbitrary sets of monomials in a formal sum, addition is harmless but multiplication can become meaningless. Hahn series provide the right support condition.

\section{Why support matters}

Consider a formal sum
\[
 T=\sum_{\m\in\M}a_{\m}\m.
\]
Its \emph{support} is
\[
 \supp T=\{\m\in\M:a_{\m}\ne0\}.
\]
For an ordinary Laurent series in $x^{-1}$,
\[
 T=a_0+a_1x^{-1}+a_2x^{-2}+\cdots,
\]
the support descends
\[
 1\succdom x^{-1}\succdom x^{-2}\succdom\cdots.
\]
There is a largest term, and after removing any finite initial segment there is still a largest remaining term.

By contrast,
\[
 x+x^2+x^3+\cdots
\]
has no largest monomial as $x\to+\infty$. It does not describe a series proceeding from dominant terms to smaller corrections.

\begin{definition}[Reverse well-ordered support]
A subset $A\subseteq\M$ is \emph{reverse well ordered} if every nonempty subset of $A$ has a largest element. A formal sum with reverse well-ordered support is called a \emph{Hahn series} or a \emph{well-based generalized series}.
\end{definition}

Some authors order monomials in the opposite direction and say simply ``well ordered.'' The content is the same: the support must move one way, from larger monomials toward smaller ones, without an infinite climb back upward.

\section{Examples and nonexamples}

\begin{example}
The support
\[
 \{1,x^{-1},x^{-2},\ldots\}
\]
is reverse well ordered. Every nonempty subset corresponds to a nonempty set of nonnegative integers, which has a least index; that index gives the largest monomial.
\end{example}

\begin{example}
The support
\[
 \{x^{1/n}:n=1,2,3,\ldots\}
\]
is reverse well ordered. The whole set has largest element $x$, and every subset has a least denominator index. The monomials approach $1$ from above, but this is allowed.
\end{example}

\begin{example}
The support
\[
 \{x^{-1/n}:n=1,2,3,\ldots\}
\]
is not reverse well ordered. It climbs
\[
 x^{-1}\precdom x^{-1/2}\precdom x^{-1/3}\precdom\cdots
\]
toward $1$ without attaining a largest member.
\end{example}

\begin{example}
The support
\[
 \{\e^{-n x}:n=0,1,2,\ldots\}
\]
is reverse well ordered, so
\[
 \sum_{n\ge0}a_n\e^{-nx}
\]
is a valid Hahn series in the monomial group generated by $\e^{-x}$.
\end{example}

\section{The Hahn field}

Given an ordered abelian group $\M$ and a coefficient field $K$, write
\[
 K[[\M]]
\]
for the set of all formal sums with coefficients in $K$ and reverse well-ordered support. Addition is coefficientwise:
\[
 \left(\sum a_{\m}\m\right)+
 \left(\sum b_{\m}\m\right)
 =\sum(a_{\m}+b_{\m})\m.
\]
Multiplication is the Cauchy product:
\begin{equation}\label{eq:hahn-product}
 \left(\sum a_{\m}\m\right)
 \left(\sum b_{\n}\n\right)
 =\sum_{\mathfrak p\in\M}
 \left(\sum_{\m\n=\mathfrak p}a_{\m}b_{\n}\right)\mathfrak p.
\end{equation}
The inner coefficient sum must be finite, and the resulting support must again be reverse well ordered. A classical ordered-group lemma ensures both facts.

\begin{lemma}[Neumann's lemma, informal form]\label{lem:neumann}
Products and finite sums of reverse well-ordered supports remain reverse well ordered. Moreover, for supports satisfying the Hahn condition, a fixed product monomial receives only finitely many contributions in the Cauchy product.
\end{lemma}

A full proof belongs to ordered-group theory, but the basic mechanism is familiar from power series. In
\[
 (1+a_1t+a_2t^2+\cdots)(1+b_1t+b_2t^2+\cdots),
\]
the coefficient of $t^N$ receives only the $N+1$ contributions with indices $i+j=N$. The support condition generalizes this finite-decomposition property.

\section{A worked Hahn product}

Let $\mu=x^{-1}$ and $\nu=\e^{-x}$. Consider
\[
 A=1+2\mu+3\nu+5\mu\nu,
 \qquad
 B=1-\mu+4\nu.
\]
Then
\begin{align*}
 AB={}&1+\mu+(3+4)\nu
 -2\mu^2 +(5+8-3)\mu\nu\\
 &+12\nu^2-5\mu^2\nu+20\mu\nu^2\\
 ={}&1+\mu+7\nu-2\mu^2+10\mu\nu
 +12\nu^2-5\mu^2\nu+20\mu\nu^2.
\end{align*}
To place this in dominance order, note that every power of $x^{-1}$ is larger than $\e^{-x}$. Thus the beginning is actually
\[
 1+\mu-2\mu^2+\cdots+7\nu+10\mu\nu-5\mu^2\nu+\cdots.
\]
A finite display can hide a long block of monomials between two terms. The formal support, not visual adjacency, determines the order.

\section{Formal summability of families}

Transseries operations often require summing a family $(T_i)_{i\in I}$ rather than a single sequence of scalar terms.

\begin{definition}[Summable family]
A family $(T_i)_{i\in I}$ of Hahn series is \emph{formally summable} if:
\begin{enumerate}
\item the union $\bigcup_i\supp T_i$ is reverse well ordered;
\item each monomial belongs to $\supp T_i$ for only finitely many indices $i$.
\end{enumerate}
Then $\sum_iT_i$ is defined coefficientwise.
\end{definition}

The second condition prevents the coefficient of one monomial from becoming an undefined infinite scalar sum. In a purely formal field we do not want to decide whether
\(
 1-1+1-1+\cdots
\)
converges; we arrange that no such coefficient sum occurs.

\section{Geometric inversion}

Suppose $S\precdom1$. Then the family
\[
 1,-S,S^2,-S^3,\ldots
\]
is formally summable under the standard transseries support conditions, and
\begin{equation}\label{eq:formal-geometric}
 \frac1{1+S}=1-S+S^2-S^3+\cdots.
\end{equation}
The proof is purely algebraic:
\[
 (1+S)\sum_{n=0}^{N}(-S)^n=1+(-1)^NS^{N+1}.
\]
At every fixed monomial, the remainder eventually becomes too small to contribute, so the formal limit is $1$.

\begin{workedbox}
Invert
\[
 T=x^3+2x+1.
\]
Factor the leading term:
\[
 T=x^3(1+2x^{-2}+x^{-3}).
\]
Set $S=2x^{-2}+x^{-3}\precdom1$. Then
\begin{align*}
 T^{-1}
 &=x^{-3}(1-S+S^2-S^3+\cdots)\\
 &=x^{-3}-2x^{-5}-x^{-6}+4x^{-7}+4x^{-8}
 -7x^{-9}+\cdots.
\end{align*}
No numerical convergence is invoked. This is an identity in the Hahn field of formal Laurent series.
\end{workedbox}

\section{Truncation}

A Hahn series has meaningful initial segments. If $\n\in\M$, define the truncation above $\n$ by
\[
 T_{\succdom\n}=\sum_{\m\succdom\n}a_{\m}\m.
\]
Depending on the support, this truncation may be finite or infinite. Formal algorithms are often interpreted as stabilization of every truncation: once enough iterations have been performed, all coefficients above any chosen threshold stop changing.

This topology of coefficient stabilization is different from numerical convergence. It is closer to the $t$-adic topology of formal power series, where a sequence converges when more and more low-degree coefficients agree.

\section{Exercises}

\begin{exercise}
Classify each support as reverse well ordered or not:
\[
 \{x^{-n^2}:n\ge0\},\qquad
 \{x^{1-1/n}:n\ge2\},\qquad
 \{\e^{-x/n}:n\ge1\},\qquad
 \{\e^{-nx}:n\ge0\}.
\]
\end{exercise}

\begin{exercise}
Using formal geometric inversion, compute $1/(1+x^{-1}+x^{-3})$ through order $x^{-6}$.
\end{exercise}

\begin{exercise}
Explain why the family $T_n=x^{-n}+x^{-1}$ is not formally summable over $n\ge1$, even though the union of its supports is reverse well ordered.
\end{exercise}

\begin{summarybox}
Hahn series are generalized series with reverse well-ordered support. This condition gives a largest term and makes Cauchy products and coefficientwise infinite sums well defined. Formal convergence means stabilization by dominance, not convergence of numerical values.
\end{summarybox}


\chapter{Grids, ratio sets, and multidimensional bookkeeping}\label{ch:grids}

The full class of well-based supports is elegant but can be difficult to compute with. A useful and historically important subfield consists of \emph{grid-based} transseries, whose monomials lie in a finitely generated multiplicative lattice.

\section{Ratio monomials}

Choose finitely many small monomials
\[
 \boldsymbol\mu=(\mu_1,\ldots,\mu_k),
 \qquad \mu_i\precdom1.
\]
A product
\[
 \boldsymbol\mu^{\mathbf n}
 =\mu_1^{n_1}\cdots\mu_k^{n_k},
 \qquad \mathbf n=(n_1,\ldots,n_k)\in\Z^k,
\]
is a lattice monomial. The tuple $\boldsymbol\mu$ is called a \emph{ratio set} because each generator is a small ratio between neighboring scales.

\begin{definition}[Grid]
Given a base monomial $\m_0$ and a lower multi-index
$\mathbf n_0\in\Z^k$, the grid generated by $\boldsymbol\mu$ is
\[
 \mathfrak J(\m_0,\boldsymbol\mu,\mathbf n_0)
 =\left\{\m_0\boldsymbol\mu^{\mathbf n}:
 \mathbf n\in\Z^k,\ \mathbf n\ge\mathbf n_0\right\},
\]
where the inequality is componentwise.
A series is \emph{grid based} if its support is contained in some such grid.
\end{definition}

The lower bound permits finitely many negative powers of the generators but forces the support eventually to move in the small directions.

\section{The basic two-generator grid}

Let
\[
 \mu=x^{-1},
 \qquad
 \nu=\e^{-x}.
\]
Then the grid consists of
\[
 x^{-m}\e^{-nx},
 \qquad m,n\in\Z
\]
above chosen lower bounds. A typical transseries is
\begin{equation}\label{eq:two-grid-series}
 T=\sum_{n=0}^{\infty}\e^{-nx}
   \sum_{m=m_0(n)}^{\infty}a_{n,m}x^{-m}.
\end{equation}
Each fixed $n$ is an exponential sector; inside it sits an inverse-power series. The grid coordinates $(m,n)$ make multiplication a discrete convolution.

\begin{figure}[htbp]
\centering
\begin{tikzpicture}[x=0.85cm,y=0.68cm,>=Latex]
  \draw[->] (-0.3,0) -- (7.4,0) node[right] {$m$ (powers of $x^{-1}$)};
  \draw[->] (0,-0.3) -- (0,5.7) node[above] {$n$ (powers of $\e^{-x}$)};
  \foreach \i in {0,...,7}{\draw[gray!35] (\i,0)--(\i,5.2);}
  \foreach \j in {0,...,5}{\draw[gray!35] (0,\j)--(7.0,\j);}
  \foreach \i in {0,...,7}{\foreach \j in {0,...,5}{\fill[MediumBlue] (\i,\j) circle (1.25pt);}}
  \draw[thick,DeepTeal,->] (1,1)--(2,1) node[midway,below] {$\times x^{-1}$};
  \draw[thick,DeepTeal,->] (1,1)--(1,2) node[midway,left] {$\times\e^{-x}$};
  \node[anchor=west] at (4.6,4.7) {$x^{-m}\e^{-nx}$};
\end{tikzpicture}
\caption{A two-dimensional grid generated by $x^{-1}$ and $\e^{-x}$.}
\label{fig:two-dimensional-grid}
\end{figure}

Dominance is not simply lexicographic in $(m,n)$, but in this example the exponential coordinate decides first:
\[
 x^{-m}\e^{-nx}\succdom x^{-p}\e^{-qx}
 \quad\text{if }n<q,
\]
regardless of finite $m,p$. If $n=q$, then the smaller inverse-power exponent is larger.

\section{Why grids help algorithms}

On a finite-dimensional grid:

\begin{itemize}
\item monomials have integer coordinate vectors;
\item multiplication adds vectors;
\item truncation becomes a region in $\Z^k$;
\item fixed-point iteration can be analyzed by showing that every correction moves in a positive combination of small directions;
\item a computer algebra system can store a finite truncation in a sparse map keyed by integer vectors.
\end{itemize}

This is the transseries analogue of representing a multivariate power series by exponent tuples.

\section{A bivariate geometric expansion}

Let $S=2\mu-3\nu$, with $\mu=x^{-1}$ and $\nu=\e^{-x}$. Then
\[
 \frac1{1+S}=1-S+S^2-S^3+\cdots.
\]
Expanding by grid degree gives
\begin{align*}
 \frac1{1+2\mu-3\nu}
 ={}&1-2\mu+3\nu+4\mu^2-12\mu\nu+9\nu^2\\
 &-8\mu^3+36\mu^2\nu-54\mu\nu^2+27\nu^3+\cdots.
\end{align*}
Every monomial corresponds to a lattice point. The coefficient of $\mu^m\nu^n$ is obtained from finitely many multinomial contributions.

\section{Formal convergence on a grid}

A sequence $T_0,T_1,T_2,\ldots$ converges formally if, for every monomial threshold, its truncation above that threshold eventually stabilizes. On a grid generated by small ratios, a sufficient pattern is
\[
 T_{j+1}-T_j\in \mu_1^{r_{1,j}}\cdots\mu_k^{r_{k,j}}K[[\M]],
\]
where the correction vectors move indefinitely into the positive cone.

For the ordinary iteration
\[
 S_{j+1}=A+S_j^2,
 \qquad A\precdom1,
\]
starting at $S_0=0$, we get
\[
 S_1=A,
 \quad S_2=A+A^2,
 \quad S_3=A+A^2+2A^3+A^4,
 \quad\ldots.
\]
The coefficient of any fixed power $A^N$ eventually stops changing. This is formal convergence even if $A$ is merely a symbol.

\section{Grid based is not the same as well based}

Every grid-based support used in the standard construction is well based, but not every well-based support fits inside a finitely generated grid. For example,
\[
 \{x^{1/n}:n\ge1\}
\]
is reverse well ordered, yet the exponents $1/n$ do not lie in a finitely generated additive semigroup of $\R$ with a common positive step. Thus
\[
 \sum_{n\ge1}a_nx^{1/n}
\]
may be a valid well-based Hahn series while not being grid based.

The larger well-based field has cleaner closure properties and greater generality. The grid-based field has more explicit combinatorics. Edgar's expositions carefully track both settings \cite{edgar-beginners,edgar-composition}.

\section{Ratio sets as proof certificates}

In practical transseries algebra, a ratio set can do more than generate support. It can certify comparisons. To prove $A\precdom B$, one may show that every monomial in $A/B$ is a product of positive powers of the small generators. Similarly, an operator $\Phi$ is contractive if differences $\Phi(S)-\Phi(T)$ gain at least one small ratio relative to $S-T$.

This finite certificate is useful because direct comparison of complicated nested exponentials can be cumbersome. The ratio set records the handful of scale reductions that the calculation actually uses.

\begin{workedbox}
Suppose
\[
 \Phi(S)=a\mu+b\nu+c\mu S+dS^2,
 \qquad \mu,\nu\precdom1.
\]
Then
\[
 \Phi(S)-\Phi(T)
 =c\mu(S-T)+d(S-T)(S+T).
\]
If $S,T\precdom1$, both multipliers $c\mu$ and $d(S+T)$ are small. Hence the difference gains a small factor. Iteration therefore pushes errors to smaller and smaller grid monomials.
\end{workedbox}

\section{Exercises}

\begin{exercise}
Let $\mu=x^{-1}$ and $\nu=\e^{-x}$. Expand
\[
 (1-\mu-\nu)^{-1}
\]
through total grid degree $m+n\le3$.
\end{exercise}

\begin{exercise}
Show that the support $\{x^{-m}\e^{-nx}:m,n\ge0\}$ is reverse well ordered under asymptotic dominance.
\end{exercise}

\begin{exercise}
Explain why the support $\{x^{-1/n}:n\ge1\}$ cannot be contained in a grid generated by finitely many power monomials $x^{-a_1},\ldots,x^{-a_k}$ with $a_i>0$.
\end{exercise}

\begin{summarybox}
A ratio set is a finite list of small monomials. Its integer products form a multidimensional grid that turns support manipulation into lattice bookkeeping. Grid-based transseries are computationally convenient but form a proper subcollection of all well-based transseries.
\end{summarybox}

\chapter{Building logarithmic-exponential transseries by height and depth}\label{ch:construction}

We now describe the inductive construction behind the standard field $\T$ of logarithmic-exponential transseries. The construction is easier to understand if one keeps two counters:

\begin{itemize}
\item \emph{exponential height}: how many genuinely large exponential layers occur;
\item \emph{logarithmic depth}: how many iterated logarithms of the main variable are needed.
\end{itemize}

The details differ slightly among authors and between grid-based and well-based variants. The following version captures the common architecture used in Edgar's exposition and in the modern differential-algebraic literature \cite{edgar-beginners,vdhoeven-book,adh-book}.

\section{Height zero: generalized power series}

Begin with the monomial group
\[
 \G_0=\{x^b:b\in\R\}.
\]
The height-zero field is the Hahn field
\begin{equation}\label{eq:T0}
 \T_0=\R[[\G_0]].
\end{equation}
Its elements are generalized power series
\[
 T=\sum_{b\in A}a_bx^b,
\]
where $A\subseteq\R$ is reverse well ordered in the usual order of exponents.

This already goes beyond ordinary Laurent series. Exponents may be real and need not lie in an arithmetic progression. Examples include
\[
 x^{\sqrt2}+3x^{1/7}-5+x^{-\pi}+\cdots
\]
and the well-based but nongrid series
\[
 \sum_{n\ge1}2^{-n}x^{1/n}.
\]

\section{Purely large exponents become new monomials}

Suppose $\T_n$ has been constructed. Let $L\in\T_n$ be purely large. We introduce a new monomial $\e^L$. To retain ordinary power factors, define
\begin{equation}\label{eq:G-next}
 \G_{n+1}
 =\left\{x^b\e^L:
 b\in\R,\ L\in\T_n\text{ purely large}\right\}.
\end{equation}
We allow $L=0$, so $\G_n$ embeds in later levels. The order is determined by the sign of the leading term of the logarithmic quotient:
\[
 \log\!\left(\frac{x^a\e^A}{x^b\e^B}\right)
 =A-B+(a-b)\log x.
\]
If this transseries is negative and infinitely large, then
$x^a\e^A\precdom x^b\e^B$; if it is positive and infinitely large, the
inequality is reversed; if it is zero, the monomials coincide. More plainly,
compare $A$ and $B$ first, and if they cancel to logarithmic size, compare the
powers of $x$.

The next field is
\begin{equation}\label{eq:T-next}
 \T_{n+1}=\R[[\G_{n+1}]].
\end{equation}
Thus one alternates:
\[
 \begin{gathered}
  \boxed{\text{form well-based sums}}\\[-0.15ex]
  \Downarrow\\[-0.15ex]
  \boxed{\text{exponentiate purely large sums}}\\[-0.15ex]
  \Downarrow\\[-0.15ex]
  \boxed{\text{form well-based sums again}}.
 \end{gathered}
\]

\section{Examples of height}

\begin{center}
\small
\begin{tabularx}{\textwidth}{@{}p{0.35\textwidth}p{0.14\textwidth}X@{}}
\toprule
Expression & Height & Reason \\
\midrule
$x^{\pi}+x^{-1}$ & $0$ & generalized powers only \\
$\e^{x^2}+\e^{-x}$ & $1$ & exponentials of height-zero large exponents \\
$\e^{\e^x}$ & $2$ & exponential of a height-one purely large term \\
$\exp(\e^{\e^x}+x)$ & $3$ & outer exponential raises the height again \\
$\e^x(1+\e^{-x}+\e^{-2x}+\cdots)$ & $1$ & the parenthesized part is a series on a height-one grid \\
\bottomrule
\end{tabularx}
\end{center}

An expression's visual nesting can overstate its genuine height. For example,
\[
 \exp(x+\e^{-x})
 =\e^x\exp(\e^{-x})
 =\e^x\left(1+\e^{-x}+\frac12\e^{-2x}+\cdots\right).
\]
The small exponential in the exponent is expanded as a power series; it does not create a new large exponential monomial. The resulting transseries has height $1$, not $2$.

\section{The log-free union}

The union
\begin{equation}\label{eq:logfree-T}
 \T_{\mathrm{lf}}=\bigcup_{n\ge0}\T_n
\end{equation}
is called the \emph{log-free} transseries field. Every one of its elements has finite exponential height, although there is no uniform bound on the heights of all elements.

The phrase ``finite height'' is crucial. The formal infinite tower
\[
 \exp\!\bigl(\exp\!\bigl(\exp(\cdots)\bigr)\bigr)
\]
is not an element of the standard field. Enlarged fields with superexponentials or infinite-height structures exist, but they require extra construction; Schmeling's thesis is a foundational source in that direction \cite{schmeling}.

\section{Adding logarithmic depth}

The log-free field contains $x$ and exponentials built from $x$, but not yet every desired iterated logarithm as a basic scale. Let
\[
 \log_0x=x,
 \qquad
 \log_1x=\log x,
 \qquad
 \log_2x=\log\log x,
 \quad\ldots.
\]
For $m\ge0$, compose log-free transseries with $\log_mx$:
\[
 \T_{n,m}=\{T\circ\log_m:T\in\T_n\}.
\]
Then form the union over finite height and finite depth:
\begin{equation}\label{eq:full-T}
 \T=\bigcup_{m\ge0}\bigcup_{n\ge0}\T_{n,m}.
\end{equation}
The precise nesting of these unions is handled carefully in formal treatments, but the practical meaning is simple: any individual transseries uses only finitely many iterated logarithms and finitely many large exponential layers.

Examples:
\begin{align*}
 \log x+\frac1{\log x}+\frac1{(\log x)^2}+\cdots
 &\quad\text{has depth }1,\\
 \e^{\sqrt{\log x}}(\log\log x)^3
 &\quad\text{uses depth }2,\\
 \exp\!\bigl(\exp(\sqrt{\log\log x})\bigr)
 &\quad\text{has finite height and depth }2.
\end{align*}

\section{A tree view of monomials}

A transmonomial can be viewed as a finite syntax tree:

\begin{figure}[htbp]
\centering
\begin{tikzpicture}[
 level distance=1.0cm,
 level 1/.style={sibling distance=3.3cm},
 level 2/.style={sibling distance=1.7cm},
 every node/.style={draw,rounded corners,fill=PaleBlue,inner sep=3pt},
 edge from parent/.style={draw,-Latex}]
\node {$x^{7}\exp(L)$}
 child {node {$x^7$}}
 child {node {$L=x^2+\e^{\sqrt{x}}$}
   child {node {$x^2$}}
   child {node {$\e^{\sqrt{x}}$}
     child {node {$\sqrt{x}$}}}}
;
\end{tikzpicture}
\caption{A finite-height monomial as a syntax tree. Supports then form well-based forests of such trees.}
\label{fig:monomial-tree}
\end{figure}

The construction is recursive on these trees. Ordering two outer exponentials requires ordering their exponent subtrees. Differentiation follows the chain rule down the tree. Exponentiating a purely large series wraps it in a new node.

\section{Why the exponent must be purely large}

Any transseries $A$ has a unique additive decomposition
\[
 A=L+c+S
\]
into purely large, constant, and small parts. Then
\[
 \e^A=\e^c\e^L\e^S.
\]
Only $\e^L$ needs to be declared a new monomial. The constant $\e^c$ is a coefficient, and $\e^S$ expands as
\[
 1+S+\frac{S^2}{2!}+\cdots.
\]
If arbitrary small and constant pieces were hidden inside monomials, the same transseries would have many incompatible monomial representations. The purely large rule gives a canonical normal form.

\section{Finite complexity versus infinite support}

A transseries may have infinitely many terms while each monomial has finite syntactic complexity. This is analogous to an ordinary power series: the series has infinitely many monomials $x^n$, but each monomial is a finite expression.

The standard field permits:
\[
 \sum_{n\ge0}\e^{-nx},
 \qquad
 \sum_{n\ge0}x^{-n}\e^{-\e^x},
 \qquad
 \sum_{n\ge1}x^{1/n},
\]
when the chosen well-based variant supports them. It does not automatically permit one term of height $n$ for every $n$ if those terms violate the support architecture. Enlargements such as omega-series and hyperseries address broader closure goals; we return to them in Chapter~\ref{ch:beyond-LE}.

\section{Exercises}

\begin{exercise}
Assign the smallest plausible exponential height and logarithmic depth to each expression:
\[
 \e^{x+1/x},\quad
 \e^{\e^x+x},\quad
 \log\log x+x^{-1},\quad
 \exp\!\bigl(\sqrt{\log x}+\e^{-\log x}\bigr).
\]
Remember to expand small pieces in exponents when appropriate.
\end{exercise}

\begin{exercise}
Write $\exp(x+2+x^{-1}+\e^{-x})$ as a constant coefficient times a monomial times a small formal series through terms of dominance at least $\e^xx^{-2}$.
\end{exercise}

\begin{exercise}
Why would allowing $\exp(c+L)$ and $\exp(L)$ as unrelated monomials destroy a canonical coefficient-monomial decomposition?
\end{exercise}

\begin{summarybox}
The transseries field is built in stages. Height zero consists of generalized powers. Purely large elements at one stage are exponentiated to create monomials at the next. Finite composition with iterated logarithms adds logarithmic depth. Constant and small exponent pieces are kept out of monomials to preserve canonical form.
\end{summarybox}


\chapter{Arithmetic, inversion, and real powers}\label{ch:arithmetic}

The Hahn-field architecture makes addition and multiplication automatic. Division, roots, and general real powers follow from one canonical factorization.

\section{Canonical multiplicative decomposition}

Let $T\ne0$. Extract its leading term:
\[
 T=a\m+\text{smaller terms},
 \qquad a=\lc(T),\quad \m=\lm(T).
\]
Factor it:
\begin{equation}\label{eq:multiplicative-decomposition}
 T=a\m(1+S),
 \qquad S\precdom1.
\end{equation}
This decomposition is unique. It plays the role of writing a nonzero formal power series as
\[
 a_kt^k(1+\text{higher powers of }t).
\]

\section{Inversion}

Using the formal geometric series,
\begin{equation}\label{eq:inverse-general}
 T^{-1}=a^{-1}\m^{-1}
 \left(1-S+S^2-S^3+\cdots\right).
\end{equation}
Every nonzero transseries therefore has an inverse.

\begin{workedbox}
Invert
\[
 T=\e^x+x^2+1.
\]
The leading monomial is $\e^x$, so
\[
 T=\e^x\left(1+x^2\e^{-x}+\e^{-x}\right).
\]
Set $S=(x^2+1)\e^{-x}$. Then
\begin{align*}
 T^{-1}
 & =\e^{-x}\left(1-S+S^2-S^3+\cdots\right)\\
 & =\e^{-x}-(x^2+1)\e^{-2x}
 +(x^2+1)^2\e^{-3x}
 -(x^2+1)^3\e^{-4x}+\cdots.
\end{align*}
Each correction is exponentially smaller than the preceding sector.
\end{workedbox}

\section{Integer and rational powers}

For $r\in\Z$, $T^r$ is defined by multiplication and inversion. If $T>0$ and $q\ge1$ is an integer, a $q$th root can be constructed from \cref{eq:multiplicative-decomposition}:
\[
 T^{1/q}=a^{1/q}\m^{1/q}(1+S)^{1/q}.
\]
The binomial series is
\begin{equation}\label{eq:binomial-formal}
 (1+S)^\alpha
 =\sum_{n=0}^{\infty}\binom{\alpha}{n}S^n,
 \qquad
 \binom{\alpha}{n}
 =\frac{\alpha(\alpha-1)\cdots(\alpha-n+1)}{n!}.
\end{equation}
It is formally summable whenever $S\precdom1$. For rational $\alpha$, this gives algebraic roots; after logarithms are defined, the same formula agrees with $\exp(\alpha\log(1+S))$ for every real $\alpha$.

\begin{example}
For $x\succdom1$,
\begin{align*}
 \sqrt{x^2+x}
 &=x\sqrt{1+x^{-1}}\\
 &=x\left(1+\frac1{2x}-\frac1{8x^2}
 +\frac1{16x^3}-\frac5{128x^4}+\cdots\right)\\
 &=x+\frac12-\frac1{8x}+\frac1{16x^2}
 -\frac5{128x^3}+\cdots.
\end{align*}
The familiar binomial asymptotic expansion is an identity in the Hahn field.
\end{example}

\section{Order and sign}

Define $T>0$ if $\lc(T)>0$. This order extends the real order and is compatible with field operations. In particular:
\[
 S<T\quad\Longleftrightarrow\quad T-S>0.
\]
An infinitesimal can be positive without being zero:
\[
 0<x^{-1}<1,
 \qquad
 0<\e^{-x}<x^{-N}
 \quad\text{for every fixed }N.
\]
Similarly, $x$ is infinitely large because $x>n$ for every real constant $n$ in the formal order.

The field is therefore \emph{non-Archimedean}: the Archimedean principle that every number is bounded by an integer fails.

\section{Valuation language}

The leading monomial defines a valuation-like measure of size. One may assign to $T$ the value corresponding to $\lm(T)$, with the convention that smaller monomials have larger valuation. Then
\[
 v(ST)=v(S)+v(T),
 \qquad
 v(S+T)\ge\min\{v(S),v(T)\},
\]
with equality unless leading terms cancel.

For beginners, leading-monomial language is usually more concrete. Valuations become valuable when studying differential equations abstractly, because they turn asymptotic size into algebraic inequalities.

\section{Cancellation and the next term}

Suppose
\[
 S=3\e^x+2x+\cdots,
 \qquad
 T=-3\e^x+5x+\cdots.
\]
Then
\[
 S+T=7x+\cdots.
\]
The leading exponential cancels, and the sum drops to a much smaller scale. Algorithms must therefore combine equal monomials before determining the leading term. A sparse dictionary representation naturally enforces this.

\begin{warningbox}
It is unsafe to infer the leading term of a sum from the leading terms of the summands without checking cancellation. Dominance is not a floating-point approximation; coefficients are exact data.
\end{warningbox}

\section{Real closedness}

A major structural theorem says that the real transseries field is real closed. Informally, this means it has the same algebraic completeness properties as the real numbers.

\begin{theorem}[Real closedness, stated without proof]\label{thm:real-closed}
The standard field $\T$ of real logarithmic-exponential transseries is real closed. Consequently:
\begin{enumerate}
\item every positive transseries has a square root;
\item every polynomial of odd degree over $\T$ has a root in $\T$;
\item every polynomial over $\T$ factors into linear and irreducible quadratic factors.
\end{enumerate}
\end{theorem}

At the Hahn-field level, divisibility of the monomial group and real closedness of the coefficient field are the key ingredients. The inductive exponential-logarithmic construction preserves the required structure. Full proofs are found in the generalized-series and transseries literature \cite{vdDMM-1997,vdDMM-2001,vdhoeven-book}.

Real closedness is an existence theorem, not an automatic recipe. Dominant balance and fixed-point iteration provide practical expansions of roots; we develop them in Part~IV.

\section{Newton's method as formal algebra}

Suppose $P(Y)\in\T[Y]$ and an approximation $Y_0$ makes $P(Y_0)$ small relative to $P'(Y_0)$. The Newton correction is
\[
 Y_1=Y_0-\frac{P(Y_0)}{P'(Y_0)}.
\]
If the correction gains dominance order, iteration often stabilizes coefficient by coefficient. This is not numerical convergence in a metric; it is convergence in the valuation topology.

For instance, to solve $Y^2=x^2+x$, start with $Y_0=x$. Then
\[
 Y_1=x-\frac{x^2-(x^2+x)}{2x}
 =x+\frac12.
\]
A second step yields
\[
 Y_2=x+\frac12-\frac{(x+1/2)^2-(x^2+x)}{2x+1}
 =x+\frac12-\frac{1}{8x}+\bigO(x^{-2}),
\]
recovering the binomial expansion progressively.

\section{Exercises}

\begin{exercise}
Invert $x^2-3x+1$ through order $x^{-6}$.
\end{exercise}

\begin{exercise}
Expand $(1+2x^{-1}+\e^{-x})^{1/2}$ through all monomials at least as large as $x^{-2}\e^{-x}$.
\end{exercise}

\begin{exercise}
Let $S,T>0$ with $S\asym T$. Show formally that $S^\alpha\asym T^\alpha$ for any fixed real $\alpha$, assuming the logarithm and exponential constructions of the next chapter.
\end{exercise}

\begin{summarybox}
Every nonzero transseries factors uniquely as coefficient times leading monomial times $1+$ an infinitesimal. Geometric and binomial series then provide inverses and roots. The resulting ordered field is non-Archimedean and, in its standard form, real closed.
\end{summarybox}


\chapter{Exponentials and logarithms without numerical convergence}\label{ch:explog}

The most important formal operations are defined by separating the part that must become a new monomial from the part that can be expanded as an ordinary power series.

\section{The exponential of a transseries}

Write the canonical additive decomposition
\[
 T=L+c+S,
\]
where $L$ is purely large, $c\in\R$, and $S\precdom1$. Define
\begin{equation}\label{eq:formal-exp-definition}
 \exp(T)=\e^c\,\e^L
 \sum_{n=0}^{\infty}\frac{S^n}{n!}.
\end{equation}
Each factor has a distinct role:

\begin{itemize}
\item $\e^L$ is a transmonomial introduced by the height construction;
\item $\e^c$ is an ordinary real coefficient;
\item $\sum S^n/n!$ is a well-based formal series because $S$ is infinitesimal.
\end{itemize}

No claim is made that substituting a finite numerical $x$ into the series converges. Formal summability follows from support moving to smaller monomials as powers of $S$ increase.

\begin{workedbox}
Expand
\[
 \exp\!\left(x+2+x^{-1}+3\e^{-x}\right).
\]
Here
\[
 L=x,
 \qquad c=2,
 \qquad S=x^{-1}+3\e^{-x}.
\]
Therefore
\[
 \exp(T)=\e^2\e^x\left(1+S+\frac{S^2}{2}+\frac{S^3}{6}+\cdots\right).
\]
The beginning, in dominance order, is
\begin{align*}
 \e^2\e^x\biggl(&1+x^{-1}+\frac12x^{-2}+\frac16x^{-3}+\cdots\\
 &+3\e^{-x}+3x^{-1}\e^{-x}
 +\frac32x^{-2}\e^{-x}+\cdots
 +\frac92\e^{-2x}+\cdots\biggr).
\end{align*}
After multiplying by $\e^x$, the sector $3\e^{-x}$ inside the parentheses contributes the constant term $3\e^2$. Formal dominance must be assessed after all monomial factors are combined.
\end{workedbox}

\section{Why $\exp$ respects addition}

Given $A$ and $B$, one expects
\[
 \exp(A+B)=\exp(A)\exp(B).
\]
The proof reduces to the ordinary formal identity for small parts, together with multiplication of monomials for large parts and real exponentials for constants. The support lemmas justify rearranging the finite contribution to each monomial. Thus the exponential is a homomorphism from the additive group of $\T$ to its positive multiplicative group.

\section{The logarithm of a positive transseries}

Let $T>0$. Its multiplicative decomposition has positive leading coefficient:
\[
 T=a\m(1+S),
 \qquad a>0,\quad S\precdom1.
\]
Define
\begin{equation}\label{eq:formal-log-definition}
 \log T
 =\log a+\log\m
 +\sum_{n=1}^{\infty}\frac{(-1)^{n+1}}{n}S^n.
\end{equation}
The logarithm of a monomial is already encoded in its construction. If
\[
 \m=x^b\e^L,
\]
then
\begin{equation}\label{eq:log-monomial}
 \log\m=b\log x+L.
\end{equation}
The small factor uses the familiar series
\[
 \log(1+S)=S-\frac{S^2}{2}+\frac{S^3}{3}-\cdots.
\]

\begin{workedbox}
Find $\log(\e^x+x^2)$. Factor the leading monomial:
\[
 \e^x+x^2=\e^x(1+x^2\e^{-x}).
\]
Therefore
\begin{align*}
 \log(\e^x+x^2)
 &=x+\log(1+x^2\e^{-x})\\
 &=x+x^2\e^{-x}-\frac12x^4\e^{-2x}
 +\frac13x^6\e^{-3x}-\cdots.
\end{align*}
The correction is exponentially small even though $x^2$ by itself is large.
\end{workedbox}

\section{Formal inverse identities}

The definitions are designed so that
\[
 \log(\exp T)=T
\]
for all real transseries $T$, and
\[
 \exp(\log T)=T
\]
for all positive transseries $T$. On the small part, these are the standard identities of formal power series. The canonical decompositions ensure that large and constant pieces are recovered exactly.

\begin{proposition}
The exponential is an order-preserving group isomorphism
\[
 (\T,+)\longrightarrow(\T_{>0},\times),
\]
with inverse logarithm.
\end{proposition}

This is one reason transseries are more than a collection of asymptotic notations. They form an ordered exponential field.

\section{General real powers}

For $T>0$ and $\alpha\in\R$, define
\begin{equation}\label{eq:real-power}
 T^\alpha=\exp(\alpha\log T).
\end{equation}
Using $T=a\m(1+S)$,
\[
 T^\alpha=a^\alpha\m^\alpha(1+S)^\alpha,
\]
and the last factor is the binomial series \cref{eq:binomial-formal}.

\begin{example}
Let
\[
 T=x^3\e^{\sqrt{x}}(1+2x^{-1}+\e^{-x}).
\]
Then
\[
 T^{\sqrt2}
 =x^{3\sqrt2}\e^{\sqrt{2x}}
 \left(1+2x^{-1}+\e^{-x}\right)^{\sqrt2}.
\]
The small factor begins
\[
 1+\sqrt2(2x^{-1}+\e^{-x})
 +\frac{\sqrt2(\sqrt2-1)}2(2x^{-1}+\e^{-x})^2+\cdots.
\]
\end{example}

\section{Nested examples}

The recursive construction handles expressions that would be awkward in a single flat formula.

\subsection*{Example 1: a small perturbation of a huge exponential}

Let
\[
 T=\exp\!\left(\e^x+x^{-1}\right).
\]
Then
\[
 T=\e^{\e^x}\exp(x^{-1})
 =\e^{\e^x}\left(1+x^{-1}+\frac1{2x^2}+\cdots\right).
\]
The monomial has height $2$, while the power corrections remain in the same sector.

\subsection*{Example 2: a logarithm peels off the top layer}

Let
\[
 T=3x^7\exp\!\left(\e^x-2\sqrt{x}\right)
 \left(1+x^{-1}+\e^{-x}\right).
\]
Then
\begin{align*}
 \log T={}&\log3+7\log x+\e^x-2\sqrt{x}\\
 &+\left(x^{-1}+\e^{-x}\right)
 -\frac12\left(x^{-1}+\e^{-x}\right)^2+\cdots.
\end{align*}
The logarithm converts multiplicative layers into additive layers and lowers the outer exponential height.

\section{Complex logarithms}

The real ordered field has a logarithm only for positive elements. One can adjoin $\ii$ and work with complex coefficients, but then there is no compatible total order and logarithms become branch dependent:
\[
 \Log T=\log|T|+\ii\arg T.
\]
Resurgent analysis necessarily uses complex directions and branches. The formal ordered-field chapters avoid this complication until Part~V.

\section{Exercises}

\begin{exercise}
Expand $\exp(x+\log x+x^{-1})$ through order $\e^x x^{-3}$.
\end{exercise}

\begin{exercise}
Compute the first four nonzero correction terms in
\[
 \log(x^5+\e^{-x}).
\]
\end{exercise}

\begin{exercise}
Show formally that
\[
 \log(AB)=\log A+\log B
\]
for positive $A,B$ by reducing to their leading coefficients, monomials, and small factors.
\end{exercise}

\begin{summarybox}
To exponentiate, split a transseries into purely large, constant, and small parts. The large part creates a monomial, the constant becomes a coefficient, and the small part is expanded. To take a logarithm, factor a positive transseries into coefficient, monomial, and $1+$ infinitesimal. These definitions make $\exp$ and $\log$ exact formal inverses.
\end{summarybox}


\chapter{What formal equality does and does not mean}\label{ch:formal-meaning}

Before differentiating or solving equations, it is worth separating several notions that are often compressed into the same equals sign.

\section{Four levels of statement}

\subsection{Identity of formal transseries}

Two transseries are equal if every monomial has the same coefficient. For example,
\[
 (1+x^{-1})^{-1}=1-x^{-1}+x^{-2}-x^{-3}+\cdots
\]
is an exact identity in the formal field. It does not depend on choosing a numerical value of $x$.

\subsection{Asymptotic expansion of a function}

A function $f$ may satisfy
\[
 f(x)\asym T(x)
\]
in the sense that every truncation of $T$ approximates $f$ to the order of the next omitted monomial. This is an analytic theorem and may hold only in a sector of the complex plane.

\subsection{Equality of germs}

Two functions define the same germ at $+\infty$ if they agree for all sufficiently large $x$. A Hardy field is a field of such germs. A transseries may map to a germ under an asymptotic expansion map, but constructing an injective map for broad classes is a substantial theorem, not notation.

\subsection{Equality after summation}

A divergent formal transseries may be assigned an analytic function by a summation method. Directional Borel sums can differ across a Stokes ray. Thus one should write
\[
 \lateralsum{\theta}\trans F
\]
rather than silently replacing $\trans F$ by a unique function.

\section{The danger of plugging in numbers}

Take
\[
 T(x)=\sum_{n\ge0}\frac{n!}{x^{n+1}}.
\]
As a formal transseries it is perfectly valid. At every fixed finite $x$, its numerical terms eventually grow, so the numerical sum diverges. Writing $T(10)$ without specifying truncation or summation is meaningless.

By contrast,
\[
 G(x)=\sum_{n\ge0}x^{-n}
\]
converges numerically for $x>1$ and equals $x/(x-1)$. It is simultaneously a formal identity and a convergent analytic representation. These are two independent virtues.

\section{Truncation is the native observation operation}

The natural way to inspect a transseries is not to evaluate it but to truncate it by dominance. Given a threshold $\m$, retain all terms larger than $\m$:
\[
 T_{\succdom\m}=\sum_{\n\succdom\m}a_{\n}\n.
\]
An algorithm is correct to order $\m$ if its output has the same truncation as the exact formal answer above that threshold.

This viewpoint resembles exact arithmetic with lazily generated coefficients. One asks for as much of the ordered support as a later calculation needs.

\section{Closure properties of the standard field}

The standard logarithmic-exponential transseries field is closed under
\[
 +,\ - ,\ \times,\ /,\ \exp,\ \log\text{ on positives},
\]
and, as later chapters explain, under differentiation, antiderivatives, and suitable composition. It is real closed and supports a rich differential algebra \cite{vdhoeven-book,adh-book}.

Yet it is not literally ``all asymptotic expressions.'' Important phenomena outside the basic real ordered field include:

\begin{itemize}
\item oscillatory monomials such as $\e^{\ii x}$, which do not possess an eventual real order;
\item arbitrary infinitely nested exponentials and logarithms;
\item transmonomials requiring superexponentials or hyperserial structure;
\item functions with wild zero sets or oscillations incompatible with Hardy-field behavior;
\item nonperturbative analytic data, such as Stokes constants, not determined by formal field operations alone.
\end{itemize}

Enlargements and complex variants handle portions of this list, but no single framework erases all distinctions.

\section{A useful discipline for calculations}

At each step, ask four questions:

\begin{enumerate}
\item What is the monomial group and dominance order?
\item Is the support well based or grid based so that the operation is formally defined?
\item Is the statement only formal, or is an actual asymptotic function involved?
\item If a divergent series is being assigned a value, what summation direction or prescription is used?
\end{enumerate}

These questions prevent most conceptual mistakes in introductory transseries work.

\begin{checkpointbox}
Formal algebra answers ``which coefficients must occur?'' Asymptotic analysis answers ``how well do truncations approximate a function?'' Summability answers ``which function, if any, is reconstructed from a divergent object?'' Resurgence answers ``how are the sectors and their discontinuities related?''
\end{checkpointbox}

\section{Exercises}

\begin{exercise}
For each statement below, classify it as a formal identity, an asymptotic assertion, an equality of functions, or a summation assertion:
\begin{enumerate}
\item $(1-S)^{-1}=\sum_{n\ge0}S^n$ for $S\precdom1$;
\item $\Ei(x)\e^{-x}\asym\sum_{n\ge0}n!x^{-n-1}$;
\item $\sum_{n\ge0}x^{-n}=x/(x-1)$ for $x>1$;
\item $\lateralsum{0^+}\trans F-\lateralsum{0^-}\trans F=2\pi\ii\e^{-x}$.
\end{enumerate}
\end{exercise}

\begin{exercise}
Give two distinct functions with the same asymptotic expansion in powers of $x^{-1}$, and then give a transseries scale that distinguishes them.
\end{exercise}

\begin{summarybox}
Formal identity, asymptotic expansion, equality of germs, and equality after summation are different claims. A transseries is natively inspected by dominance truncation. Numerical substitution is justified only when convergence or an explicit summation rule has been established.
\end{summarybox}

\part{Calculus, composition, and recursive construction}

\chapter{Differentiation}\label{ch:differentiation}

Formal differentiation is one of the strongest reasons to impose disciplined support conditions. We want to differentiate term by term, preserve the transseries field, and retain the product, quotient, and chain rules.

\section{Differentiate monomials first}

A general log-free transmonomial can be written
\[
 \m=x^a\e^L,
\]
where $L$ is purely large at a lower height. Its derivative is
\begin{equation}\label{eq:monomial-derivative}
 \m'=\left(\frac{a}{x}+L'\right)\m.
\end{equation}
This recursively reduces differentiation of a monomial to differentiation of its exponent.

At logarithmic depth, the usual chain rule applies. For example,
\begin{align*}
 \frac{\dd}{\dd x}(\log x)^a
 &=\frac{a}{x\log x}(\log x)^a,\\
 \frac{\dd}{\dd x}\e^{\sqrt{\log x}}
 &=\frac{1}{2x\sqrt{\log x}}\e^{\sqrt{\log x}},\\
 \frac{\dd}{\dd x}\e^{-\e^x}
 &=-\e^x\e^{-\e^x}.
\end{align*}

The multiplier $\m'/\m$ is the logarithmic derivative. It often changes the power or logarithmic prefactor but preserves the essential exponential scale.

\section{Termwise differentiation}

For
\[
 T=\sum_{\m}a_{\m}\m,
\]
define
\begin{equation}\label{eq:termwise-derivative}
 T'=\sum_{\m}a_{\m}\m'.
\end{equation}
The right side must be a formally summable family. This is not automatic for arbitrary generalized series, but the transseries monomial hierarchy is constructed so that differentiation preserves well-based support.

\begin{theorem}[Closure under differentiation]\label{thm:diff-closure}
The standard logarithmic-exponential transseries field $\T$ carries a derivation $T\mapsto T'$ satisfying:
\begin{align*}
 (S+T)'&=S'+T',\\
 (ST)'&=S'T+ST',\\
 (T^{-1})'&=-T^{-2}T',\\
 (\exp T)'&=T'\exp T,\\
 (\log T)'&=T'/T\qquad(T>0).
\end{align*}
Its constant field is $\R$: if $T'=0$, then $T\in\R$.
\end{theorem}

The support proof is technical because differentiation may replace one monomial by a whole transseries times that monomial. At a conceptual level, the exponent and logarithmic-depth recursion ensures that derivatives move within a controlled collection of lower-height factors. See \cite{vdhoeven-book,edgar-beginners,adh-book} for rigorous treatments.

\section{Power and logarithmic series}

At height zero, differentiation is familiar:
\[
 \left(\sum_ba_bx^b\right)'
 =\sum_bba_bx^{b-1}.
\]
Examples:
\begin{align*}
 \left(x^{\sqrt2}+3+x^{-\pi}\right)'
 &=\sqrt2x^{\sqrt2-1}-\pi x^{-\pi-1},\\
 \left(\log x+\frac1{\log x}
 +\frac{2}{(\log x)^2}+\cdots\right)'
 &=\frac1x-\frac1{x(\log x)^2}
 -\frac4{x(\log x)^3}+\cdots.
\end{align*}

A constant term disappears, which can reveal a much smaller leading derivative. This is ordinary cancellation, but across transseries scales the drop can be dramatic.

\begin{example}
Let
\[
 T=x+100+\e^{-x^2}.
\]
Then $T\asym x$, while
\[
 T'=1-2x\e^{-x^2}\asym1.
\]
The leading scale drops from $x$ to $1$.
\end{example}

\begin{example}
Let
\[
 T=1+\e^{-\e^x}.
\]
The nonconstant part is extraordinarily small, but
\[
 T'=-\e^x\e^{-\e^x}.
\]
The derivative is larger than $\e^{-\e^x}$ by a factor $\e^x$, yet it remains smaller than every $\e^{-Nx}$.
\end{example}

\section{Differentiating an exponential sector expansion}

Suppose
\[
 T=\sum_{k\ge0}C^k\e^{-kAx}T_k(x),
\]
where $A>0$ and each $T_k$ is an inverse-power series. Then
\begin{equation}\label{eq:sector-derivative}
 T'=\sum_{k\ge0}C^k\e^{-kAx}
 \left(T_k'-kA T_k\right).
\end{equation}
Differentiation preserves the sector index $k$. Nonlinear multiplication, by contrast, convolves sector indices. This separation is what makes coefficient equations for nonlinear ODEs manageable.

\section{The chain rule}

When composition is defined,
\begin{equation}\label{eq:formal-chain-rule}
 (T\circ S)'=(T'\circ S)S'.
\end{equation}
For example,
\[
 T(x)=\e^{x^2+x^{-1}},
 \qquad
 S(x)=\log x.
\]
Then
\[
 (T\circ S)(x)=
 \exp\!\left((\log x)^2+(\log x)^{-1}\right),
\]
and
\[
 (T\circ S)'=
 \left(2\log x-(\log x)^{-2}\right)\frac1x
 \exp\!\left((\log x)^2+(\log x)^{-1}\right).
\]

\section{Logarithmic differentiation as a simplifier}

For a product of complicated factors, compute $T'/T$:
\[
 T=x^a(\log x)^b\exp(L).
\]
Then
\begin{equation}\label{eq:general-log-derivative}
 \frac{T'}{T}=\frac{a}{x}+\frac{b}{x\log x}+L'.
\end{equation}
This avoids repeated product rules and immediately reveals the derivative's scale.

\begin{workedbox}
Let
\[
 T=x^{-3}(\log x)^7\exp\!\left(-x^2+\e^{\sqrt{x}}\right).
\]
Then
\[
 \frac{T'}T=-\frac3x+\frac7{x\log x}-2x
 +\frac{\e^{\sqrt{x}}}{2\sqrt{x}}.
\]
The last term dominates $2x$, so
\[
 T'\asym
 \frac{\e^{\sqrt{x}}}{2\sqrt{x}}T.
\]
Even though $T$ contains the decaying factor $\e^{-x^2}$, the nested exponential in its exponent controls the logarithmic derivative.
\end{workedbox}

\section{Differential polynomials}

A differential polynomial in an unknown $Y$ is a polynomial expression in
\[
 Y,Y',Y'',\ldots,Y^{(r)}
\]
with transseries coefficients. For example,
\[
 P(Y)=Y''+xY'-\e^xY^2+x^{-1}.
\]
Substituting a transseries $Y$ and collecting monomials gives another transseries. Solving $P(Y)=0$ means making every monomial coefficient vanish. Dominant balance identifies the first term; recursive cancellation determines the rest.

The differential-algebraic strength of $\T$ is much deeper than closure under differentiation. The field satisfies powerful existence and model-theoretic properties discussed in Chapter~\ref{ch:hfields}.

\section{Exercises}

\begin{exercise}
Differentiate
\[
 x^4(\log x)^{-2}\exp\!\left(x^{1/3}-\e^{-x}\right)
\]
using a logarithmic derivative.
\end{exercise}

\begin{exercise}
Let
\[
 T=\sum_{n\ge0}n!x^{-n-1}+C\e^{-x}.
\]
Verify directly that $T'+T=x^{-1}$ as a formal identity.
\end{exercise}

\begin{exercise}
Suppose $S\precdom1$. Differentiate the formal identity
\[
 \log(1+S)=S-\frac{S^2}{2}+\frac{S^3}{3}-\cdots
\]
and recover $(\log(1+S))'=S'/(1+S)$.
\end{exercise}

\begin{summarybox}
Differentiate monomials through logarithmic derivatives, then differentiate well-based sums term by term. The transseries support architecture guarantees closure and the usual calculus rules. Sector differentiation preserves exponential sector labels, while nonlinear products later mix them.
\end{summarybox}


\chapter{Formal integration and asymptotic antiderivatives}\label{ch:integration}

Integration is subtler than differentiation. The derivative of a single monomial is usually a monomial times a nonconstant transseries, so its antiderivative is often an infinite expansion rather than one term.

\section{Elementary power cases}

For $a\ne-1$,
\[
 \int x^a\dd x=\frac{x^{a+1}}{a+1}.
\]
At the resonant exponent $a=-1$,
\[
 \int x^{-1}\dd x=\log x.
\]
Then another logarithmic hierarchy appears:
\begin{align*}
 \int\frac{1}{x\log x}\dd x&=\log\log x,\\
 \int\frac{1}{x\log x\log\log x}\dd x&=\log_3x.
\end{align*}
These are reminders that logarithmic depth is forced by integration, not added merely for decoration.

\section{Antidifferentiating a fast monomial}

Let $\m$ be a positive monomial and set
\[
 A=\frac{\m'}{\m}=(\log\m)'.
\]
We seek an antiderivative in the form
\begin{equation}\label{eq:integral-ansatz}
 F=\frac{\m}{A}(1+U),
 \qquad U\precdom1.
\end{equation}
Define
\[
 \delta=\frac{A'}{A^2}.
\]
Differentiate:
\begin{align*}
 F'
 &=\m(1+U)+\frac{\m}{A}U'
   -\frac{\m A'}{A^2}(1+U)\\
 &=\m\left[1+U+\frac{U'}A-\delta(1+U)\right].
\end{align*}
The condition $F'=\m$ becomes
\begin{equation}\label{eq:integration-fixed-point}
 U=\delta+\delta U-\frac{U'}A.
\end{equation}
When $\delta\precdom1$ and differentiation divided by $A$ makes terms smaller, this is a contractive formal fixed-point equation. Starting with $U_0=0$ and iterating the right side produces successively smaller corrections.

\begin{intuitionbox}
The first guess $\m/A$ comes from treating $A=(\log\m)'$ as locally constant: since $\m'=A\m$, division by $A$ almost reverses differentiation. The infinite correction compensates for the fact that $A$ itself varies.
\end{intuitionbox}

\section{Example: $\int\e^{x^2}\dd x$}

Here
\[
 \m=\e^{x^2},
 \qquad A=2x,
 \qquad \delta=\frac{2}{4x^2}=\frac1{2x^2}.
\]
Seek
\[
 F=\frac{\e^{x^2}}{2x}
 \left(1+c_1x^{-2}+c_2x^{-4}+\cdots\right).
\]
Substitution or repeated use of \cref{eq:integration-fixed-point} gives
\[
 c_1=\frac12,
 \quad c_2=\frac34,
 \quad c_3=\frac{15}{8},
 \quad c_4=\frac{105}{16},
\]
and in general
\[
 c_n=\frac{(2n-1)!!}{2^n}.
\]
Therefore
\begin{equation}\label{eq:int-exp-x2}
 \int\e^{x^2}\dd x
 \asym\frac{\e^{x^2}}{2x}
 \left(1+\frac1{2x^2}+\frac3{4x^4}
 +\frac{15}{8x^6}+\frac{105}{16x^8}+\cdots\right).
\end{equation}
The series diverges factorially, but it is a valid formal antiderivative.

\begin{workedbox}
Differentiate the first three corrections:
\[
 F_2=\frac{\e^{x^2}}{2x}
 \left(1+\frac1{2x^2}+\frac3{4x^4}\right).
\]
A direct calculation gives
\[
 F_2'=\e^{x^2}\left(1-\frac{15}{8x^6}\right).
\]
The first omitted coefficient is exactly the one needed to cancel the residual term.
\end{workedbox}

\section{Example: $\int\e^{\e^x}\dd x$}

Let $\m=\e^{\e^x}$. Then
\[
 A=\e^x,
 \qquad \delta=\e^{-x}.
\]
The fixed-point expansion produces factorial coefficients:
\begin{equation}\label{eq:int-double-exp}
 \int\e^{\e^x}\dd x
 \asym \e^{\e^x-x}
 \left(1+\e^{-x}+2\e^{-2x}+6\e^{-3x}
 +24\e^{-4x}+\cdots\right).
\end{equation}
Indeed, substituting $t=\e^x$ gives
\[
 \int\e^{\e^x}\dd x=\int\frac{\e^t}{t}\dd t=\Ei(t),
\]
and the familiar factorial expansion of $\Ei(t)$ reappears with $t=\e^x$.

This example is a useful scale elevator: the same divergent pattern occurs one exponential level higher.

\section{Slow monomials and integration by parts}

If $\m$ varies slowly, the approximation $\m/A$ may not be the right starting point. For example,
\[
 \int\frac{\log x}{x^2}\dd x
 =-\frac{\log x+1}{x}.
\]
One may integrate powers first and treat logarithms by ordinary integration by parts.

A practical integration algorithm classifies the leading monomial according to the size of its logarithmic derivative, chooses a leading antiderivative, subtracts its derivative, and repeats on the smaller residual. Van der Hoeven's algorithmic treatment makes this systematic \cite{vdhoeven-book}.

\section{Termwise integration and constants}

It is tempting to write
\[
 \int\sum_{\m}a_{\m}\m\dd x
 =\sum_{\m}a_{\m}\int\m\dd x.
\]
This is often correct, but formal summability of the family of antiderivatives must be checked. The general closure theorem is stronger than a naive termwise rule.

\begin{theorem}[Antiderivative closure, stated without proof]\label{thm:antiderivative}
For every $T\in\T$ there exists $F\in\T$ such that $F'=T$. Any two antiderivatives differ by a real constant.
\end{theorem}

Together with real closedness and closure under exponential integration, this is part of the statement that $\T$ is \emph{Liouville closed}. The differential-algebraic meaning is discussed in Chapter~\ref{ch:hfields}.

\section{An inverse-power recurrence}

Suppose we want
\[
 \int x^\alpha\e^{\lambda x}\dd x,
 \qquad \lambda\ne0.
\]
Set
\[
 F=\frac{x^\alpha\e^{\lambda x}}{\lambda}
 \sum_{n\ge0}c_nx^{-n},
 \qquad c_0=1.
\]
Differentiation gives the recurrence
\[
 c_{n+1}=-\frac{\alpha-n}{\lambda}c_n,
\]
so
\begin{equation}\label{eq:exp-power-integral}
 \int x^\alpha\e^{\lambda x}\dd x
 \asym\frac{x^\alpha\e^{\lambda x}}{\lambda}
 \left(1-\frac{\alpha}{\lambda x}
 +\frac{\alpha(\alpha-1)}{\lambda^2x^2}-\cdots\right).
\end{equation}
For nonnegative integer $\alpha$, the series terminates and gives the elementary exact antiderivative. For general $\alpha$, it is infinite and typically divergent.

\section{Exercises}

\begin{exercise}
Derive the recurrence for \cref{eq:int-exp-x2} directly by substituting
\[
 F=\e^{x^2}\sum_{n\ge0}a_nx^{-2n-1}.
\]
\end{exercise}

\begin{exercise}
Find the first four terms of an asymptotic antiderivative of
\[
 x^3\e^{-x^2}.
\]
In this case the exact integral is elementary; compare the results.
\end{exercise}

\begin{exercise}
Apply \cref{eq:exp-power-integral} to $\alpha=-1$ and $\lambda=1$. Explain the relationship with the Euler factorial series.
\end{exercise}

\begin{summarybox}
A monomial's logarithmic derivative suggests its leading antiderivative. Correcting $\m/A$ leads to a contractive fixed-point equation and often a divergent but useful transseries. Logarithmic depth appears at resonant power exponents, and the full transseries field is closed under antiderivatives.
\end{summarybox}


\chapter{Composition and formal Taylor expansion}\label{ch:composition}

Composition is where transseries become a language of functions rather than merely an ordered field. It is also one of the technically delicate operations, because substituting an infinite series into another infinite series can destroy support control if done indiscriminately.

\section{Which inner arguments are allowed?}

For the standard real transseries composition, the inner argument $S$ is usually required to be positive and infinitely large:
\[
 S>0,
 \qquad S\succdom1.
\]
This matches the interpretation that $T(x)$ describes behavior at $x=+\infty$ and we are replacing $x$ by another quantity tending to $+\infty$.

At the elementary level, define composition recursively:
\begin{align*}
 c\circ S&=c,\\
 x\circ S&=S,\\
 (A+B)\circ S&=(A\circ S)+(B\circ S),\\
 (AB)\circ S&=(A\circ S)(B\circ S),\\
 (\exp A)\circ S&=\exp(A\circ S),\\
 (\log A)\circ S&=\log(A\circ S),
\end{align*}
followed by extension to well-based sums. The last extension is the hard part; Edgar proves appropriate support and convergence properties in \cite{edgar-composition}.

\section{First examples}

\begin{example}
If
\[
 T(x)=x^2+\e^{-x},
 \qquad
 S(x)=\log x+x^{-1},
\]
then
\[
 (T\circ S)(x)=(\log x+x^{-1})^2
 +\exp(-\log x-x^{-1}).
\]
The exponential term becomes
\[
 x^{-1}\exp(-x^{-1})
 =x^{-1}-x^{-2}+\frac12x^{-3}-\cdots.
\]
\end{example}

\begin{example}
For $T(x)=\log x$ and $S=x+x^{-1}$,
\begin{align*}
 T\circ S
 &=\log(x+x^{-1})\\
 &=\log x+\log(1+x^{-2})\\
 &=\log x+x^{-2}-\frac12x^{-4}
 +\frac13x^{-6}-\cdots.
\end{align*}
\end{example}

\section{Composition with a small perturbation}

Suppose
\[
 S=S_0+U,
 \qquad U\precdom S_0,
\]
and the derivatives of $T$ are well behaved relative to $U$. One expects a formal Taylor expansion
\begin{equation}\label{eq:formal-taylor}
 T(S_0+U)
 =\sum_{n=0}^{\infty}\frac{T^{(n)}(S_0)}{n!}U^n.
\end{equation}
For ordinary analytic functions this follows from convergence near a point. In transseries it is a theorem about support and formal summability. Its hypotheses depend on the relative depth of $T$, $S_0$, and $U$; careless use in ``deep'' cases can fail.

For the common case $U\precdom1$ and $S_0=x$, the formula behaves exactly as intuition suggests.

\begin{workedbox}
Expand $\e^{x+x^{-1}}$ around $x$:
\[
 T(x)=\e^x,
 \qquad U=x^{-1}.
\]
Since $T^{(n)}(x)=\e^x$,
\[
 T(x+U)=\e^x\sum_{n\ge0}\frac{x^{-n}}{n!}
 =\e^x\left(1+x^{-1}+\frac1{2x^2}+\cdots\right).
\]
This agrees with the exponential decomposition of Chapter~\ref{ch:explog}.
\end{workedbox}

\section{A less trivial Taylor calculation}

Let $T(x)=\log x$ and $U=\e^{-x}$. Since
\[
 T^{(n)}(x)=(-1)^{n-1}(n-1)!x^{-n}
 \qquad(n\ge1),
\]
we obtain
\begin{align*}
 \log(x+\e^{-x})
 &=\log x+
 \sum_{n\ge1}\frac{(-1)^{n-1}}n
 x^{-n}\e^{-nx}\\
 &=\log x+x^{-1}\e^{-x}
 -\frac12x^{-2}\e^{-2x}
 +\frac13x^{-3}\e^{-3x}-\cdots.
\end{align*}
Again this is the same result obtained by factoring $x$ and expanding $\log(1+x^{-1}\e^{-x})$.

\section{Associativity and the chain rule}

For admissible positive infinitely large arguments,
\[
 (T\circ S)\circ U=T\circ(S\circ U).
\]
Composition has identity $x$:
\[
 T\circ x=x\circ T=T.
\]
Together with \cref{eq:formal-chain-rule}, these properties make suitable transseries into a group under composition whenever compositional inverses exist.

\section{Compositional inverses}

Suppose $T$ is positive, infinitely large, and strictly increasing in the formal sense $T'>0$. Under broad conditions there is a transseries $S$ such that
\[
 T\circ S=S\circ T=x.
\]
The leading monomial of $S$ is found by reversing the dominant behavior of $T$. Corrections can then be obtained by fixed-point or Newton iteration.

Examples:
\begin{align*}
 T=x^a\quad(a>0)
 &\Longrightarrow\quad T^{-1}=x^{1/a},\\
 T=\e^x
 &\Longrightarrow\quad T^{-1}=\log x,\\
 T=x+\log x
 &\Longrightarrow\quad T^{-1}=x-\log x+\frac{\log x}{x}+\cdots.
\end{align*}
The inverse of $x\e^x$ is Lambert's $W$ function and provides our main worked example in Chapter~\ref{ch:lambertW}.

\section{Taylor's theorem in the transseries setting}

A modern formal Taylor theorem can be summarized as follows.

\begin{theorem}[Formal Taylor principle, schematic version]\label{thm:taylor-schematic}
Let $T,S,U\in\T$, with $S$ positive and infinitely large and $U$ sufficiently small relative to the variation scale of $T$ at $S$. Then the family
\[
 \left(\frac{T^{(n)}\circ S}{n!}U^n\right)_{n\ge0}
\]
is formally summable and its sum equals $T\circ(S+U)$.
\end{theorem}

The phrase ``sufficiently small'' cannot be deleted. Edgar distinguishes shallow, moderate, and deep substitutions according to comparisons between $U$ and logarithmic derivatives of monomials \cite{edgar-composition,edgar-fractional}. Recent work continues to refine Taylor principles for generalized series and broader transserial settings \cite{bagayoko-mantova-taylor,mantova-monotonicity}.

\section{Why composition can fail naively}

Consider a formal sum $T=\sum_n\m_n$. After substitution, it could happen that infinitely many distinct $\m_n\circ S$ collapse onto the same monomial, creating an undefined infinite coefficient sum. Or the substituted support might acquire an infinite increasing chain. The composition theorems rule out such behavior under their hypotheses.

A second danger is using a Taylor expansion when $U$ is not small on the derivative scale. For a very deep monomial $\m$, the quantity
\[
 \frac{\m'}{\m}U
\]
may be large even when $U\precdom x$. Then successive Taylor terms need not decrease.

\begin{warningbox}
The condition $U\precdom S$ is not by itself enough for every Taylor expansion. Check the logarithmic derivative: a useful first test is whether $U(T'/T)\circ S\precdom1$ for the monomials that dominate $T$.
\end{warningbox}

\section{Exercises}

\begin{exercise}
Compute $\log(x+\log x)$ through terms of order $(\log x)^3/x^3$.
\end{exercise}

\begin{exercise}
Let $T(x)=x^3$ and $S=x+\e^{-x}$. Verify the formal Taylor expansion directly and observe that it terminates.
\end{exercise}

\begin{exercise}
Find the first three corrections to the compositional inverse of $T=x+\log x$.
\end{exercise}

\begin{summarybox}
Composition substitutes one positive infinitely large transseries into another, with support conditions doing the work that analytic convergence would do for ordinary functions. Formal Taylor expansion is valid when the perturbation is small on the relevant derivative scale. Composition supports chain rules, inverses, and iteration.
\end{summarybox}

\chapter{Lambert's \texorpdfstring{$W$}{W} function as a transseries inverse}\label{ch:lambertW}

Lambert's function $W$ is defined by
\begin{equation}\label{eq:W-definition}
 W(x)\e^{W(x)}=x.
\end{equation}
It is the compositional inverse of $T\mapsto T\e^T$. Its expansion at $+\infty$ is a model calculation: dominant balance finds the logarithmic scales, and a contractive equation generates all coefficients.

\section{Find the dominant term}

Taking logarithms of \cref{eq:W-definition} gives
\begin{equation}\label{eq:W-log-equation}
 W+\log W=\log x.
\end{equation}
Since $\log W\precdom W$ for a large $W$, the first approximation is
\[
 W\asym\log x.
\]
Substituting this back into the smaller term suggests
\[
 W\approx\log x-\log\log x.
\]
Introduce the standard abbreviations
\begin{equation}\label{eq:L1L2}
 L_1=\log x,
 \qquad
 L_2=\log\log x.
\end{equation}
We therefore write
\begin{equation}\label{eq:W-u-ansatz}
 W=L_1-L_2+u,
 \qquad u\precdom L_2.
\end{equation}

\section{Derive a fixed-point equation}

Because
\[
 W=L_1\left(1-\frac{L_2-u}{L_1}\right),
\]
we have
\[
 \log W=L_2+
 \log\left(1-\frac{L_2-u}{L_1}\right).
\]
Substitute into \cref{eq:W-log-equation}. The $L_1$ and $L_2$ terms cancel, leaving
\begin{equation}\label{eq:W-fixed-point}
 u=-\log\left(1-\frac{L_2-u}{L_1}\right).
\end{equation}
The ratio $L_2/L_1$ is small. The right side can therefore be expanded and iterated.

A tempting alternative is to set $W=L_1+Q$ and use
\[
 Q=-\log(L_1+Q).
\]
That also works after proper normalization, but \cref{eq:W-fixed-point} has already removed the large correction $-L_2$ and is more strongly contractive.

\section{Coefficient extraction}

Set
\[
 r=L_1^{-1},
 \qquad \ell=L_2,
 \qquad
 u=\sum_{n\ge1}p_n(\ell)r^n.
\]
Equation \cref{eq:W-fixed-point} becomes
\[
 u=-\log\bigl(1-r(\ell-u)\bigr)
 =\sum_{k\ge1}\frac{r^k(\ell-u)^k}{k}.
\]
Equate powers of $r$.

At order $r$,
\[
 p_1=\ell.
\]
At order $r^2$,
\[
 p_2=-p_1+\frac{\ell^2}{2}
 =\frac{\ell(-2+\ell)}2.
\]
Continuing gives
\begin{align*}
 p_3(\ell)&=\frac{\ell(6-9\ell+2\ell^2)}6,\\
 p_4(\ell)&=\frac{\ell(-12+36\ell-22\ell^2+3\ell^3)}{12},\\
 p_5(\ell)&=\frac{\ell(60-300\ell+350\ell^2-125\ell^3+12\ell^4)}{60}.
\end{align*}
Thus
\begin{align}
 W(x)\asym{}&L_1-L_2+\frac{L_2}{L_1}
 +\frac{L_2(-2+L_2)}{2L_1^2}\notag\\
 &+\frac{L_2(6-9L_2+2L_2^2)}{6L_1^3}\notag\\
 &+\frac{L_2(-12+36L_2-22L_2^2+3L_2^3)}{12L_1^4}\notag\\
 &+\frac{L_2(60-300L_2+350L_2^2-125L_2^3+12L_2^4)}{60L_1^5}
 +\cdots.
\label{eq:W-expansion}
\end{align}
This is a logarithmic-power transseries: each power of $1/L_1$ carries a polynomial in $L_2$.

\section{Why the recursion stabilizes}

Let
\[
 \Phi(u)=-\log\left(1-\frac{L_2-u}{L_1}\right).
\]
Then
\[
 \Phi'(u)=-\frac{1/L_1}{1-(L_2-u)/L_1}.
\]
Formally, $\Phi'(u)\asym-L_1^{-1}\precdom1$. Therefore a difference between two approximations gains a factor of order $1/L_1$ after one iteration. The coefficient of $L_1^{-N}$ eventually becomes permanent.

Starting with $u_0=0$:
\begin{align*}
 u_1&=-\log(1-L_2/L_1),\\
 u_2&=-\log\left(1-\frac{L_2-u_1}{L_1}\right),\\
 &\ \vdots
\end{align*}
Each iterate contains a longer correct initial segment of \cref{eq:W-expansion}.

\section{A quick formal verification}

Differentiate \cref{eq:W-definition}:
\[
 W'\e^W+W\e^W W'=1.
\]
Since $W\e^W=x$,
\begin{equation}\label{eq:W-derivative}
 W'(x)=\frac{W(x)}{x(1+W(x))}.
\end{equation}
Substituting the first terms
\[
 W=L_1-L_2+\bigO(L_2/L_1)
\]
gives
\[
 W'=\frac1x\left(1-\frac1{L_1}
 +\frac{1-L_2}{L_1^2}+\cdots\right),
\]
consistent with termwise differentiation of \cref{eq:W-expansion}.

\section{Accuracy and truncation}

Unlike the Euler factorial series, the classical large-$x$ expansion of $W$ in inverse powers of $L_1$ has a more delicate convergence/asymptotic story depending on the domain and organization of terms. For transseries purposes, the safe statement is formal-asymptotic: for fixed truncation order, the residual in \cref{eq:W-log-equation} moves to a smaller logarithmic scale.

At $x=\e^{100}$,
\[
 L_1=100,
 \qquad L_2=\log100\approx4.60517.
\]
Already
\[
 L_1-L_2+L_2/L_1
\]
provides a strong approximation because $L_2/L_1\approx0.046$.

\section{The real branch \texorpdfstring{$W_{-1}$}{W-1} near zero}

For $-\e^{-1}<x<0$, Lambert $W$ has two real branches. The branch $W_{-1}(x)$ tends to $-\infty$ as $x\to0^-$. Write
\[
 x=-\e^{-\eta},
 \qquad \eta\to+\infty,
 \qquad W_{-1}(x)=-v.
\]
Then
\[
 v\e^{-v}=\e^{-\eta},
 \qquad
 v-\log v=\eta.
\]
Let $L=\log\eta$. Reversing the same logarithmic calculation gives
\begin{align}\label{eq:Wminus1-expansion}
 W_{-1}(-\e^{-\eta})\asym{}&-\eta-L-\frac{L}{\eta}
 +\frac{L(-2+L)}{2\eta^2}\\
 &-\frac{L(6-9L+2L^2)}{6\eta^3}+\cdots.
\notag
\end{align}
The two real branches thus share the same universal polynomial pattern after the appropriate large logarithm and signs are chosen.

\section{The branch point uses a different scale}

Near $x=-\e^{-1}$, the derivative \cref{eq:W-derivative} blows up because $1+W=0$. The correct local variable is a square root, not an inverse logarithm. Set
\[
 p=\sqrt{2(1+\e x)}.
\]
Then the two branches have Puiseux expansions
\begin{equation}\label{eq:W-branchpoint}
 W_{0,-1}(x)
 =-1\pm p-\frac{p^2}{3}
 \pm\frac{11p^3}{72}-\frac{43p^4}{540}+\cdots.
\end{equation}
This is a useful reminder: the dominant balance determines the correct transseries scale. A logarithmic expansion suited to infinity is not transported unchanged to a branch point.

\section{A general inversion recipe}

For a monotone large transseries $F$:

\begin{algorithm}[Formal compositional inversion]
\begin{enumerate}
\item Find a leading inverse $G_0$ from dominant balance.
\item Write $G=G_0+U$ with $U\precdom G_0$.
\item Substitute into $F(G)=x$.
\item Isolate the dominant linear occurrence of $U$ and rewrite as $U=\Phi(U)$.
\item Iterate until the desired dominance threshold stabilizes.
\item Verify by composing the truncated answer back into $F$.
\end{enumerate}
\end{algorithm}

Lambert $W$ is the prototype because the inverse necessarily mixes two logarithmic levels.

\section{Exercises}

\begin{exercise}
Starting from \cref{eq:W-fixed-point}, derive $p_3(\ell)$ explicitly.
\end{exercise}

\begin{exercise}
Substitute the expansion through $L_1^{-2}$ into $W+\log W-L_1$ and show that the residual is of order $L_2^3/L_1^3$.
\end{exercise}

\begin{exercise}
Derive the first two corrections in \cref{eq:Wminus1-expansion} directly from $v-\log v=\eta$.
\end{exercise}

\begin{summarybox}
Lambert $W$ is found by peeling away the dominant logarithm and the next iterated logarithm, then solving a contractive fixed-point equation in $1/\log x$. Its coefficient polynomials illustrate how transseries naturally combine several nested scales. Different singular points require different scale hierarchies.
\end{summarybox}


\chapter{Formal fixed points and coefficient stabilization}\label{ch:fixed-points}

Fixed-point iteration in transseries resembles numerical iteration only superficially. The relevant notion of closeness is dominance: an error becomes smaller when its leading monomial moves downward in the support order.

\section{A valuation contraction principle}

Let $\Phi$ act on a class of transseries. Suppose there is a small monomial $\mu\precdom1$ such that
\begin{equation}\label{eq:formal-contraction}
 \Phi(S)-\Phi(T)\preceqdom \mu(S-T)
\end{equation}
for all $S,T$ in the class. Then each iteration gains at least one factor of $\mu$.

\begin{theorem}[Formal contraction principle, schematic form]\label{thm:formal-contraction}
Under suitable well-based support hypotheses, if \cref{eq:formal-contraction} holds and $T_{n+1}=\Phi(T_n)$, then the differences $T_{n+1}-T_n$ move to successively smaller monomials. The sequence converges coefficientwise to a unique fixed point $T=\Phi(T)$ in the class.
\end{theorem}

The rigorous versions use ratio sets, valuations, or complete ultrametric spaces. For calculation, the operational test is simple: every pass must push the unresolved residual beyond the current truncation.

\section{The Catalan fixed point}

Consider
\begin{equation}\label{eq:Catalan-fixed}
 S=A+S^2,
 \qquad A\precdom1.
\end{equation}
Start with $S_0=0$:
\begin{align*}
 S_1&=A,\\
 S_2&=A+A^2,\\
 S_3&=A+A^2+2A^3+A^4,\\
 S_4&=A+A^2+2A^3+5A^4+6A^5+\cdots.
\end{align*}
The stable coefficients are the Catalan numbers:
\begin{equation}\label{eq:Catalan-series}
 S=A+A^2+2A^3+5A^4+14A^5+42A^6+\cdots.
\end{equation}
Algebraically,
\[
 S=\frac{1-\sqrt{1-4A}}2,
\]
and the branch is selected by $S\precdom1$.

For two small candidates $S,T$,
\[
 \Phi(S)-\Phi(T)=S^2-T^2=(S+T)(S-T),
\]
and $S+T\precdom1$, so the map is contractive.

\section{Residual correction}

A flexible method is to work with the residual. Suppose $P(Y)=0$ and $Y_0$ is an approximation. Write
\[
 Y=Y_0+U.
\]
Expand
\[
 P(Y_0+U)=P(Y_0)+P'(Y_0)U+\text{higher powers of }U.
\]
If $P'(Y_0)$ has an inverse and the nonlinear remainder is smaller, isolate
\begin{equation}\label{eq:residual-fixed}
 U=-\frac{P(Y_0)}{P'(Y_0)}+\text{smaller terms involving }U.
\end{equation}
This is Newton's correction reorganized as a fixed-point equation.

\section{A square-root example}

Solve
\[
 Y^2=x^2+x+\e^{-x}
\]
with $Y>0$. Dominant balance gives $Y_0=x$. Write
\[
 Y=x(1+S).
\]
Then
\[
 (1+S)^2=1+x^{-1}+x^{-2}\e^{-x},
\]
so
\[
 2S+S^2=x^{-1}+x^{-2}\e^{-x}
\]
and
\begin{equation}\label{eq:sqrt-fixed}
 S=\frac12x^{-1}+\frac12x^{-2}\e^{-x}-\frac12S^2.
\end{equation}
Iteration gives
\begin{align*}
 S={}&\frac1{2x}-\frac1{8x^2}+\frac1{16x^3}
 -\frac5{128x^4}+\cdots\\
 &+\frac12x^{-2}\e^{-x}
 -\frac14x^{-3}\e^{-x}
 +\frac3{16}x^{-4}\e^{-x}+\cdots
 +\bigO(\e^{-2x}).
\end{align*}
The power sector is the binomial expansion of $\sqrt{1+x^{-1}}$; the exponential forcing generates its own correction sector.

\section{Fixed points for differential operators}

The integration equation \cref{eq:integration-fixed-point} has the form
\[
 U=\delta+\mathcal{A}(U),
 \qquad
 \mathcal{A}(U)=\delta U-U'/A.
\]
If both multiplication by $\delta$ and the operator $A^{-1}\partial$ make monomials smaller, then $\mathcal{A}$ is contractive. The solution is a Neumann series in the operator:
\[
 U=(I-\mathcal{A})^{-1}\delta
 =\delta+\mathcal{A}\delta+\mathcal{A}^2\delta+\cdots.
\]
This is the operator analogue of $1/(1-S)=1+S+S^2+\cdots$.

\section{Newton iteration versus plain iteration}

A fixed-point rearrangement can converge slowly in dominance order if it gains only one small ratio per step. Newton iteration often doubles the amount of correct support, because the new residual is quadratic in the old error:
\[
 Y_{n+1}=Y_n-\frac{P(Y_n)}{P'(Y_n)}.
\]
For formal power series, this underlies fast algorithms for reciprocal, square root, exponential, and reversion. The same idea extends to finite truncations of grid-based transseries, although support comparison and composition are more expensive.

\section{How an iteration fails}

An equation can be algebraically equivalent yet computationally useless. In the Lambert problem, writing
\[
 Q=\e^{-Q}-\log x
\]
and starting from $Q_0=0$ gives
\[
 Q_1=1-\log x,
 \qquad
 Q_2=x/\e-\log x,
\]
which moves to a larger scale instead of a smaller one. The map is not contractive in the intended region.

A good rearrangement isolates the variable through an operator whose derivative is small. This is why
\[
 u=-\log\left(1-\frac{L_2-u}{L_1}\right)
\]
works: differentiation with respect to $u$ introduces $1/L_1\precdom1$.

\section{A practical stopping rule}

Suppose the goal is an answer through threshold $\m_*$. Stop when
\[
 T_{n+1}-T_n\precdom\m_*.
\]
Then all terms larger than or equal to $\m_*$ have stabilized, assuming the contraction estimate persists. This is the formal analogue of stopping a numerical iteration when the error estimate falls below tolerance.

\begin{lstlisting}[style=pseudo,caption={Dominance-driven fixed-point iteration}]
input: map Phi, initial approximation T, monomial threshold cutoff
repeat:
    U := truncate(Phi(T), cutoff)
    if truncate(U - T, cutoff) = 0:
        return U
    T := U
\end{lstlisting}

The operation \texttt{truncate} is essential. Exact untruncated intermediates may be infinite and need not be materialized.

\section{Exercises}

\begin{exercise}
Use fixed-point iteration to solve $S=A(1+S)^2$ through $A^5$.
\end{exercise}

\begin{exercise}
Derive a contractive equation for the positive solution of $Y^3=x^3+x$ and compute the first four corrections.
\end{exercise}

\begin{exercise}
For \cref{eq:Catalan-fixed}, prove uniqueness among infinitesimal solutions by subtracting two fixed-point equations.
\end{exercise}

\begin{summarybox}
Formal fixed-point iteration converges when differences gain small monomial factors. Coefficients stabilize one dominance layer at a time. The art lies in choosing a rearrangement whose linearized operator is small; algebraic equivalence alone does not guarantee a useful iteration.
\end{summarybox}


\chapter{Fractional iteration and Abel's equation}\label{ch:fractional-iteration}

Composition permits integer iterates
\[
 T^{[0]}=x,
 \qquad
 T^{[1]}=T,
 \qquad
 T^{[2]}=T\circ T,
 \quad\ldots.
\]
A natural question is whether the discrete index can be replaced by a real parameter $s$ so that
\begin{equation}\label{eq:iteration-group}
 T^{[s+t]}=T^{[s]}\circ T^{[t]},
 \qquad
 T^{[1]}=T.
\end{equation}
Such a family is an iteration group.

\section{Abel's linearizing coordinate}

Suppose there is a transseries $V$ satisfying Abel's equation
\begin{equation}\label{eq:Abel}
 V\circ T=V+1.
\end{equation}
Then define
\begin{equation}\label{eq:fractional-via-Abel}
 T^{[s]}=V^{-1}\circ(V+s).
\end{equation}
Indeed,
\begin{align*}
 T^{[s]}\circ T^{[t]}
 &=V^{-1}\circ(V+s)\circ V^{-1}\circ(V+t)\\
 &=V^{-1}\circ(V+s+t)=T^{[s+t]}.
\end{align*}
Abel's coordinate turns composition by $T$ into ordinary translation by $1$.

\section{Translation}

For
\[
 T(x)=x+1,
\]
take $V(x)=x$. Then
\[
 T^{[s]}(x)=x+s.
\]
This is the simplest possible iteration group.

\section{Scaling}

Let
\[
 T(x)=ax,
 \qquad a>0,\quad a\ne1.
\]
Take
\[
 V(x)=\frac{\log x}{\log a}.
\]
Then
\[
 V(ax)=V(x)+1,
\]
so
\begin{equation}\label{eq:fractional-scaling}
 T^{[s]}(x)=a^s x.
\end{equation}
The familiar real power $a^s$ is exactly a fractional iterate of multiplication by $a$.

\section{Power maps}

Let
\[
 T(x)=x^a,
 \qquad a>0,\quad a\ne1.
\]
Now
\[
 V(x)=\frac{\log\log x}{\log a}
\]
satisfies
\[
 V(x^a)=\frac{\log(a\log x)}{\log a}=V(x)+1.
\]
Therefore
\begin{equation}\label{eq:fractional-power-map}
 T^{[s]}(x)=x^{a^s}.
\end{equation}
At $s=1/2$, the compositional square root of $x^a$ is $x^{\sqrt a}$:
\[
 (x^{\sqrt a})\circ(x^{\sqrt a})=x^a.
\]

\section{Near-identity maps and iterative logarithms}

If
\[
 T=x+\text{small correction},
\]
one may seek a formal vector field
\[
 \mathcal{V}=v(x)\frac{\dd}{\dd x}
\]
whose time-one flow is $T$. Formally,
\begin{equation}\label{eq:flow-exp}
 T=\exp(\mathcal{V})x
 =x+\mathcal{V}x+\frac{\mathcal{V}^2x}{2!}+\cdots.
\end{equation}
The operator $\mathcal{V}$ is called an iterative logarithm, and
\[
 T^{[s]}=\exp(s\mathcal{V})x.
\]
This is composition's analogue of writing $a^s=\exp(s\log a)$.

For example, if
\[
 T=x+a x^{-1}+\bigO(x^{-2}),
\]
then the leading vector field is
\[
 \mathcal{V}\asym a x^{-1}\partial_x,
\]
and the fractional iterate begins
\[
 T^{[s]}=x+sax^{-1}+\cdots.
\]
Higher coefficients must correct the difference between a single displacement and a flow.

\section{Shallow, moderate, and deep iteration}

The Taylor expansion used to compose $T$ with a perturbation depends on how the perturbation compares with logarithmic derivatives of the monomials. Edgar classifies cases as shallow, moderate, and deep \cite{edgar-fractional}. Roughly:

\begin{itemize}
\item In a shallow case, each composition correction is clearly smaller and a direct Taylor-recursive construction works.
\item In a moderate case, leading derivative effects are comparable and must be resummed or normalized.
\item In a deep case, a nominally small displacement can produce large relative changes in some monomials, so the naive Taylor iteration fails.
\end{itemize}

The classification is not a judgment of difficulty alone; it is a comparison of transseries scales.

\section{Why the exponential map is special}

For $T(x)=\e^x$, integer iteration raises exponential height:
\[
 x,\quad\e^x,\quad\e^{\e^x},\quad\ldots.
\]
A fractional iterate would be a form of tetration. It cannot be represented by simply interpolating coefficients inside one fixed finite-height sector. Standard logarithmic-exponential transseries are not designed to contain a globally satisfactory half-iterate of the exponential. Larger structures involving superexponentials, omega-series, surreal analysis, or other extensions are needed for such questions \cite{schmeling,berarducci-mantova-germs}.

\section{Conjugacy is the unifying idea}

Abel's equation is one instance of conjugating a complicated map to a simple normal form:
\[
 V\circ T\circ V^{-1}=x+1.
\]
Near a fixed point with multiplier $\lambda\ne0,1$, Schr\"oder's equation instead seeks
\[
 V\circ T=\lambda V,
\]
leading to
\[
 T^{[s]}=V^{-1}(\lambda^sV).
\]
Transseries provide the asymptotic language for constructing such conjugacies when ordinary Taylor series are insufficient.

\section{Exercises}

\begin{exercise}
Verify \cref{eq:fractional-power-map} directly from the group law.
\end{exercise}

\begin{exercise}
Find an Abel coordinate and fractional iterates for $T(x)=x+c$ with nonzero real $c$.
\end{exercise}

\begin{exercise}
Suppose $T=V^{-1}\circ(V+1)$. Show that replacing $V$ by $V+c$ does not change the resulting iteration group.
\end{exercise}

\begin{summarybox}
Fractional iteration is obtained by conjugating composition to translation or scaling. Abel's equation gives $T^{[s]}=V^{-1}(V+s)$. Simple maps yield familiar formulas, while the exponential map points beyond finite-height logarithmic-exponential transseries toward tetration and larger transserial fields.
\end{summarybox}

\part{Solving equations with transseries}

\chapter{Dominant balance and transseries Newton methods}\label{ch:dominant-balance}

A zero sum cannot have exactly one uncancelled largest monomial. This elementary observation is the transseries version of dominant balance and the starting point for solving algebraic and differential equations.

\section{The largest terms must cancel}

Suppose
\[
 T_1+T_2+\cdots+T_r=0.
\]
If one $T_j$ were strictly larger than all the others, its leading monomial could not be canceled. Therefore at least two summands must share the maximal dominance scale, and their leading coefficients must add to zero.

For a polynomial equation
\[
 P(Y)=a_0+a_1Y+\cdots+a_dY^d=0,
\]
we guess a leading monomial $Y\asym c\m$. Each term has leading monomial
\[
 \lm(a_j)\m^j.
\]
A viable guess makes the maximum among these monomials occur at least twice. This is the valuation-theoretic form of the Newton polygon rule.

\begin{algorithm}[Dominant balance]
\begin{enumerate}
\item List the monomial scale of every term under a trial $Y\asym c\m$.
\item Equate pairs of candidate dominant scales to determine $\m$.
\item Reject candidates for which another term becomes strictly larger.
\item Solve the leading coefficient equation for $c$.
\item Write $Y=c\m(1+S)$ with $S\precdom1$ and derive a fixed-point equation for $S$.
\end{enumerate}
\end{algorithm}

\section{A cubic with an exponential coefficient}

Consider
\begin{equation}\label{eq:cubic-exp}
 Y^3-xY-\e^x=0.
\end{equation}
The dominant exponential suggests balancing $Y^3$ with $\e^x$:
\[
 Y\asym\e^{x/3}.
\]
Then
\[
 xY\asym x\e^{x/3}\precdom\e^x,
\]
so the third term is indeed smaller. The real leading coefficient satisfies $c^3=1$, giving $c=1$.

Set
\[
 Y=\e^{x/3}(1+u),
 \qquad
 \eps=x\e^{-2x/3}\precdom1.
\]
After division by $\e^x$, \cref{eq:cubic-exp} becomes
\begin{equation}\label{eq:cubic-u}
 (1+u)^3-1=\eps(1+u).
\end{equation}
Expand:
\[
 3u+3u^2+u^3=\eps+\eps u.
\]
Solving recursively gives
\begin{equation}\label{eq:cubic-u-series}
 u=\frac{\eps}{3}-\frac{\eps^3}{81}
 +\frac{\eps^4}{243}-\frac{4\eps^6}{6561}
 +\frac{5\eps^7}{19683}+\cdots.
\end{equation}
Therefore the real root has transseries
\begin{align}\label{eq:cubic-root}
 Y={}&\e^{x/3}
 +\frac{x}{3}\e^{-x/3}
 -\frac{x^3}{81}\e^{-5x/3}
 +\frac{x^4}{243}\e^{-7x/3}\\
 &-\frac{4x^6}{6561}\e^{-11x/3}
 +\frac{5x^7}{19683}\e^{-13x/3}+\cdots.
\notag
\end{align}
The other two roots over $\C$ begin with $\omega\e^{x/3}$, where $\omega^3=1$ and $\omega\ne1$; their correction equations use the corresponding coefficient.

\section{How to obtain the coefficients}

Write
\[
 u=\sum_{n\ge1}a_n\eps^n.
\]
Substitute into \cref{eq:cubic-u} and compare powers. At order $\eps$,
\[
 3a_1=1,
\]
so $a_1=1/3$. At order $\eps^2$,
\[
 3a_2+3a_1^2=a_1,
\]
so $a_2=0$. At order $\eps^3$,
\[
 3a_3+6a_1a_2+a_1^3=a_2,
\]
so $a_3=-1/81$. The zero coefficients are genuine structural cancellations, not missing work.

\section{A richer fifth-degree example}

Edgar's introductory paper studies
\begin{equation}\label{eq:Edgar-quintic}
 Y^5+\e^xY^2-xY-9=0
\end{equation}
by exactly this method \cite{edgar-beginners}. Let $Y\asym ax^b\e^L$. The exponential parts of the four terms are
\[
 5L,
 \qquad x+2L,
 \qquad L,
 \qquad0.
\]
A dominant cancellation is possible only at
\[
 L=\frac{x}{3}
 \qquad\text{or}\qquad
 L=-\frac{x}{2}.
\]

For $L=x/3$, balancing $Y^5$ with $\e^xY^2$ gives
\[
 a^5+a^2=0,
 \qquad a^3=-1.
\]
There are three complex large roots; one is real with $a=-1$. Its beginning is
\begin{equation}\label{eq:Edgar-large-root}
 Y=-\e^{x/3}-\frac{x}{3}\e^{-x}
 +3\e^{-4x/3}+\frac{2x^2}{9}\e^{-7x/3}+\cdots.
\end{equation}
For $L=-x/2$, balancing $\e^xY^2$ with $9$ gives $a=\pm3$, producing two small real roots:
\begin{equation}\label{eq:Edgar-small-roots}
 Y=\pm3\e^{-x/2}+\frac{x}{2}\e^{-x}
 \pm\frac{x^2}{24}\e^{-3x/2}+\cdots.
\end{equation}
Together these balances account for all five roots. A seemingly difficult quintic becomes a collection of recursive small-correction problems.

\section{The Newton polygon picture}

Suppose the coefficient $a_j$ has a valuation or exponential order $v_j$. Plot the points
\[
 (j,v_j).
\]
A candidate root valuation corresponds to a slope of a supporting edge of the lower convex hull, with sign conventions depending on how the valuation is oriented. Points on the same edge represent terms that tie for dominance. Their coefficients form the leading polynomial for $c$.

\begin{figure}[htbp]
\centering
\begin{tikzpicture}[x=1.05cm,y=0.82cm,>=Latex]
  \draw[->] (-0.3,0)--(6.0,0) node[right] {$j$};
  \draw[->] (0,-0.3)--(0,5.6) node[above] {coefficient valuation};
  \coordinate (A) at (0,4.5);
  \coordinate (B) at (2,2.2);
  \coordinate (C) at (5,1.0);
  \coordinate (D) at (1,4.0);
  \coordinate (E) at (3,3.4);
  \draw[thick,DeepTeal] (A)--(B)--(C);
  \foreach \P/\lab in {A/$a_0$,B/$a_2$,C/$a_5$,D/$a_1$,E/$a_3$}{
    \fill[MediumBlue] (\P) circle (2pt);
    \node[above right] at (\P) {\lab};
  }
  \draw[dashed,gray] (D)--(A);
  \draw[dashed,gray] (E)--(B);
  \node[text width=3.2cm,align=center] at (4.3,4.8)
    {supporting edges encode possible dominant balances};
\end{tikzpicture}
\caption{A schematic Newton polygon. Only points on a supporting edge participate in a leading cancellation.}
\label{fig:newton-polygon}
\end{figure}

For transseries coefficients, a single real-valued valuation may not capture the full nested hierarchy, but the same ``maximum attained twice'' principle survives in the ordered monomial group.

\section{Newton correction after dominant balance}

Once $Y_0$ contains several terms, compute
\[
 R=P(Y_0).
\]
If $P'(Y_0)$ is nonzero, set
\[
 \Delta=-\frac{R}{P'(Y_0)}.
\]
Then
\[
 P(Y_0+\Delta)=\bigO(\Delta^2P''(Y_0)),
\]
so the residual often drops rapidly in dominance. This is an efficient way to extend a long expansion without re-deriving a fixed-point equation at every order.

\section{Degenerate balances}

If the leading coefficient polynomial has a multiple root, then $P'(Y_0)$ may be unusually small and ordinary Newton correction loses strength. New scales, including fractional powers or logarithms, can appear. Near a turning point or algebraic branch point, Puiseux series are the simplest example.

For differential equations, resonance can play the same role: the linearized operator fails to invert on a particular monomial, forcing logarithmic factors or additional free parameters.

\section{Exercises}

\begin{exercise}
Find all leading balances for
\[
 Y^4-\e^xY+x^2=0.
\]
How many real leading coefficients arise from each balance?
\end{exercise}

\begin{exercise}
Derive the coefficient of $\eps^4$ in \cref{eq:cubic-u-series}.
\end{exercise}

\begin{exercise}
For \cref{eq:Edgar-quintic}, verify that no value of $L$ other than $x/3$ or $-x/2$ can make the maximal exponential order occur twice.
\end{exercise}

\begin{summarybox}
A valid leading balance makes at least two maximal monomials cancel. After the leading monomial and coefficient are fixed, the equation becomes a small-correction fixed point or Newton iteration. The method is a transseries generalization of Newton polygons and Puiseux expansions.
\end{summarybox}


\chapter{Linear differential equations and WKB transseries}\label{ch:linear-odes}

Linear equations already exhibit divergent sectors, exponentially small complementary solutions, turning points, and Stokes phenomena. Their structure is simpler than that of nonlinear equations because sectors do not multiply into new sectors.

\section{First-order equations}

Consider
\begin{equation}\label{eq:first-order-linear}
 y'+a(x)y=b(x).
\end{equation}
Let
\[
 A(x)=\int a(x)\dd x.
\]
Formally,
\[
 (\e^A y)'=\e^A b,
\]
so
\begin{equation}\label{eq:first-order-solution}
 y=\e^{-A}\left(C+\int\e^A b\dd x\right).
\end{equation}
The homogeneous monomial $\e^{-A}$ is the scale on which the integration constant lives. A power-series particular solution may be blind to it, exactly as in Chapter~\ref{ch:euler}.

\begin{example}
For
\[
 y'+xy=1,
\]
we have $A=x^2/2$ and
\[
 y=\e^{-x^2/2}
 \left(C+\int\e^{x^2/2}\dd x\right).
\]
Formal integration gives
\[
 y\asym\frac1x+\frac1{x^3}+\frac3{x^5}
 +\frac{15}{x^7}+\cdots+C\e^{-x^2/2}.
\]
The coefficients are odd double factorials, and the missing constant lies on the scale $\e^{-x^2/2}$.
\end{example}

\section{Second-order normal form}

A broad class of equations can be transformed into
\begin{equation}\label{eq:WKB-normal}
 y''=Q(x)y.
\end{equation}
Set
\[
 y=\exp\left(\int p(x)\dd x\right).
\]
Then
\[
 y' =py,
 \qquad
 y''=(p'+p^2)y,
\]
so $p$ satisfies the Riccati equation
\begin{equation}\label{eq:WKB-Riccati}
 p'+p^2=Q.
\end{equation}
When $Q\succdom1$ is positive, dominant balance gives
\[
 p\asym\pm\sqrt Q.
\]
The two signs generate growing and decaying WKB sectors.

\section{The first WKB corrections}

Let $q=\sqrt Q$ and choose $s=\pm1$. Write
\[
 p=sq+r_0+r_1+\cdots,
\]
with descending corrections. Substitution into \cref{eq:WKB-Riccati} gives at the first correction
\[
 2sq r_0+sq'=0,
\]
so
\begin{equation}\label{eq:WKB-r0}
 r_0=-\frac{q'}{2q}=-\frac{Q'}{4Q}.
\end{equation}
At the next order,
\[
 2sq r_1+r_0'+r_0^2=0,
\]
which yields
\begin{equation}\label{eq:WKB-r1}
 r_1=s\left(
 \frac{Q''}{8Q^{3/2}}-
 \frac{5(Q')^2}{32Q^{5/2}}
 \right).
\end{equation}
Integrating $r_0$ produces the universal amplitude $Q^{-1/4}$. Thus
\begin{equation}\label{eq:WKB-leading}
 y_{\pm}(x)\asym
 Q(x)^{-1/4}
 \exp\left(\pm\int^x\sqrt{Q(t)}\dd t\right)
 \left(1+\text{smaller corrections}\right).
\end{equation}

\section{The Airy equation}

For
\begin{equation}\label{eq:Airy}
 y''=xy,
\end{equation}
we have $Q=x$ and
\[
 \int\sqrt{x}\dd x=\frac23x^{3/2}.
\]
The formal solutions are
\begin{align}\label{eq:Airy-WKB}
 y_{\pm}(x)\asym{}&x^{-1/4}
 \exp\left(\pm\frac23x^{3/2}\right)\\
 &\times\left(
 1\pm\frac5{48}x^{-3/2}
 +\frac{385}{4608}x^{-3}
 \pm\frac{85085}{663552}x^{-9/2}
 +\cdots\right).
\notag
\end{align}
Up to normalization, the decaying sign corresponds to $\Ai(x)$ and the growing sign to $\Bi(x)$ on the positive real axis.

The coefficient series diverge factorially. In the complex plane, which exponential is dominant depends on the argument of $x$, and the actual solution's decomposition changes across Stokes rays. Chapter~\ref{ch:stokes} returns to this example.

\section{Deriving the Airy recurrence}

Take
\[
 y=x^{-1/4}\e^{s\zeta}
 \sum_{n\ge0}a_nx^{-3n/2},
 \qquad
 \zeta=\frac23x^{3/2},
 \qquad s=\pm1,
 \qquad a_0=1.
\]
Substitute into $y''=xy$. Collecting powers gives a recurrence equivalent to
\begin{equation}\label{eq:Airy-recurrence}
 a_{n+1}=s\,
 \frac{(6n+1)(6n+5)}{48(n+1)}a_n.
\end{equation}
For the decaying solution $s=-1$, signs alternate; for the growing solution $s=+1$, they do not. The magnitude behaves roughly like a constant times $n!$, explaining divergence.

\section{Factorization of the differential operator}

If $p$ solves \cref{eq:WKB-Riccati}, then noncommutative operator multiplication gives
\begin{equation}\label{eq:operator-factorization}
 (\partial+p)(\partial-p)
 =\partial^2-p'-p^2
 =\partial^2-Q,
\end{equation}
where $\partial=\dd/\dd x$. Thus a transseries solution of the Riccati equation formally factors the second-order operator.

Conversely, constructing an asymptotic factorization is equivalent to constructing WKB logarithmic derivatives. This bridge is valuable in symbolic differential algebra.

\section{Turning points}

The WKB scale assumes $Q$ stays away from zero. If $Q(x_0)=0$, then $Q^{-1/4}$ blows up and the correction hierarchy collapses. Near a simple turning point, rescaling reduces the equation to Airy's equation. A local Airy approximation is then matched to outer WKB expansions.

Turning points are places where the correct transseries changes its monomial architecture. They also emit Stokes curves in the complex plane.

\section{Inhomogeneous equations}

For
\[
 y''-Qy=f,
\]
a particular transseries can be found by variation of parameters or by formally inverting the differential operator. If $y_+,y_-$ are homogeneous WKB solutions with Wronskian $W$, then
\[
 y_p=-y_+\int\frac{y_-f}{W}\dd x
 +y_-\int\frac{y_+f}{W}\dd x.
\]
Each integral may require its own asymptotic transseries. The full solution includes two independent constants multiplying $y_+$ and $y_-$.

\section{Exercises}

\begin{exercise}
Find the formal particular solution and homogeneous scale for
\[
 y'+x^2y=x.
\]
\end{exercise}

\begin{exercise}
Derive \cref{eq:WKB-r1} from the Riccati equation.
\end{exercise}

\begin{exercise}
Use \cref{eq:Airy-recurrence} to reproduce the first three correction coefficients in \cref{eq:Airy-WKB}.
\end{exercise}

\begin{summarybox}
Linear first-order equations attach the integration constant to a homogeneous exponential scale. Second-order equations lead through a Riccati substitution to WKB transseries with two exponential sectors. Their divergent coefficients and changing dominance in the complex plane prepare the Stokes phenomenon.
\end{summarybox}

\chapter{Nonlinear equations and the tower of sectors}\label{ch:nonlinear-odes}

A nonlinear equation does more than attach a single homogeneous exponential to a perturbative solution. Products of the first exponential correction generate its powers, and each power carries a new asymptotic series.

\section{A Riccati model}

Consider
\begin{equation}\label{eq:Riccati-model}
 y'+y=y^2+\frac1x.
\end{equation}
We first seek a perturbative inverse-power sector
\begin{equation}\label{eq:phi0}
 \phi_0(x)=\sum_{n=1}^{\infty}a_nx^{-n}.
\end{equation}
Substitution gives
\[
 \phi_0'+\phi_0=\phi_0^2+x^{-1}.
\]
The coefficient of $x^{-1}$ gives $a_1=1$. For $m\ge2$,
\begin{equation}\label{eq:phi0-recurrence}
 a_m=(m-1)a_{m-1}+\sum_{i=1}^{m-1}a_i a_{m-i}.
\end{equation}
Hence
\begin{equation}\label{eq:phi0-coeffs}
 \phi_0=x^{-1}+2x^{-2}+8x^{-3}+44x^{-4}
 +296x^{-5}+2312x^{-6}+\cdots.
\end{equation}
The rapid coefficient growth already suggests divergence.

\section{Linearize to discover the first exponential}

Write
\[
 y=\phi_0+\eta.
\]
Subtracting the equation for $\phi_0$ gives
\begin{equation}\label{eq:Riccati-eta}
 \eta'+(1-2\phi_0)\eta=\eta^2.
\end{equation}
Ignoring the quadratic term for an infinitesimal $\eta$, the linearized equation has solution
\[
 \eta\asym C\exp\left(-\int(1-2\phi_0)\dd x\right).
\]
Since $\phi_0\asym x^{-1}$,
\[
 -\int(1-2\phi_0)\dd x=-x+2\log x+\text{smaller terms}.
\]
Therefore the first beyond-all-orders scale is
\begin{equation}\label{eq:Riccati-first-scale}
 C\e^{-x}x^2\left(1+\text{inverse-power corrections}\right).
\end{equation}
The power $x^2$ is not guessed arbitrarily; it comes from integrating the $2/x$ correction in the linearized coefficient.

\section{The full one-parameter transseries}

Because the equation contains $y^2$, powers of the first correction interact. Use the ansatz
\begin{equation}\label{eq:Riccati-full-transseries}
 \trans y(x;C)=\sum_{k=0}^{\infty}C^k\e^{-kx}\phi_k(x),
\end{equation}
where each $\phi_k$ is an inverse-power or Laurent-type formal series. Substitute into \cref{eq:Riccati-model}. Differentiation gives
\[
 \frac{\dd}{\dd x}\left(\e^{-kx}\phi_k\right)
 =\e^{-kx}(\phi_k'-k\phi_k),
\]
and squaring convolves sector labels:
\[
 y^2=\sum_{k\ge0}C^k\e^{-kx}
 \sum_{i+j=k}\phi_i\phi_j.
\]
Thus every sector satisfies
\begin{equation}\label{eq:Riccati-sector-equation}
 \phi_k'+(1-k)\phi_k
 =\sum_{i+j=k}\phi_i\phi_j+\frac{\delta_{k0}}x.
\end{equation}
This single formula is the bookkeeping payoff of the transseries ansatz.

\section{The one-instanton sector}

For $k=1$, \cref{eq:Riccati-sector-equation} becomes
\[
 \phi_1'=2\phi_0\phi_1.
\]
Therefore
\[
 \phi_1=K\exp\left(2\int\phi_0\dd x\right).
\]
Absorb the normalization $K$ into $C$. Termwise integration of \cref{eq:phi0-coeffs} gives
\begin{equation}\label{eq:phi1}
 \phi_1=x^2\left(
 1-\frac4x-\frac8{x^3}-\frac{52}{x^4}
 -\frac{384}{x^5}-\frac{3200}{x^6}+\cdots
 \right).
\end{equation}
The coefficient of $x^{-2}$ inside the parentheses happens to vanish.

\section{The two-instanton sector}

For $k=2$,
\[
 \phi_2'-\phi_2=2\phi_0\phi_2+\phi_1^2.
\]
The dominant balance is
\[
 -\phi_2\asym\phi_1^2\asym x^4,
\]
so $\phi_2\asym-x^4$. Continuing gives
\begin{equation}\label{eq:phi2}
 \phi_2=x^4\left(
 -1+\frac6x-\frac6{x^2}+\frac8{x^3}
 +\frac{48}{x^4}+\frac{384}{x^5}+\cdots
 \right).
\end{equation}
The sector $C^2\e^{-2x}\phi_2$ was not present in the linearized equation. It is created by the square of the first sector.

\section{What the parameter means}

The first-order ODE has a one-dimensional family of local solutions. The transseries parameter $C$ is the formal coordinate on that family near the chosen asymptotic regime. It is determined by an initial condition or boundary condition after analytic reconstruction.

A Stokes constant is different. It is a property of the equation and normalization describing how the parameter $C$ changes when the summation direction crosses a Stokes ray. Schematically,
\[
 C\longmapsto C+S_1
\]
for a simple one-parameter problem. Confusing $C$ with $S_1$ obscures which data are solution-dependent and which are equation-dependent.

\section{A transasymptotic glimpse}

The leading behavior of \cref{eq:Riccati-full-transseries} is
\[
 y\approx\phi_0+
 z-z^2+z^3-\cdots,
 \qquad
 z=Cx^2\e^{-x},
\]
where the signs follow from the leading terms of $\phi_k$. The geometric pattern suggests
\begin{equation}\label{eq:Riccati-transasymptotic-leading}
 y\approx\phi_0+\frac{z}{1+z}
\end{equation}
when $z$ is not extremely small even though $\e^{-x}$ is. A denominator zero at $z=-1$ predicts the approximate location of movable poles:
\[
 Cx^2\e^{-x}\approx-1.
\]
This reorganization, which sums over sectors before expanding deeply within each sector, is called a \emph{transasymptotic} expansion. Chapter~\ref{ch:transasymptotics} develops the idea further.

\section{General sector algebra}

For an equation with polynomial nonlinearities, a product of a $k$-sector and an $\ell$-sector lies in sector $k+\ell$. Thus the nonnegative integers form the sector semigroup. For several independent exponential actions, the labels become multi-indices in $\N^r$.

If a cubic nonlinearity and a symmetry suppress even powers, one may obtain only odd sectors. If the linearized problem has several eigenvalues, each produces its own fundamental action. Resonance occurs when different combinations land on the same action.

\section{Divergence within every sector}

The outer sum over $k$ is only one infinite direction. Each $\phi_k$ usually has its own divergent large-order expansion:
\[
 \phi_k(x)=x^{\beta_k}\sum_{n\ge0}a_{k,n}x^{-n}.
\]
A resurgent transseries is therefore often doubly infinite. The large-order behavior of $a_{k,n}$ is controlled by neighboring sectors, and the Borel singularities of $\phi_k$ encode transitions to them.

\section{Exercises}

\begin{exercise}
Derive recurrence \cref{eq:phi0-recurrence} and compute $a_5$.
\end{exercise}

\begin{exercise}
Starting from $\phi_1'=2\phi_0\phi_1$, reproduce the coefficients $-4,0,-8$ in \cref{eq:phi1}.
\end{exercise}

\begin{exercise}
Use dominant balance in the $k$th sector equation to conjecture the leading power and sign of $\phi_k$.
\end{exercise}

\begin{summarybox}
Linearization finds the first exponential correction; nonlinearity generates all of its powers. Sector equations are obtained by convolution of sector labels. The free transseries parameter records the solution, while Stokes constants later describe direction-dependent changes of that parameter.
\end{summarybox}


\chapter{Singular perturbations and boundary layers}\label{ch:singular-perturbation}

Transseries also arise when a small parameter multiplies the highest derivative. An ordinary power series in that parameter often solves the differential equation in the bulk but cannot satisfy all boundary conditions. Exponentials such as $\e^{-x/\eps}$ repair the mismatch.

\section{A first boundary-layer equation}

Let $0<\eps\ll1$ and consider
\begin{equation}\label{eq:boundary-layer-model}
 \eps y'(x)+y(x)=x,
 \qquad y(0)=0.
\end{equation}
If one simply sets $\eps=0$, the reduced equation gives
\[
 y_0=x.
\]
This leading term happens to satisfy the boundary value at $x=0$, but that coincidence does not persist at the next algebraic order. To see the nonuniform correction, compute the regular expansion carefully. Write
\[
 y\asym y_0+\eps y_1+\eps^2y_2+\cdots.
\]
Substitution gives
\[
 y_0=x,
 \qquad y_1=-y_0'=-1,
 \qquad y_2=-y_1'=0,
\]
so the outer solution is
\[
 y_{\mathrm{out}}=x-\eps.
\]
This fails the boundary condition by $-\eps$.

The exact solution is
\begin{equation}\label{eq:boundary-layer-exact}
 y=x-\eps+\eps\e^{-x/\eps}.
\end{equation}
The exponentially small term for fixed $x>0$ is invisible to every power of $\eps$, yet at $x=0$ it is exactly large enough to restore the boundary condition.

\begin{warningbox}
A term can be exponentially small in the outer region and order one in a rescaled boundary layer. ``Small'' always refers to a specified asymptotic regime.
\end{warningbox}

\section{The inner variable}

Near $x=0$, introduce
\[
 X=\frac{x}{\eps}.
\]
Then $y'(x)=\eps^{-1}y_X$, and \cref{eq:boundary-layer-model} becomes
\[
 y_X+y=\eps X.
\]
At leading order in the layer,
\[
 y_X+y=0,
\]
whose solution is $A\e^{-X}=A\e^{-x/\eps}$. Matching to the outer discrepancy fixes $A=\eps$.

The inner and outer descriptions are two projections of the same transseries
\[
 y=x-\eps+\eps\e^{-x/\eps}.
\]

\section{General forcing}

For
\begin{equation}\label{eq:boundary-general}
 \eps y'+y=f(x),
\end{equation}
formal inversion of $1+\eps\partial$ gives
\begin{equation}\label{eq:outer-neumann}
 y_{\mathrm{out}}
 =(1+\eps\partial)^{-1}f
 =f-\eps f'+\eps^2f''-\eps^3f'''+\cdots.
\end{equation}
The homogeneous correction is
\[
 C\e^{-x/\eps}.
\]
Hence
\begin{equation}\label{eq:boundary-general-transseries}
 \trans y=f-\eps f'+\eps^2f''-\cdots+C\e^{-x/\eps}.
\end{equation}
For analytic $f$, the derivative coefficients often grow factorially, so even the outer series may diverge. The boundary layer and Borel divergence then coexist.

\section{Two boundary layers}

For a second-order problem such as
\[
 \eps y''+a(x)y'+b(x)y=f(x)
\]
on a finite interval, different characteristic roots may generate exponentials attached to the left and right boundaries. A term small from one end can grow toward the other. Boundary conditions select the allowed linear combination.

The sign of the action matters. For $\e^{-A(x)/\eps}$, a layer decays where $\Re A(x)>0$. Curves where $\Re A=0$ separate regions of growth and decay and are closely related to Stokes geometry.

\section{A nonlinear layer}

Consider
\begin{equation}\label{eq:nonlinear-layer}
 \eps y'=y-y^2,
 \qquad 0<y<1.
\end{equation}
The exact family is
\[
 y(x)=\frac{1}{1+C\e^{-x/\eps}}.
\]
For $x/\eps\to+\infty$,
\begin{equation}\label{eq:nonlinear-layer-transseries}
 y=1-C\e^{-x/\eps}+C^2\e^{-2x/\eps}
 -C^3\e^{-3x/\eps}+\cdots.
\end{equation}
This is the same nonlinear sector-generation mechanism seen in Chapter~\ref{ch:nonlinear-odes}, now with action $x/\eps$.

Near a point where $C\e^{-x/\eps}$ is order one, the sector expansion should be resummed into the closed denominator. That is a simple transasymptotic transition.

\section{Matched asymptotics versus transseries}

Classical matched asymptotic expansions construct separate outer and inner series and compare their overlap. A transseries aims to retain all relevant scales in one formal object. The approaches are complementary:

\begin{itemize}
\item matching is geometrically intuitive and often easiest for a small number of regions;
\item transseries give systematic algebra for multiple exponentially small effects, nonlinear sector interactions, and Stokes switching;
\item exponential asymptotics uses late terms of the outer expansion to determine the beyond-all-orders contributions needed for uniformity.
\end{itemize}

\section{Turning points and action integrals}

For WKB-type singular perturbations,
\[
 \eps^2y''=Q(x)y,
\]
the exponential actions are
\[
 A_{\pm}(x)=\pm\int^x\sqrt{Q(t)}\dd t,
\]
and the sectors contain $\e^{A_\pm/\eps}$. At a turning point $Q=0$, the two actions coalesce, ordinary WKB breaks down, and an Airy scaling resolves the local structure.

The action is the analogue of the monomial exponent $x$ in $\e^{-x}$. Its zeros, equalities, and complex phases organize the global asymptotic geometry.

\section{Exercises}

\begin{exercise}
Solve $\eps y'+y=1$ with $y(0)=0$ exactly and identify its outer and inner expansions.
\end{exercise}

\begin{exercise}
For \cref{eq:boundary-general}, determine $C$ in terms of a boundary value $y(0)=y_b$, using a truncation of the outer series.
\end{exercise}

\begin{exercise}
Expand the exact nonlinear solution \cref{eq:nonlinear-layer-transseries} about the other equilibrium $y=0$ in a region where $C\e^{-x/\eps}$ is large.
\end{exercise}

\begin{summarybox}
Singular perturbations force exponentials such as $\e^{-x/\eps}$ that are invisible to power series in $\eps$ but essential near boundaries. Inner scaling, matching, and transseries are compatible descriptions. Nonlinearity again generates towers of exponential sectors.
\end{summarybox}


\chapter{Several actions, resonance, and logarithmic sectors}\label{ch:multi-parameter}

Systems and higher-order equations usually carry more than one independent exponentially small mode. Their transseries parameters and actions become vectors, and arithmetic relations among actions create resonance.

\section{A multi-parameter template}

A common formal ansatz is
\begin{equation}\label{eq:multi-template}
 \trans F(x;\boldsymbol\sigma)
 =\sum_{\mathbf k\in\N^r}
 \boldsymbol\sigma^{\mathbf k}
 \e^{-\mathbf k\cdot\mathbf A x}
 x^{\mathbf k\cdot\boldsymbol\beta}
 \Phi_{\mathbf k}(x),
\end{equation}
where
\[
 \boldsymbol\sigma^{\mathbf k}
 =\sigma_1^{k_1}\cdots\sigma_r^{k_r},
 \qquad
 \mathbf k\cdot\mathbf A
 =k_1A_1+\cdots+k_rA_r.
\]
Each fundamental linearized mode contributes an action $A_j$ and a parameter $\sigma_j$. Polynomial nonlinearities add multi-indices:
\[
 \mathbf k+\boldsymbol\ell.
\]
Thus the sector labels form the additive semigroup $\N^r$.

\section{Linear systems}

For
\[
 \mathbf y'=M(x)\mathbf y+\mathbf f(x),
\]
the eigenvalues of the leading part of $M$ determine exponential factors
\[
 \exp\left(\int\lambda_j(x)\dd x\right).
\]
If $M$ is asymptotically diagonalizable, each eigenmode carries its own formal power corrections. Off-diagonal terms couple the recurrences but do not by themselves create sums of actions; nonlinear terms do.

If two eigenvalues coalesce, diagonalization may fail and polynomial or logarithmic prefactors arise, just as Jordan blocks produce factors of $x$ multiplying $\e^{\lambda x}$ in constant-coefficient systems.

\section{A simple resonance that forces a logarithm}

Consider
\begin{equation}\label{eq:log-resonance}
 y'+y=\frac{\e^{-x}}x.
\end{equation}
The homogeneous mode is $\e^{-x}$. The forcing lies in the same exponential sector. Multiplying by $\e^x$ gives
\[
 (\e^xy)'=\frac1x,
\]
so
\begin{equation}\label{eq:log-resonance-solution}
 y=\e^{-x}(C+\log x).
\end{equation}
The logarithm appears because the forcing resonates with the kernel of the linear operator $\partial+1$. Trying an ansatz $a(x)\e^{-x}$ reduces the problem to $a'=1/x$, whose antiderivative raises logarithmic depth.

\begin{intuitionbox}
Resonance means that the operator expected to determine a coefficient cannot be inverted within the originally guessed monomial class. The remedy is usually a generalized prefactor: a power of $x$, a logarithm, or an additional polynomial in a transseries parameter.
\end{intuitionbox}

\section{Commensurate actions}

Suppose a system has fundamental actions $A_1=A$ and $A_2=2A$. The sector $\e^{-2Ax}$ can arise in two ways:
\[
 \sigma_2\e^{-2Ax}
 \qquad\text{and}\qquad
 \sigma_1^2\e^{-2Ax}.
\]
Nonlinear equations mix these contributions. A clean normal form may require replacing the raw sector coefficient by a polynomial in parameters,
\[
 \sigma_2+c\sigma_1^2,
\]
or, in stronger resonance, by terms involving $x$ or $\log x$.

When actions satisfy no nontrivial integer relation, sector labels remain easier to separate. This is the nonresonant setting assumed in many rigorous Borel-summability theorems.

\section{Logarithmic transmonomials in resurgent ansatzes}

A resonant transseries may take the form
\begin{equation}\label{eq:log-sector-template}
 \sum_{\mathbf k}
 \boldsymbol\sigma^{\mathbf k}\e^{-\mathbf k\cdot\mathbf A/z}
 z^{\beta_{\mathbf k}}
 \sum_{m=0}^{M_{\mathbf k}}(\log z)^m
 \sum_{n\ge0}a_{\mathbf k,m,n}z^n.
\end{equation}
The logarithmic degree can grow with the sector index. Bookkeeping must then include three coordinates: action, logarithmic power, and perturbative order.

\section{Parameter counting}

An $n$th-order ODE has an $n$-dimensional local solution space under generic conditions. A transseries representation should contain the same number of independent free parameters, although asymptotic restrictions may set some to zero in a given sector. Nonlinear generation of infinitely many sectors does not create infinitely many independent constants; most sector coefficients are recursively determined by the finite parameter set.

This is a useful consistency check. If a first-order generic ODE appears to require two unrelated transseries parameters, one has probably double-counted a normalization or overlooked a relation.

\section{Bridge to Stokes transformations}

In a multi-parameter problem, crossing a Stokes ray acts on the parameter vector:
\[
 \boldsymbol\sigma
 \longmapsto
 \Stokes_{\theta}(\boldsymbol\sigma).
\]
The map can be triangular and nonlinear because sectors generate one another. Resonance can introduce polynomial and logarithmic terms in this parameter transformation. Alien calculus packages these changes as differential operators on transseries parameters; Chapter~\ref{ch:resurgence} gives a gentle introduction.

\section{Exercises}

\begin{exercise}
Solve
\[
 y'-2y=\e^{2x}x^{-1}
\]
and identify the resonant logarithmic factor.
\end{exercise}

\begin{exercise}
In a two-action transseries with actions $A$ and $\sqrt2A$, explain why distinct multi-indices cannot produce the same total action.
\end{exercise}

\begin{exercise}
A second-order nonlinear ODE has two independent linearized decaying modes. Write a schematic two-parameter sector ansatz through total sector degree $2$.
\end{exercise}

\begin{summarybox}
Several linearized modes lead to multi-parameter transseries indexed by $\N^r$. Nonlinear products add sector indices. Integer relations among actions create resonance, mixing sectors and often forcing powers or logarithms. The number of independent parameters still matches the differential equation's solution-space dimension.
\end{summarybox}

\part{From formal series to analytic functions}

\chapter{Borel transforms and Laplace reconstruction}\label{ch:borel}

A formal transseries can be manipulated without assigning a numerical value to it. Analysis begins when we ask a harder question:

\begin{quote}
Can a divergent formal sector be turned into an actual analytic function, in a way that respects the differential equation and recovers the formal expansion?
\end{quote}

For factorially divergent series, the basic answer is \emph{Borel--Laplace summation}. The idea is beautifully simple. Divide away the factorial growth, analytically continue the resulting power series, and then use a Laplace integral to put the factorials back.

This chapter develops that procedure slowly and fixes conventions carefully. Different books shift indices in slightly different ways; no mathematics depends on the choice, but every sign and factor does.

\section{Why ordinary convergence fails}

Consider a formal inverse-power series
\begin{equation}\label{eq:borel-formal-series}
 \trans f(x)=\sum_{n=0}^{\infty}a_nx^{-n-1},
 \qquad x\longrightarrow\infty.
\end{equation}
If $a_n$ grows approximately like $n!$, then for every fixed $x$ the terms eventually increase:
\[
 \left|\frac{a_{n+1}x^{-n-2}}{a_nx^{-n-1}}\right|
 \approx\frac{n+1}{|x|}.
\]
The smallest term occurs near $n\approx |x|$. Truncating there can give extraordinary accuracy, but the infinite series itself does not converge.

Factorial divergence is not random noise. It says that the coefficients have exactly the growth produced by repeated Laplace integration:
\begin{equation}\label{eq:laplace-moments}
 \int_0^{\infty}e^{-x\xi}\xi^n\dd\xi
 =\frac{n!}{x^{n+1}},
 \qquad \Re x>0.
\end{equation}
Equation~\eqref{eq:laplace-moments} is the entire motivation for the Borel transform.

\section{The formal Borel transform}

\begin{definition}[Borel transform]
For a formal series of the form \eqref{eq:borel-formal-series}, its order-one formal Borel transform is
\begin{equation}\label{eq:borel-definition}
 \Borel\trans f(\xi)
 :=\sum_{n=0}^{\infty}\frac{a_n}{n!}\xi^n.
\end{equation}
\end{definition}

If $a_n$ has at most factorial growth, the new coefficients may have merely geometric growth, so \eqref{eq:borel-definition} has a positive radius of convergence. A divergent series in $x^{-1}$ has become an ordinary analytic germ at $\xi=0$.

\begin{example}
For
\[
 \trans f(x)=\sum_{n\ge0}n!\,x^{-n-1},
\]
we obtain
\[
 \Borel\trans f(\xi)=\sum_{n\ge0}\xi^n=\frac1{1-\xi}
\]
near $\xi=0$. The factorial explosion has become a single pole at $\xi=1$.
\end{example}

The new variable $\xi$ is called the \emph{Borel variable}. The $\xi$-plane is often more informative than the original $x$-plane: singularities in the Borel plane reveal exponentially small scales in the original problem.

\section{The Laplace transform puts the factorials back}

Let $\theta$ be a direction in the Borel plane. Write
\[
 e^{\ii\theta}\R_+=\{re^{\ii\theta}:r\ge0\}.
\]
If a function $g(\xi)$ can be continued along this ray and grows no faster than exponentially, define its directional Laplace transform by
\begin{equation}\label{eq:directional-laplace}
 \Laplace_\theta g(x)
 :=\int_0^{e^{\ii\theta}\infty}e^{-x\xi}g(\xi)\dd\xi.
\end{equation}
The integral converges when $\Re(xe^{\ii\theta})$ is sufficiently positive.

Expanding $g(\xi)=\sum a_n\xi^n/n!$ near the origin and using \eqref{eq:laplace-moments} formally gives
\[
 \Laplace_\theta g(x)
 \asym\sum_{n\ge0}a_nx^{-n-1}.
\]
This motivates the central definition.

\begin{definition}[Directional Borel sum]
Suppose that $\Borel\trans f$ converges near $0$, admits analytic continuation along the ray of angle $\theta$, and has suitable exponential growth there. The directional Borel sum of $\trans f$ is
\begin{equation}\label{eq:borel-sum-definition}
 \lateralsum{\theta}\trans f
 :=\Laplace_\theta\Borel\trans f.
\end{equation}
\end{definition}

Three logically separate steps are hidden in this notation:
\begin{enumerate}
 \item the formal Borel transform must converge near $0$;
 \item that germ must be analytically continued far enough along the chosen ray;
 \item the continued function must satisfy a growth bound making the Laplace integral converge.
\end{enumerate}
A formal series can fail at any one of these stages.

\begin{intuitionbox}
The Borel transform converts \emph{large-order growth} into \emph{finite-distance singularities}. The Laplace transform converts those singularities back into exponentially small contributions. In this sense, the Borel plane is a microscope for effects that ordinary power asymptotics cannot see.
\end{intuitionbox}

\section{Gevrey order: measuring factorial growth}

\begin{definition}[Gevrey--1 series]
The formal series \eqref{eq:borel-formal-series} is Gevrey--1 if there are constants $C,A>0$ such that
\begin{equation}\label{eq:gevrey-one}
 |a_n|\le CA^n n!
 \qquad(n\ge0).
\end{equation}
\end{definition}

Condition \eqref{eq:gevrey-one} guarantees that the Borel transform converges for $|\xi|<A^{-1}$. More generally, Gevrey order $s$ corresponds roughly to growth like $\Gamma(1+sn)$. Ordinary analytic Taylor series have Gevrey order $0$; the most common divergent expansions from rank-one irregular differential equations have Gevrey order $1$.

A function $f$ on a sector has a Gevrey--1 asymptotic expansion $\sum a_nx^{-n-1}$ if its remainders satisfy a uniform factorial estimate of the form
\begin{equation}\label{eq:gevrey-remainder}
 \left|f(x)-\sum_{n=0}^{N-1}a_nx^{-n-1}\right|
 \le C A^N N!\,|x|^{-N-1}
\end{equation}
through every slightly smaller subsector. This estimate is much stronger than the ordinary Poincar\'e condition for each fixed $N$.

\section{The Euler series from beginning to end}

Return to the equation from \cref{ch:euler}:
\begin{equation}\label{eq:euler-borel-ode}
 y'+y=\frac1x.
\end{equation}
Its formal perturbative solution is
\begin{equation}\label{eq:euler-borel-series}
 \trans y_0(x)=\sum_{n=0}^{\infty}\frac{n!}{x^{n+1}}.
\end{equation}
The Borel transform is
\begin{equation}\label{eq:euler-borel-transform}
 \Borel\trans y_0(\xi)=\frac1{1-\xi}.
\end{equation}
Away from the positive real ray, there is no obstruction to a directional Laplace integral. Along the positive ray, however, the pole at $\xi=1$ lies directly on the path.

For $x>0$, define upper and lower lateral sums by paths passing just above or below the pole:
\begin{equation}\label{eq:euler-lateral-sums}
 \lateralsum{0^\pm}\trans y_0(x)
 :=\int_0^{\infty e^{\pm\ii0}}
 \frac{e^{-x\xi}}{1-\xi}\dd\xi.
\end{equation}
Their average is the Cauchy principal value,
\begin{equation}\label{eq:euler-principal-value}
 \frac12\left(\lateralsum{0^+}+\lateralsum{0^-}\right)\trans y_0
 =e^{-x}\Ei(x),
\end{equation}
with the standard real branch of $\Ei$ for $x>0$. The two lateral paths differ by a small loop around the pole, giving
\begin{equation}\label{eq:euler-stokes-jump}
 \lateralsum{0^+}\trans y_0-
 \lateralsum{0^-}\trans y_0
 =2\pi\ii e^{-x}.
\end{equation}
The ambiguity is not arbitrary. It is precisely a multiple of the missing homogeneous solution of \eqref{eq:euler-borel-ode}.

\begin{workedbox}
The exact family of solutions is
\[
 y(x)=e^{-x}\Ei(x)+Ce^{-x}.
\]
Represent the same analytic solution using either lateral sum:
\[
 y=\lateralsum{0^+}\trans y_0+C_+e^{-x}
  =\lateralsum{0^-}\trans y_0+C_-e^{-x}.
\]
By \eqref{eq:euler-stokes-jump},
\[
 C_+=C_--2\pi\ii.
\]
Thus crossing the singular direction changes the transseries coordinate, not the underlying solution.
\end{workedbox}

This example contains the whole elementary Stokes mechanism:
\begin{itemize}
 \item coefficient growth $n!$;
 \item a Borel singularity at the action $A=1$;
 \item an exponentially small scale $e^{-Ax}=e^{-x}$;
 \item two lateral reconstructions;
 \item a jump proportional to the homogeneous exponential;
 \item a compensating change of the free transseries parameter.
\end{itemize}

\section{The operational dictionary}

Borel transformation converts several operations into simple operations in the $\xi$-plane. Assume for the moment that all series have zero constant term in $x$ and use convention \eqref{eq:borel-definition}.

\subsection{Differentiation becomes multiplication}

From
\[
 \frac{\dd}{\dd x}\sum_{n\ge0}a_nx^{-n-1}
 =-\sum_{n\ge0}(n+1)a_nx^{-n-2},
\]
one finds
\begin{equation}\label{eq:borel-derivative}
 \Borel(\trans f')=-\xi\,\Borel\trans f.
\end{equation}
This is why differential equations often become algebraic or convolution equations after Borel transformation.

\subsection{Products become convolution}

Define
\begin{equation}\label{eq:borel-convolution}
 (g*h)(\xi):=\int_0^\xi g(\eta)h(\xi-\eta)\dd\eta.
\end{equation}
A beta-integral calculation gives
\begin{equation}\label{eq:borel-product}
 \Borel(\trans f\trans g)
 =(\Borel\trans f)*(\Borel\trans g).
\end{equation}
Indeed,
\[
 \int_0^\xi\frac{\eta^m}{m!}
 \frac{(\xi-\eta)^n}{n!}\dd\eta
 =\frac{\xi^{m+n+1}}{(m+n+1)!}.
\]
The shift by one matches multiplication of $x^{-m-1}$ and $x^{-n-1}$.

\subsection{Multiplication by \texorpdfstring{$x^{-1}$}{1/x} becomes integration}

One similarly obtains
\begin{equation}\label{eq:borel-xinverse}
 \Borel(x^{-1}\trans f)(\xi)
 =\int_0^\xi\Borel\trans f(\eta)\dd\eta.
\end{equation}

These rules let us translate many formal ODE recurrences into analytic integral equations. Convolution also explains why Borel singularities add: if $g$ is singular near $A$ and $h$ near $B$, then $g*h$ naturally develops singular behavior near $A+B$.

\section{Why the Laplace sum has the advertised asymptotic series}

Split the Laplace integral at a small positive $r$:
\[
 \int_0^{e^{\ii\theta}\infty}
 =\int_0^{re^{\ii\theta}}+\int_{re^{\ii\theta}}^{e^{\ii\theta}\infty}.
\]
On the short segment, expand
\[
 g(\xi)=\sum_{n=0}^{N-1}\frac{a_n}{n!}\xi^n+R_N(\xi).
\]
Integrating the polynomial gives the first $N$ inverse powers. Analyticity bounds $R_N(\xi)$ by a constant times $|\xi|^N$, producing a remainder of order $|x|^{-N-1}$. The far segment is exponentially small because of $e^{-x\xi}$. This is the basic proof behind Watson-type lemmas.

The exact hypotheses matter when one wants uniform Gevrey bounds, continuation over wide sectors, or uniqueness. For a first encounter, the important direction is:
\[
 \boxed{\begin{gathered}
  \text{analytic continuation and an exponential growth bound}\\[-0.2ex]
  \Downarrow\\[-0.2ex]
  \text{a function with the stated formal asymptotic series}
 \end{gathered}}
\]

\section{Uniqueness and exponentially flat functions}

An asymptotic expansion does not usually determine a function uniquely. The function $e^{-x}$ has the zero inverse-power expansion as $x\to+\infty$. More generally, exponentially small functions can be added without changing any coefficient.

On sufficiently wide complex sectors, analyticity and growth restrictions can restore uniqueness. A typical Watson uniqueness theorem says, roughly, that a Gevrey--1 expansion on a sector wider than $\pi$ determines the function uniquely within the appropriate growth class. The width $\pi$ is natural because $e^{-1/z}$ can decay on sectors of opening up to $\pi$ around the positive real $z$-axis.

This does not contradict Stokes phenomena. Different sectorial sums live on overlapping sectors separated by singular directions; each is unique in its own admissible summation problem, while their analytic continuations may differ by exponentially flat terms on the overlap.

\section{Other Gevrey orders and acceleration}

Not every problem has simple $n!$ growth. Suppose coefficients behave like $\Gamma(1+n/k)$. One can use a Borel transform of order $k$, followed by a matching Laplace transform. If several factorial scales occur in succession, \emph{multisummability} performs several transformations at different levels.

Still faster growth may require \emph{acceleration}: a nonlinear change of the Borel variable that turns excessive growth into manageable exponential growth. \`Ecalle's theory organizes these procedures into a flexible summation machine. Costin's generalized Borel summation gives a closely related analytic framework for broad classes of nonlinear differential systems \cite{costin-1995,costin-1998,costin-book}.

\section{Borel summation is not automatic}

The phrase ``take the Borel sum'' can conceal serious obstacles:
\begin{enumerate}
 \item the series may not be Gevrey of the expected order;
 \item the Borel transform may have a natural boundary preventing continuation;
 \item continuation may exist but grow too rapidly for an ordinary Laplace transform;
 \item singularities may lie on the desired direction, forcing lateral or median prescriptions;
 \item a complete transseries may require summing infinitely many sectors in a controlled topology.
\end{enumerate}

For particular differential equations, these are theorems to be proved, not formal consequences of writing down an ansatz.

\section{Exercises}

\begin{exercise}
Compute the Borel transforms of
\[
 \sum_{n\ge0}(-1)^n n!x^{-n-1}
 \quad\text{and}\quad
 \sum_{n\ge0}A^{-n-1}n!x^{-n-1}.
\]
Locate their nearest Borel singularities.
\end{exercise}

\begin{exercise}
Use the beta integral to prove \eqref{eq:borel-product} directly for two monomials and then extend it bilinearly.
\end{exercise}

\begin{exercise}
Apply the Borel transform to $y'+y=x^{-1}$ and solve the resulting algebraic equation for $\Borel\trans y$.
\end{exercise}

\begin{exercise}
Suppose $g$ is analytic near the positive ray and satisfies $|g(\xi)|\le Ce^{K|\xi|}$. Show that \eqref{eq:directional-laplace} converges for real $x>K$ when $\theta=0$.
\end{exercise}

\begin{summarybox}
Borel transformation divides out factorial growth; analytic continuation exposes finite Borel singularities; Laplace transformation reconstructs an analytic function with the original asymptotic expansion. In the Euler example, a pole at $\xi=1$ produces a jump proportional to $e^{-x}$. Summability requires analytic hypotheses beyond formal transseries algebra.
\end{summarybox}


\chapter{Stokes phenomena: when exponentially small terms switch}\label{ch:stokes}

An asymptotic formula can change its visible form as the argument of a complex variable rotates. A term that is tiny in one direction may become oscillatory on another and dominant on a third. The exact analytic function remains smooth; what changes is the efficient asymptotic description.

This direction-dependent rearrangement is the \emph{Stokes phenomenon}. Transseries are its natural bookkeeping language.

\section{A geometric picture of exponentials}

Consider
\begin{equation}\label{eq:stokes-exponential}
 e^{-A/z},
 \qquad z=re^{\ii\varphi},\quad A=|A|e^{\ii\alpha}.
\end{equation}
Its magnitude is
\begin{equation}\label{eq:stokes-magnitude}
 \left|e^{-A/z}\right|
 =\exp\left[-\frac{|A|}{r}\cos(\alpha-\varphi)\right].
\end{equation}
Therefore:
\begin{itemize}
 \item it is exponentially small where $\cos(\alpha-\varphi)>0$;
 \item it is exponentially large where $\cos(\alpha-\varphi)<0$;
 \item it is oscillatory with unit magnitude where $\cos(\alpha-\varphi)=0$.
\end{itemize}
The rays $\varphi=\alpha\pm\pi/2$ separate decay from growth.

Meanwhile, the Borel--Laplace integral
\[
 \int e^{-\xi/z}\Borel\trans f(\xi)\dd\xi
\]
meets a singularity at $\xi=A$ when its integration ray has angle $\alpha$. Thus the direction $\arg z=\arg A$ is where lateral Borel sums differ.

\begin{warningbox}
Authors do not all use the words \emph{Stokes line} and \emph{anti-Stokes line} in the same way. Some call the Borel-singular direction a Stokes line; others reserve ``Stokes line'' for equal phase and call equal magnitude anti-Stokes. The invariant content is the pair of equations
\[
 \Im(A/z)=0
 \quad\text{and}\quad
 \Re(A/z)=0.
\]
Always check which one an author is naming.
\end{warningbox}

\section{Lateral sums and discontinuous asymptotic coordinates}

Suppose $\Borel\trans f$ has singularities on the ray of angle $\theta$. Deform the Laplace path slightly above or below that ray:
\begin{equation}\label{eq:lateral-general}
 \lateralsum{\theta^\pm}\trans f
 :=\int_0^{e^{\ii(\theta\pm0)}\infty}
 e^{-\xi/z}\Borel\trans f(\xi)\frac{\dd\xi}{z}.
\end{equation}
Here we use the small-$z$ normalization, with $\dd\xi/z$, so that $\xi^n$ reconstructs $n!z^n$. This is equivalent to the large-$x$ convention of \cref{ch:borel} after $z=1/x$.

The difference
\begin{equation}\label{eq:disc-definition}
 \operatorname{Disc}_\theta\trans f
 :=\lateralsum{\theta^+}\trans f-
   \lateralsum{\theta^-}\trans f
\end{equation}
is exponentially small in a region where the original perturbative series is useful. The two sums have the same ordinary power expansion, so no finite algebraic-order calculation can detect the jump.

A full one-parameter transseries has the form
\begin{equation}\label{eq:stokes-full-transseries}
 \trans\Phi(z,\sigma)
 =\sum_{k\ge0}\sigma^k e^{-kA/z}
 z^{k\beta}\trans\phi_k(z).
\end{equation}
The same analytic solution can be represented on the two sides of a singular direction by different parameter values:
\begin{equation}\label{eq:stokes-parameter-match}
 \lateralsum{\theta^+}\trans\Phi(z,\sigma_+)
 =\lateralsum{\theta^-}\trans\Phi(z,\sigma_-).
\end{equation}
In the simplest normalization,
\begin{equation}\label{eq:simple-stokes-shift}
 \sigma_+=\sigma_-+S_\theta,
\end{equation}
where $S_\theta$ is a Stokes constant. Depending on the convention for the lateral discontinuity and the normalization of $\phi_1$, the sign may be reversed.

\section{The Euler jump revisited}

For the Euler series, the Borel singularity is the simple pole
\[
 \frac1{1-\xi}
\]
at $A=1$. The residue theorem gives
\[
 \lateralsum{0^+}\trans y_0-
 \lateralsum{0^-}\trans y_0=2\pi\ii e^{-x}.
\]
The Stokes multiplier is therefore $2\pi\ii$ in the normalization where the first nonperturbative sector is exactly $e^{-x}$. If that sector were normalized as $7e^{-x}$, the reported constant would be $2\pi\ii/7$. Stokes constants are meaningful only together with a stated sector normalization.

\section{A local singularity produces a whole asymptotic sector}

A simple pole produces a pure exponential. More complicated Borel singularities produce exponentials times powers and series.

Suppose near $\xi=A$ the analytic continuation has a branch behavior
\begin{equation}\label{eq:borel-local-branch}
 \Borel\trans\phi_0(\xi)
 \sim c(\xi-A)^{-\beta}
 \left(1+b_1(\xi-A)+b_2(\xi-A)^2+\cdots\right).
\end{equation}
The discontinuity across a cut from $A$ is supported near that point. Shift $\xi=A+u$ in the Laplace integral. The factor splits as
\[
 e^{-(A+u)/z}=e^{-A/z}e^{-u/z}.
\]
The resulting jump has the form
\begin{equation}\label{eq:borel-branch-sector}
 e^{-A/z}z^{1-\beta}
 \left(c_0+c_1z+c_2z^2+\cdots\right).
\end{equation}
Thus the local expansion around a Borel singularity is not merely an obstruction. It contains the fluctuation series of another transseries sector.

\section{Airy's equation as a map of Stokes geometry}

Airy's equation
\begin{equation}\label{eq:airy-stokes-equation}
 y''-zy=0
\end{equation}
has WKB exponentials
\begin{equation}\label{eq:airy-exponentials}
 \exp\left(\pm\frac23 z^{3/2}\right).
\end{equation}
Their magnitude ratio is controlled by
\[
 \Re\left(\frac43z^{3/2}\right).
\]
The rays of equal magnitude satisfy
\begin{equation}\label{eq:airy-equal-magnitude}
 \Re(z^{3/2})=0,
 \qquad
 \arg z=\frac{(2k+1)\pi}{3}.
\end{equation}
The rays of equal phase satisfy
\begin{equation}\label{eq:airy-equal-phase}
 \Im(z^{3/2})=0,
 \qquad
 \arg z=\frac{2k\pi}{3}.
\end{equation}
These six rays divide the plane into sectors in which one WKB exponential dominates the other.

For $|\arg z|<\pi$, one standard asymptotic form is
\begin{equation}\label{eq:airy-ai-asymptotic}
 \Ai(z)\asym
 \frac{e^{-\frac23z^{3/2}}}{2\sqrt\pi\,z^{1/4}}
 \sum_{n\ge0}(-1)^n\frac{u_n}{(\frac23z^{3/2})^n},
\end{equation}
inside suitable closed subsectors. After analytic continuation, a multiple of the opposite exponential appears in the efficient representation. The connection formulas among $\Ai(z)$ and its rotations are exact; the Stokes phenomenon describes how those exact formulas are reflected in asymptotic decompositions.

\begin{center}
\begin{tikzpicture}[scale=1.12]
  \draw[->,DarkGray] (-3.0,0)--(3.0,0) node[right] {$\Re z$};
  \draw[->,DarkGray] (0,-2.7)--(0,2.7) node[above] {$\Im z$};
  \foreach \a in {0,60,120,180,240,300}{
    \draw[MediumBlue,line width=0.8pt] (0,0)--(\a:2.45);
  }
  \foreach \a in {0,120,240}{
    \node[DeepBlue] at (\a:2.75) {$\Im(z^{3/2})=0$};
  }
  \foreach \a in {60,180,300}{
    \node[DeepTeal] at (\a:2.7) {$\Re(z^{3/2})=0$};
  }
  \fill (0,0) circle (1.2pt);
\end{tikzpicture}
\end{center}

The diagram is schematic: because $z^{3/2}$ is multivalued, one must specify branches when following a solution around the origin.

\section{Dominant and subdominant exponentials}

Suppose an exact solution is written asymptotically as
\[
 F(z)\asym F_{\mathrm{dom}}(z)
 +s(\varphi)F_{\mathrm{sub}}(z).
\]
When $F_{\mathrm{sub}}/F_{\mathrm{dom}}$ is exponentially tiny, changing $s$ by an order-one amount changes $F$ by less than every power of $z$. This is where an apparently discontinuous Stokes multiplier can hide.

After crossing an equal-magnitude ray, the former subdominant term may become dominant. Then its coefficient is no longer negligible and must have been chosen correctly. Stokes analysis supplies exactly that previously invisible information.

\section{Optimal truncation and the smooth reality behind the jump}

The classical asymptotic rule often states that a subdominant exponential ``switches on'' discontinuously at a Stokes ray. Exact analytic functions do not jump there. The discontinuity belongs to an idealized decomposition into an optimally truncated dominant series plus a separately named exponential remainder.

At finite but large $|z|^{-1}$, the effective multiplier changes rapidly but smoothly across a narrow angular transition region. Under broad conditions, the leading profile is an error function. This is \emph{Berry smoothing}. The angular width typically shrinks like the inverse square root of the large asymptotic parameter. Costin and Kruskal proved precise optimal-truncation and smoothing results for broad differential systems \cite{costin-kruskal}.

\begin{intuitionbox}
A Stokes jump is like the sharp limit of a rapidly changing coordinate chart. The analytic solution is the landscape; the transseries parameter is a coordinate. Crossing a special ray requires changing coordinates to keep describing the same point.
\end{intuitionbox}

\section{Median summation and real solutions}

If a formal series has real coefficients but a Borel singularity lies on the positive real ray, the two lateral sums are often complex conjugates for real positive $z$:
\[
 \overline{\lateralsum{0^+}\trans f}
 =\lateralsum{0^-}\trans f.
\]
Their average is real but may not by itself respect nonlinear multiplication. \`Ecalle's \emph{median summation} supplements lateral averaging by a half Stokes transformation. Schematically,
\begin{equation}\label{eq:median-schematic}
 \median
 =\lateralsum{0^-}\circ\Stokes_0^{-1/2}
 =\lateralsum{0^+}\circ\Stokes_0^{1/2}.
\end{equation}
The square root is understood as the exponential of half the infinitesimal Stokes operator. In favorable resurgent problems, median summation preserves reality and the algebraic operations needed by the equation.

For the elementary Euler pole, ordinary principal-value averaging already illustrates the real answer. For nonlinear problems, the transseries parameter and the full Stokes automorphism must be included.

\section{Global monodromy from local jumps}

Rotate $z$ through a full turn. One may cross several Stokes rays, applying a Stokes transformation at each. Their ordered product combines with ordinary monodromy from powers $z^\beta$ and logarithms. Exact global connection formulas impose relations among the Stokes constants.

For a linear equation, these transformations are represented by Stokes matrices. Their product, together with formal monodromy, gives the actual monodromy around the irregular singularity. For nonlinear equations, the transformations act nonlinearly on transseries parameters rather than linearly on a vector space of solutions.

\section{A practical procedure}

For a differential equation near an irregular limit:
\begin{enumerate}
 \item construct the perturbative formal series;
 \item linearize about it and identify exponential actions $A_j$;
 \item draw rays from the equations $\Re(A_j/z)=0$ and $\Im(A_j/z)=0$;
 \item Borel transform each formal sector;
 \item locate singularities at actions and their combinations;
 \item define directional sums away from singular rays;
 \item compare lateral sums and infer parameter transformations;
 \item use boundary, reality, or global monodromy conditions to fix the parameters.
\end{enumerate}
The hard steps are analytic continuation and singularity control. The geometric map tells us where to look.

\section{Exercises}

\begin{exercise}
For $A=2+2\ii$, find the rays on which $e^{-A/z}$ has unit magnitude and the ray on which a Borel singularity at $\xi=A$ obstructs radial Laplace summation.
\end{exercise}

\begin{exercise}
Assume $\Borel\trans f$ has a simple pole $r/(\xi-A)$ on a positive ray. Compute the difference of upper and lower lateral Laplace integrals, keeping track of the exponential $e^{-A/z}$.
\end{exercise}

\begin{exercise}
Derive the six rays \eqref{eq:airy-equal-magnitude}--\eqref{eq:airy-equal-phase} directly from $z=re^{\ii\varphi}$.
\end{exercise}

\begin{exercise}
Explain why changing the normalization of a one-instanton sector changes the numerical Stokes constant but not the analytic jump.
\end{exercise}

\begin{summarybox}
Borel singular directions force upper and lower lateral sums. Their difference is exponentially small and is absorbed by a transformation of transseries parameters. Equal-phase and equal-magnitude rays organize the geometry, though their names vary by convention. Optimal truncation reveals a smooth finite-parameter transition whose sharp limit is the classical Stokes jump.
\end{summarybox}

\chapter{Resurgence: how sectors remember one another}\label{ch:resurgence}

The Borel transform of the Euler series has one pole, and that pole predicts the missing exponential. In nonlinear problems the picture repeats recursively: after continuing around one singularity, one encounters more singularities, and their local expansions reproduce other sectors of the same transseries.

This self-referential analytic structure is called \emph{resurgence}. The word suggests that information from one sector ``rises again'' inside the large-order behavior of another.

\section{From Borel summability to resurgence}

Borel summability along one ray asks for analytic continuation only along that ray. Resurgence asks for much more controlled continuation throughout the Borel plane.

A useful informal definition is:

\begin{definition}[Resurgent germ, informal]
A convergent Borel germ at $\xi=0$ is resurgent if it can be analytically continued along every finite path that avoids a locally finite set of singular points, and if the singular behavior reached in this way remains of a controlled type.
\end{definition}

The precise category varies. \`Ecalle's endlessly continuable germs allow the set of forbidden points to depend on the length of the path but require only finitely many obstructions below every finite path length. More restrictive ``simple resurgent'' classes permit singularities built from poles and logarithms. For elementary reasoning, the essential properties are:
\begin{enumerate}
 \item singularities are isolated enough that one can travel around them;
 \item their locations have asymptotic meaning;
 \item their local expansions can themselves be continued;
 \item convolution preserves the class under suitable hypotheses.
\end{enumerate}

\section{Actions become Borel locations}

Suppose a transseries contains the exponential scale
\[
 e^{-A/z}.
\]
Then one expects a Borel singularity displaced by $A$. If several actions $A_1,\dots,A_r$ occur, nonlinear convolution generates sums
\begin{equation}\label{eq:action-semigroup}
 \mathbf k\cdot\mathbf A
 =k_1A_1+\cdots+k_rA_r,
 \qquad \mathbf k\in\N^r.
\end{equation}
These are exactly the sector actions in \cref{ch:multi-parameter}. The additive geometry of the Borel plane mirrors multiplication of transmonomials:
\[
 e^{-A/z}e^{-B/z}=e^{-(A+B)/z}.
\]

The singularity pattern need not be only the positive semigroup. Analytic continuation of different sectors may reveal differences of actions, conjugate actions, or repeated sheets of a ramified surface. Nevertheless, \eqref{eq:action-semigroup} is the first map to draw.

\section{Darboux's principle in the Borel plane}

Large Taylor coefficients are controlled by nearby singularities. This elementary fact is the bridge from finite-distance Borel geometry to factorial divergence.

Let
\[
 g(\xi)=\sum_{n\ge0}b_n\xi^n
\]
be analytic at $0$. If its nearest singularity is a simple pole
\begin{equation}\label{eq:local-simple-pole}
 g(\xi)\sim\frac{c}{A-\xi},
\end{equation}
then
\[
 \frac{c}{A-\xi}
 =\frac{c}{A}\sum_{n\ge0}\left(\frac\xi A\right)^n,
\]
so
\begin{equation}\label{eq:darboux-pole}
 b_n\sim cA^{-n-1}.
\end{equation}
If $g=\Borel\trans f$ and $a_n=n!b_n$, this becomes
\begin{equation}\label{eq:factorial-from-pole}
 a_n\sim c\,\frac{n!}{A^{n+1}}.
\end{equation}
The action $A$ controls the geometric factor and the singularity type controls the power of $n$ multiplying $n!$.

More generally,
\begin{equation}\label{eq:branch-darboux}
 g(\xi)\sim c\left(1-\frac\xi A\right)^{-\beta}
\end{equation}
implies
\begin{equation}\label{eq:branch-coefficients}
 b_n\sim
 \frac{c}{\Gamma(\beta)}A^{-n}n^{\beta-1},
 \qquad
 a_n\sim
 \frac{c}{\Gamma(\beta)}A^{-n}\Gamma(n+\beta).
\end{equation}
This is a version of Darboux's principle or singularity analysis.

\begin{intuitionbox}
The late coefficients of a divergent series are a blurred photograph of the nearest Borel singularities. The exponential action gives the distance and direction; the factorial-over-power correction gives the singularity's local shape.
\end{intuitionbox}

\section{Large-order relations between sectors}

Write a one-action transseries as
\begin{equation}\label{eq:resurgent-one-action}
 \trans\Phi(z,\sigma)
 =\sum_{k\ge0}\sigma^ke^{-kA/z}z^{k\beta}
 \trans\phi_k(z),
 \qquad
 \trans\phi_k(z)=\sum_{n\ge0}a_{k,n}z^n.
\end{equation}
A singularity of the Borel transform of $\phi_0$ at $A$ has local data related to $\phi_1$. Consequently, the late perturbative coefficients often satisfy an expansion of the schematic form
\begin{equation}\label{eq:large-order-master}
 a_{0,n}\sim
 \frac{S_{0\to1}}{2\pi\ii}
 \sum_{m\ge0}a_{1,m}
 \frac{\Gamma(n+\beta-m)}{A^{n+\beta-m}}.
\end{equation}
Truncating the right-hand sum after a few $m$ gives successive inverse-$n$ corrections.

Every symbol in \eqref{eq:large-order-master} depends on normalization conventions:
\begin{itemize}
 \item shifting the starting power of $z$ changes $\beta$;
 \item rescaling $\phi_1$ rescales $S_{0\to1}$ inversely;
 \item choosing upper-minus-lower rather than lower-minus-upper changes a sign;
 \item logarithmic singularities can add derivatives with respect to $\beta$ and hence $\log n$ terms.
\end{itemize}
The invariant claim is that the full low-order expansion of one sector governs the full late-order expansion of a neighboring sector.

\section{A direct contour explanation}

Cauchy's coefficient formula says
\[
 b_n=\frac1{2\pi\ii}\oint\frac{g(\xi)}{\xi^{n+1}}\dd\xi.
\]
For large $n$, deform the small circle outward. Contributions arise when the contour is caught on singularities. A singularity near $A$ contributes a factor $A^{-n}$, and its local discontinuity supplies an asymptotic expansion in inverse powers of $n$. After multiplying by $n!$, one obtains \eqref{eq:large-order-master}.

This contour argument is why large-order formulas are not mysterious numerical coincidences. They are coefficient extraction performed on a multiply continued Borel transform.

\section{Alien derivatives: singularity extractors}

Ordinary differentiation measures local change in the original variable. An \emph{alien derivative} measures singular behavior after analytic continuation to a specified Borel location.

We write
\[
 \alien_A
\]
for an operator associated with the singularity at action $A$. Very roughly, $\alien_A\trans\phi$ extracts the minor or discontinuity data of the Borel transform of $\trans\phi$ at $\xi=A$, translated back to the origin. The actual definition includes choices of continuation paths and averages over ways of bypassing earlier singularities.

Alien derivatives are designed to satisfy derivation-like rules:
\begin{equation}\label{eq:alien-leibniz}
 \alien_A(FG)=(\alien_AF)G+F(\alien_AG)
\end{equation}
for the appropriately normalized operator. They interact nontrivially with ordinary differentiation because the exponential weight $e^{-A/z}$ carries the action. A common pointed operator incorporates this weight so that commutation rules become simpler.

\begin{warningbox}
Alien calculus has several notational normalizations: $\Delta_A$, pointed $\dot\Delta_A$, logarithms of Stokes automorphisms, and lateral alien operators. Formulas copied between sources can differ by $e^{-A/z}$, by $2\pi\ii$, or by a sign. The structural relations are more stable than the typography.
\end{warningbox}

\section{The Stokes automorphism}

The total discontinuity along a ray collects all Borel singularities lying on that ray. Alien derivatives package this collection into a formal automorphism. Schematically,
\begin{equation}\label{eq:stokes-automorphism-schematic}
 \Stokes_\theta
 =\exp\left(
   \sum_{\substack{A\ne0\\\arg A=\theta}}
   e^{-A/z}\alien_A
 \right).
\end{equation}
With one standard convention, lateral summation satisfies
\begin{equation}\label{eq:lateral-stokes-relation}
 \lateralsum{\theta^+}
 =\lateralsum{\theta^-}\circ\Stokes_\theta.
\end{equation}
Thus a Stokes jump is the analytic shadow of an algebraic automorphism acting on formal transseries.

Because the exponent in \eqref{eq:stokes-automorphism-schematic} is a derivation, its exponential is an automorphism. This is the mechanism that allows Stokes continuation to respect multiplication and nonlinear equations.

\section{Bridge equations}

For a transseries solution of a differential equation, alien derivatives are often expressible as ordinary derivatives with respect to transseries parameters. In a simple one-parameter normalization, a bridge equation may read
\begin{equation}\label{eq:bridge-equation}
 e^{-A/z}\alien_A\trans\Phi(z,\sigma)
 =S_1\frac{\partial}{\partial\sigma}
 \trans\Phi(z,\sigma).
\end{equation}
It is called a bridge because it connects:
\begin{itemize}
 \item analytic continuation in the Borel plane, on the left;
 \item motion through the formal family of solutions, on the right.
\end{itemize}
Exponentiating \eqref{eq:bridge-equation} gives a translation of $\sigma$ and recovers the elementary shift \eqref{eq:simple-stokes-shift}.

In resonant or multi-parameter problems, the right side is a vector field in parameter space,
\[
 \sum_j S_{A,j}(\boldsymbol\sigma)
 \frac{\partial}{\partial\sigma_j},
\]
and the Stokes map can be triangular or nonlinear.

\section{The Euler example in alien language}

For
\[
 \Borel\trans y_0(\xi)=\frac1{1-\xi},
\]
the only finite singularity is at $A=1$. Its singular data are constant, corresponding to the one-instanton sector $e^{-x}$ with no fluctuation series. The alien derivative therefore maps the perturbative sector to a constant multiple of the exponential sector. Exponentiation yields the parameter shift by $2\pi\ii$ found directly in \cref{ch:borel}.

This example is unusually finite. In a nonlinear problem, the singularity at $A$ leads to the $k=1$ sector, the singularity at $2A$ to the $k=2$ sector, and analytic continuation of each sector leads both forward and backward through a network of sectors.

\section{Convolution builds the singularity lattice}

Suppose $g$ and $h$ have singularities at $A$ and $B$. In
\[
 (g*h)(\xi)=\int_0^\xi g(\eta)h(\xi-\eta)\dd\eta,
\]
the integration contour becomes pinched when $\eta$ approaches $A$ while $\xi-\eta$ approaches $B$. This occurs near $\xi=A+B$. Hence nonlinear products generate summed singularities.

Repeated convolution creates
\[
 A,\ 2A,\ 3A,\dots
\]
or the multi-action semigroup \eqref{eq:action-semigroup}. This analytic mechanism is exactly parallel to the sector convolution derived formally in \cref{ch:nonlinear-odes}.

\section{Resonance and logarithmic singularities}

If two different paths in action space reach the same point, their singular contributions mix. Integer relations among actions can produce logarithmic branching in the Borel plane. After Laplace transformation, such logarithms lead to powers of $\log z$ and polynomial dependence on sector indices.

This is the analytic counterpart of the formal resonance described in \cref{ch:multi-parameter}. It also complicates bridge equations: a single alien derivative can connect several monomials with the same total exponential weight.

\section{Numerically reading Stokes data from coefficients}

Suppose the leading large-order law is
\[
 a_n\sim C\,A^{-n}\Gamma(n+\beta).
\]
Then the ratio
\begin{equation}\label{eq:large-order-ratio}
 \frac{a_n}{a_{n-1}}
 \sim\frac{n+\beta-1}{A}
\end{equation}
estimates $A$ from a long coefficient list. After estimating $A$ and $\beta$, normalize
\begin{equation}\label{eq:stokes-sequence}
 C_n:=\frac{a_nA^{n+\beta}}{\Gamma(n+\beta)}.
\end{equation}
The sequence tends slowly to $C$, with corrections in powers of $1/n$. Richardson extrapolation or nonlinear sequence transformations can accelerate it. Subtracting the predicted contribution of the nearest singularity then reveals the next one.

A practical large-order investigation therefore proceeds hierarchically:
\begin{enumerate}
 \item inspect coefficient ratios to guess the closest action;
 \item divide out its factorial-over-power law;
 \item fit inverse-$n$ corrections from the neighboring sector;
 \item subtract this contribution;
 \item repeat for more distant actions.
\end{enumerate}
High-precision arithmetic is often essential because subleading singularities are exponentially suppressed in coefficient order.

\section{What resurgence does not claim}

Not every transseries is resurgent. A logarithmic-exponential Hahn series may be perfectly valid as a formal ordered object without any chosen Borel transform. Conversely, a resurgent object may involve fractional powers, logarithms, oscillatory actions, or other ingredients outside a narrow grid-based LE field.

It is useful to reserve three separate adjectives:
\begin{description}
 \item[formal transseries:] an element of a specified generalized-series algebra;
 \item[Borel summable:] reconstructible in a specified direction by a specified Borel--Laplace procedure;
 \item[resurgent:] having controlled repeated analytic continuation and singularity relations in the Borel plane.
\end{description}
None of the last two follows merely from the first.

\section{Exercises}

\begin{exercise}
Let $\Borel\trans f(\xi)=(1-\xi/A)^{-2}$. Compute its Taylor coefficients and the corresponding coefficients of $\trans f(z)$ under the convention $\Borel(\sum a_nz^{n+1})=\sum a_n\xi^n/n!$.
\end{exercise}

\begin{exercise}
Starting from \eqref{eq:branch-darboux}, use the generalized binomial theorem to derive \eqref{eq:branch-coefficients}.
\end{exercise}

\begin{exercise}
Explain geometrically why convolution of singularities at $A$ and $B$ can create a singularity at $A+B$.
\end{exercise}

\begin{exercise}
Assume \eqref{eq:large-order-ratio}. Construct an estimator for $A$ from $a_n/a_{n-1}$ when $\beta$ is known, and find its first relative error if one additional $c/n$ correction is present.
\end{exercise}

\begin{summarybox}
Resurgence turns Borel singularities into channels of communication among transseries sectors. Darboux analysis converts local singular behavior into factorial large-order laws. Alien derivatives extract singular data, the Stokes automorphism combines it along a ray, and bridge equations translate analytic continuation into motion in transseries-parameter space.
\end{summarybox}


\chapter{Transasymptotics, hyperasymptotics, and movable singularities}\label{ch:transasymptotics}

A transseries has at least two notions of order. In
\[
 \sum_{k\ge0}\sigma^ke^{-kA/z}\sum_{n\ge0}a_{k,n}z^n,
\]
we may move deeper in perturbative order $n$, or move across exponential sectors $k$. Ordinary asymptotics expands first in small $z$ at fixed $k$. Near certain boundaries that ordering becomes inefficient: many exponential sectors have comparable size and should be summed together.

A \emph{transasymptotic expansion} reorganizes the same information by summing across sectors at fixed combinations such as
\begin{equation}\label{eq:transasymptotic-variable}
 \tau=\sigma z^\beta e^{-A/z}.
\end{equation}
This often reveals nonlinear transitions and movable singularities that are invisible in any finite sector truncation.

\section{An exactly solvable logistic model}

Consider
\begin{equation}\label{eq:logistic-model}
 y'=-y+y^2,
 \qquad x\longrightarrow+\infty.
\end{equation}
The zero solution is the perturbative sector. Linearization gives the small exponential $e^{-x}$. Set
\[
 \tau=\sigma e^{-x}.
\]
Since $\tau'=-\tau$, the function
\begin{equation}\label{eq:logistic-exact}
 y(x)=\frac{\tau}{1+\tau}
\end{equation}
satisfies \eqref{eq:logistic-model}. Expanding at small $\tau$ gives
\begin{equation}\label{eq:logistic-transseries}
 y=\tau-\tau^2+\tau^3-\tau^4+\cdots
 =\sum_{k\ge1}(-1)^{k-1}\sigma^ke^{-kx}.
\end{equation}
The infinite tower of nonlinear sectors is simply the Taylor expansion of a rational function of the transasymptotic variable.

The exact solution has a movable pole when
\begin{equation}\label{eq:logistic-pole}
 1+\sigma e^{-x}=0.
\end{equation}
Its location depends on the integration constant $\sigma$, hence ``movable.'' No finite truncation of \eqref{eq:logistic-transseries} has a pole. The pole emerges only after summing all sectors.

\begin{intuitionbox}
Perturbative order describes fine structure inside one sector. Transasymptotic order describes collective behavior of many sectors. A singularity of the summed function can be encoded in perfectly regular coefficients spread over infinitely many sectors.
\end{intuitionbox}

\section{Reorganizing a two-dimensional array}

Suppose
\begin{equation}\label{eq:double-transseries}
 \trans y(z,\sigma)
 =\sum_{k\ge0}\tau^k
 \sum_{n\ge0}a_{k,n}z^n,
 \qquad \tau=\sigma z^\beta e^{-A/z}.
\end{equation}
Instead of summing over $n$ inside each $k$, collect equal powers of $z$:
\begin{equation}\label{eq:transasymptotic-reorganization}
 \trans y(z,\sigma)
 =\sum_{n\ge0}z^nF_n(\tau),
 \qquad
 F_n(\tau):=\sum_{k\ge0}a_{k,n}\tau^k.
\end{equation}
The functions $F_n$ are called transasymptotic coefficient functions. Substituting \eqref{eq:transasymptotic-reorganization} into the differential equation often yields a hierarchy of algebraic or differential equations in $\tau$:
\[
 \mathcal E_0(F_0)=0,
 \qquad
 \mathcal L_{F_0}F_1=R_1(F_0),
 \qquad\dots
\]
The leading $F_0$ frequently has poles or branch points. Those singularities predict where the original solution ceases to be regular.

\section{The Riccati model revisited}

For \cref{eq:Riccati-model}, the natural collective variable was
\[
 \tau=Cx^2e^{-x}.
\]
The leading sector coefficients suggested
\[
 \phi_k(x)\sim(-1)^{k-1}x^{2k}
 \qquad(k\ge1).
\]
Thus
\[
 \sum_{k\ge1}C^ke^{-kx}\phi_k(x)
 \sim\sum_{k\ge1}(-1)^{k-1}\tau^k
 =\frac{\tau}{1+\tau}.
\]
The approximate movable-pole condition
\begin{equation}\label{eq:riccati-pole-condition}
 1+Cx^2e^{-x}=0
\end{equation}
can be inverted with Lambert $W$. Taking a square root formally,
\[
 xe^{-x/2}=\pm(-C)^{-1/2},
\]
so
\begin{equation}\label{eq:riccati-poles-lambert}
 x=-2W_j\left(\mp\frac1{2\sqrt{-C}}\right),
\end{equation}
with branch choices $j$ producing arrays of approximate complex poles. Corrections from $F_1,F_2,\dots$ refine these locations.

This is a recurring pattern: a transseries derived at infinity predicts global singularities by solving an equation in the combined exponential variable.

\section{Painlev\'e I as a nonlinear prototype}

The first Painlev\'e equation is
\begin{equation}\label{eq:painleve-I}
 u''=6u^2+z.
\end{equation}
Its solutions have no movable branch points, but they do have movable poles. For large $z$ in suitable complex sectors, dominant balance gives
\begin{equation}\label{eq:painleve-leading}
 u_0(z)\sim\pm\sqrt{-\frac z6}.
\end{equation}
Substitution generates a formal algebraic-power expansion around either branch. Linearizing $u=u_0+\eta$ gives
\begin{equation}\label{eq:painleve-linearized}
 \eta''\sim12u_0\eta.
\end{equation}
A WKB estimate therefore produces exponentials whose actions are proportional to $(-z)^{5/4}$. Nonlinearity generates their integer multiples, yielding a one- or multi-parameter transseries depending on the sector and normalization.

Reorganizing those sectors in a variable
\[
 \tau=\sigma z^\beta
 \exp\bigl(-A(-z)^{5/4}\bigr)
\]
produces meromorphic coefficient functions $F_n(\tau)$. Their pole conditions give asymptotic arrays of Painlev\'e poles. This is one of the most striking uses of transasymptotics: data computed from a remote irregular limit predict the location of movable singularities of a nonlinear special function. Detailed normalization depends on the chosen branch of $(-z)^{1/4}$ and on the sector under study \cite{olde-daalhuis,costin-book}.

\section{When the transasymptotic variable is order one}

Even if $e^{-A/z}$ is tiny, $\sigma$ or $z^\beta$ may compensate it. More importantly, analytic continuation can move into regions where $\Re(A/z)$ decreases. The crossover occurs when
\begin{equation}\label{eq:tau-order-one}
 |\tau|=|\sigma z^\beta e^{-A/z}|\asymp1.
\end{equation}
At that point, truncating at $k=0$ or $k=1$ is no longer uniform. A rational or meromorphic $F_0(\tau)$ can remain useful across the transition.

This is analogous to a boundary layer. There one introduces a stretched coordinate such as $X=x/\eps$ because ordinary powers of $\eps$ cease to be uniform. Here one introduces an exponential coordinate because the hierarchy of sectors ceases to be uniform.

\section{Hyperasymptotics: expanding the remainder}

Optimal truncation of a divergent power series leaves a remainder of the size of the least term, often
\[
 e^{-A/z}.
\]
Instead of stopping, one can expand that remainder using the neighboring saddle or Borel singularity. The new expansion is itself usually divergent and can be optimally truncated, leaving a still smaller remainder such as $e^{-B/z}$ or $e^{-2A/z}$. Repeating this process gives \emph{hyperasymptotics}.

The hierarchy is:
\begin{enumerate}
 \item truncate the perturbative sector near its least term;
 \item represent the remainder by an adjacent exponential sector;
 \item truncate that sector near its own least term;
 \item continue through the network of actions.
\end{enumerate}
Hyperasymptotics emphasizes exponentially improved numerical approximation. Resurgence emphasizes the analytic relations that make the hierarchy possible. Transseries provide the common formal container.

\section{Terminants and smooth switching}

Near a Stokes ray, the remainder after optimal truncation is often expressed using an incomplete gamma function or a closely related \emph{terminant}. Its uniform asymptotics produce the error-function profile of Berry smoothing. The terminant therefore resolves three descriptions at once:
\begin{itemize}
 \item far from the ray, one exponential multiplier is nearly constant;
 \item inside the transition layer, it varies smoothly;
 \item in the infinite-parameter limit, the smooth profile approaches a step.
\end{itemize}
This is the analytic mechanism behind the informal phrase ``an exponential switches on.''

\section{Sector summation versus perturbative summation}

There are two different infinite sums to control:
\begin{equation}\label{eq:two-summations}
 \sum_{k\ge0}\sigma^ke^{-kA/z}
 \underbrace{\left(\sum_{n\ge0}a_{k,n}z^n\right)}_{\text{usually Borel summed}}
 \qquad\text{and}\qquad
 \underbrace{\sum_{k\ge0}\tau^kF_k}_{\text{sector sum}}.
\end{equation}
Borel summability of every inner series does not automatically prove convergence of the outer sector sum. Rigorous reconstruction requires uniform estimates in $k$, or a topology in which the whole transseries summation is continuous. Costin's rank-one theory is designed to supply such control in a broad nonresonant class \cite{costin-1995,costin-1998}.

\section{Practical pole prediction}

Given a computed transseries:
\begin{enumerate}
 \item factor each sector into a common variable $\tau$;
 \item compute the leading coefficients $a_{k,0}$ for many $k$;
 \item identify or Pad\'e-approximate $F_0(\tau)=\sum a_{k,0}\tau^k$;
 \item find its nearest singularities $\tau=\tau_*$;
 \item solve $\sigma z^\beta e^{-A/z}=\tau_*$ for $z$;
 \item use $F_1,F_2,\dots$ to perturb the predicted locations;
 \item compare against numerical integration of the original equation.
\end{enumerate}
The inversion in step 5 often reduces to Lambert $W$ at leading order, exactly as in \eqref{eq:riccati-poles-lambert}.

\section{Exercises}

\begin{exercise}
Solve \eqref{eq:logistic-model} by separation of variables and match its integration constant to $\sigma$ in \eqref{eq:logistic-exact}.
\end{exercise}

\begin{exercise}
Substitute $y=\sum_{n\ge0}z^nF_n(\tau)$, with $\tau=\sigma e^{-x}$ and $z=x^{-1}$, into a simple first-order equation of your choice. Derive the leading equation for $F_0$.
\end{exercise}

\begin{exercise}
Starting from \eqref{eq:riccati-pole-condition}, derive \eqref{eq:riccati-poles-lambert} and state where branch choices enter.
\end{exercise}

\begin{exercise}
For Painlev\'e I, use \eqref{eq:painleve-leading}--\eqref{eq:painleve-linearized} to show that the WKB exponent grows as a constant times $(-z)^{5/4}$.
\end{exercise}

\begin{summarybox}
Transasymptotics sums across exponential sectors by holding $\tau=\sigma z^\beta e^{-A/z}$ fixed. It can reveal movable poles hidden across infinitely many sectors. Hyperasymptotics instead repeatedly expands optimally truncated remainders. Both exploit the same action network exposed by resurgence.
\end{summarybox}


\chapter{What analytic reconstruction theorems actually say}\label{ch:analytic-status}

Formal algebra is exact within its own rules. Analytic reconstruction is conditional: it depends on the equation, the sector, growth bounds, nonresonance, and the chosen summation method. This chapter separates broadly established conclusions from tempting but invalid shortcuts.

\section{Four levels of assertion}

For a displayed object $\trans y$, ask which of the following is meant.

\begin{enumerate}
 \item \textbf{Formal existence.} The support is admissible and all algebraic operations are defined.
 \item \textbf{Formal solution.} Substitution into the formal equation gives zero coefficient by coefficient.
 \item \textbf{Asymptotic realization.} There is an actual function $y$ whose asymptotic expansion is $\trans y$ in a specified sector and scale.
 \item \textbf{Canonical reconstruction.} A stated summation procedure assigns such a function to $\trans y$ and respects the operations needed by the problem.
\end{enumerate}
Each implication requires a theorem. A formal solution need not automatically possess a function realization, and a realization need not be unique without sectorial conditions.

\section{A representative rank-one theorem}

A broad class studied by Costin consists, after normalization, of analytic systems near infinity of the schematic form
\begin{equation}\label{eq:rank-one-system}
 \mathbf y'
 =-\Lambda\mathbf y-\frac1xB\mathbf y
 +\mathbf f_0(x)+\mathbf g(x,\mathbf y),
\end{equation}
where $\Lambda$ is a constant matrix, $\mathbf f_0$ is small at infinity, and $\mathbf g$ is analytic and at least nonlinear in $\mathbf y$ after suitable separation of terms. Under nondegeneracy and nonresonance assumptions, one constructs formal transseries
\begin{equation}\label{eq:rank-one-transseries}
 \trans{\mathbf y}(x,\mathbf C)
 =\sum_{\mathbf k\in\N^r}
 \mathbf C^{\mathbf k}e^{-\mathbf k\cdot\boldsymbol\lambda x}
 x^{\mathbf k\cdot\boldsymbol\beta}
 \trans{\mathbf y}_{\mathbf k}(x).
\end{equation}
The rigorous conclusion, stated informally but faithfully, is:

\begin{theorem}[Generalized Borel reconstruction, schematic]\label{thm:costin-schematic}
For systems in the admissible rank-one class, the sector series in \eqref{eq:rank-one-transseries} are summable by a generalized Borel procedure in nonsingular directions where the selected exponentials decay. The summed transseries are actual solutions. Conversely, solutions with the corresponding decay behavior are represented by such summed transseries. Crossing singular directions changes the transseries parameters by Stokes maps.
\end{theorem}

The exact theorem includes hypotheses on eigenvalues, resonance, analyticity, and normalization. Its significance is the two-way correspondence: transseries are not merely guessed approximations; within the theorem's class they coordinatize actual solutions \cite{costin-1995,costin-1998,costin-borel-stokes}.

\section{Why nonresonance appears}

If eigenvalue combinations coincide,
\[
 \mathbf k\cdot\boldsymbol\lambda
 =\boldsymbol\ell\cdot\boldsymbol\lambda,
 \qquad\mathbf k\ne\boldsymbol\ell,
\]
sector equations couple at the same exponential weight. Denominators used to solve recurrences may vanish, logarithmic terms may appear, and parameter transformations may cease to be simple translations. Resonant theories exist, but the cleanest general theorems state explicit restrictions because the normal form and summability estimates change.

\section{The role of sector width}

The Laplace kernel decays only when
\[
 \Re(xe^{\ii\theta})>K
\]
for the growth constant $K$ of the continued Borel transform. Hence a Borel sum naturally lives in a sector of the original variable. Rotating that sector eventually makes the kernel grow or makes the path hit singularities.

Analytic continuation of the actual solution may extend farther, but its asymptotic representation must then change. This is why connection problems require several overlapping sectorial solutions and Stokes data.

\section{Quasianalyticity and uniqueness}

A class is quasianalytic if an element with zero asymptotic expansion must itself vanish. Ordinary Poincar\'e asymptotics on a narrow sector is not quasianalytic because exponentially flat functions survive. Strong Gevrey bounds on sufficiently wide sectors can be quasianalytic. Resurgent and analyzable classes obtain uniqueness by adding continuation and growth structure, not by ignoring the flat functions.

The lesson is subtle:
\begin{quote}
A divergent series can determine a function uniquely, but only after one specifies a sufficiently rigid analytic category and a direction or median prescription.
\end{quote}

\section{A comparison table}

\begin{longtable}{>{\raggedright\arraybackslash}p{0.25\textwidth}
                  >{\raggedright\arraybackslash}p{0.31\textwidth}
                  >{\raggedright\arraybackslash}p{0.34\textwidth}}
\toprule
Formal statement & Desired analytic statement & Extra work required\\
\midrule
\endfirsthead
\toprule
Formal statement & Desired analytic statement & Extra work required\\
\midrule
\endhead
$\trans f$ has well-based support & Infinite expression defines a germ & Evaluation or summation theorem; uniform control\\
$P(\trans f)=0$ formally & $P(f)=0$ analytically & Summation must commute with operations in $P$\\
$\trans f'=(\trans f)'$ & Differentiated sum has summed derivative & Uniform convergence or Borel operational theorem\\
Formal inverse exists & Analytic function is nonzero and inverse is represented & Sectorial lower bound and multiplicative summation\\
$\trans f\circ\trans g$ is formally defined & Composition of reconstructed functions matches it & Domain control and a composition theorem\\
A formal parameter $C$ is free & Every nearby actual solution has one $C$ & Completeness/converse theorem for the ODE\\
Borel singularity at $A$ & Stokes jump contains $e^{-A/z}$ & Analytic continuation and local singularity analysis\\
Large-order formula is guessed & Stokes constant is determined & Error estimates and exclusion/subtraction of other singularities\\
\bottomrule
\end{longtable}

\section{Formal composition is especially delicate}

Addition, multiplication, and differentiation often pass through Borel summation under standard hypotheses. Composition is harder because:
\begin{itemize}
 \item the inner function changes the asymptotic domain;
 \item infinitely many Taylor terms may contribute at the same magnitude;
 \item logarithms and exponentials can move singular directions;
 \item the formal support conditions and analytic convergence conditions are different.
\end{itemize}
For standard fields of transseries, composition is a formal operation under precise support hypotheses. Recent work establishes broad Taylor-expansion theorems over generalized power-series fields and monotonicity/Taylor approximation for omega-series \cite{bagayoko-mantova-taylor,mantova-monotonicity}. These results strengthen the formal-function dictionary, but they do not turn every arbitrary Hahn series into an analytic germ.

\section{No single universal summation map}

One might hope for a map
\[
 \Sigma:\{\text{all formal transseries}\}
 \longrightarrow\{\text{all functions}\}
\]
that preserves every operation and agrees with ordinary convergence. Such a map is too much to expect without restricting both sides. Problems include natural boundaries, incompatible branch choices, nonunique flat additions, excessive growth, and set-theoretic size.

Successful theories choose a structured domain:
\begin{itemize}
 \item multisummable series for systems with finitely many Gevrey levels;
 \item resurgent transseries with controlled Borel continuation;
 \item analyzable functions obtained by generalized Borel procedures;
 \item LE-transseries interpreted as germs in Hardy fields;
 \item omega-series evaluated at positive infinite surreal numbers.
\end{itemize}
Each answers a different version of ``what function does this formal expression represent?''

\section{Proof checklist for a claimed transseries solution}

When reading or writing a serious argument, check:
\begin{enumerate}
 \item \textbf{Ambient field:} Which monomials and supports are allowed?
 \item \textbf{Formal recursion:} Why is every coefficient or sector uniquely determined?
 \item \textbf{Growth:} What Gevrey or coefficient bounds are proved?
 \item \textbf{Continuation:} Where can the Borel transform be continued?
 \item \textbf{Summation direction:} Which rays are admissible or singular?
 \item \textbf{Operations:} Does summation commute with the equation's products, derivatives, and compositions?
 \item \textbf{Parameter completeness:} Does every relevant actual solution arise?
 \item \textbf{Stokes data:} How are lateral representations matched?
 \item \textbf{Uniformity:} Are sector sums controlled when infinitely many $k$ contribute?
 \item \textbf{Domain:} In exactly which $x$- or $z$-sector is the result valid?
\end{enumerate}
A paper need not prove every item from scratch, but it should invoke a theorem whose hypotheses cover them.

\section{Exercises}

\begin{exercise}
Give an example of two distinct functions with the same inverse-power asymptotic expansion on the positive real axis.
\end{exercise}

\begin{exercise}
Explain why a summation method that does not preserve multiplication cannot automatically reconstruct solutions of a nonlinear ODE from formal solutions.
\end{exercise}

\begin{exercise}
For the rank-one template \eqref{eq:rank-one-system}, explain how the eigenvalues of $\Lambda$ determine the first exponential actions.
\end{exercise}

\begin{exercise}
Classify each claim as formal, asymptotic, or summational: ``$\log(e^T)=T$,'' ``$f(x)\sim x^{-1}$,'' and ``the two lateral sums differ by $2\pi\ii e^{-x}$.''
\end{exercise}

\begin{summarybox}
Formal solvability, asymptotic realization, and canonical summation are different levels. Rigorous reconstruction theorems prove their connection for specified classes, directions, and hypotheses. The broad rank-one theory shows that generalized Borel-summed transseries can parametrize actual nonlinear ODE solutions, but no unrestricted universal summation principle exists.
\end{summarybox}

\part{Transseries in the larger mathematical landscape}

\chapter{Hardy fields: transseries as eventual germs}\label{ch:hfields}

Formal transseries speak the language of growth. Hardy fields speak about actual functions that eventually have stable behavior. The two theories meet because both carry an order, a derivative, and a comparison relation at infinity.

\section{Germs at positive infinity}

Two functions $f$ and $g$, each defined on some interval extending arbitrarily far to the right, have the same \emph{germ at $+\infty$} if they agree for all sufficiently large $x$. We forget every bounded initial segment.

For example, the functions
\[
 f(x)=\log x
 \quad\text{for }x>0
\]
and
\[
 g(x)=
 \begin{cases}
 1729,&0<x\le100,\\
 \log x,&x>100
 \end{cases}
\]
define the same germ. Their eventual behavior is identical.

Addition, multiplication, and differentiation are defined using representatives:
\[
 [f]+[g]=[f+g],
 \qquad
 [f][g]=[fg],
 \qquad
 [f]'=[f'].
\]
Changing a representative on a bounded interval does not affect these operations.

\begin{definition}[Hardy field]
A Hardy field is a field of germs at $+\infty$ of real-valued differentiable functions that is closed under differentiation.
\end{definition}

The word \emph{field} has strong consequences. Every nonzero element must possess a reciprocal, so a representative of a nonzero germ must be nonzero for all sufficiently large $x$.

\section{Why oscillation is excluded}

The germ of $\sin x$ cannot lie in a Hardy field. It has zeros arbitrarily far to the right, so $1/\sin x$ is not represented by a function defined on any final interval $(a,\infty)$. Similarly, a nonzero element of a Hardy field has an eventual sign.

This nonoscillation is one reason Hardy fields align so well with ordered transseries. A transseries has a leading monomial and therefore a definite sign when its leading coefficient is nonzero. Persistent oscillations such as $\sin x$ require a different asymptotic language, often involving complex phases, Fourier modes, or transseries over nonordered monomial groups.

\section{Basic examples}

The rational function field
\[
 \R(x)
\]
becomes a Hardy field by ordering $x$ as positive infinite and taking ordinary differentiation. Larger examples include fields generated by suitable collections of
\[
 x^r,\quad \log x,\quad \log\log x,\quad e^x,
 \quad e^{\sqrt{x}},\quad \text{and their algebraic combinations}.
\]
The germs definable in many tame, or \emph{o-minimal}, expansions of the real field form Hardy fields. Tameness prevents wild oscillation and guarantees eventual regularity.

Not every ad hoc set of familiar functions is a field or is closed under differentiation. ``The Hardy field generated by these functions'' means the smallest field of germs containing them and closed under the required derivative operations, when such a field is under discussion.

\section{Eventual order and dominance}

For germs $f$ and $g$, define
\[
 f<g
 \quad\Longleftrightarrow\quad
 f(x)<g(x)\text{ eventually}.
\]
Because nonzero Hardy-field germs have eventual sign, this is a total order.

Let $C$ be the constant field, usually identified with $\R$. Define
\begin{align}
 f\preceq g
 &\quad\Longleftrightarrow\quad
 |f|\le c|g|\text{ eventually for some }c\in C^{>0},
 \label{eq:hardy-dominance}\\
 f\precdom g
 &\quad\Longleftrightarrow\quad
 |f|<c|g|\text{ eventually for every }c\in C^{>0}.
 \label{eq:hardy-strict-dominance}
\end{align}
Thus $f\precdom g$ means $f/g\to0$ whenever the quotient has the expected limit interpretation.

Examples are
\[
 \log x\precdom x^a\precdom e^{x^b}\precdom e^x
 \qquad(0<a,b<1).
\]

The convex hull of the constants is
\begin{equation}\label{eq:hfield-valuation-ring}
 \mathcal O=\{f:f\preceq1\},
\end{equation}
the ring of bounded germs. Its maximal ideal
\begin{equation}\label{eq:hfield-infinitesimals}
 \mathfrak o=\{f:f\precdom1\}
\end{equation}
consists of infinitesimals. The quotient $\mathcal O/\mathfrak o$ records the eventual constant part when that part exists.

\begin{warningbox}
A general bounded function need not approach a constant. The axiom $\mathcal O=C+\mathfrak o$ is therefore a genuine tameness requirement: every bounded element must be a constant plus an infinitesimal.
\end{warningbox}

\section{The definition of an \texorpdfstring{$H$}{H}-field}

\begin{definition}[$H$-field]
An $H$-field is an ordered differential field $K$ with constant field $C$ such that:
\begin{enumerate}
 \item $\mathcal O=C+\mathfrak o$;
 \item if $f$ is larger than every constant, written $f>C$, then $f'>0$.
\end{enumerate}
\end{definition}

The second axiom says that a positive infinite element is increasing. For actual germs this is a deep eventual regularity property, not true of arbitrary functions tending to infinity. For example, $x+\sin(x^2)$ tends to infinity but its derivative changes sign infinitely often. Such a germ cannot live in a Hardy field containing its derivative and satisfying the $H$-field conditions.

Hardy fields containing $\R$ are $H$-fields, and the standard transseries field $\T$ is an $H$-field. The axioms isolate the order--derivative behavior shared by these apparently different objects \cite{adh-numbers,adh-book}.

\section{Leading monomials versus valuation classes}

In a transseries, the leading monomial $\lm(T)$ determines its dominance class. In an abstract $H$-field, one may not have named monomials, but the relation
\[
 f\asymp g
 \quad\Longleftrightarrow\quad
 f\preceq g\text{ and }g\preceq f
\]
plays the same role. The quotient group of nonzero elements by this relation is the \emph{value group}.

Very roughly:
\[
 \boxed{\text{transmonomial}}
 \quad\longleftrightarrow\quad
 \boxed{\text{chosen representative of a valuation class}}.
\]
A Hahn field has a distinguished cross-section of monomials. A general Hardy field need not.

\section{Logarithmic derivatives}

For $f\ne0$, the logarithmic derivative is
\begin{equation}\label{eq:logarithmic-derivative}
 f^\dagger:=\frac{f'}f.
\end{equation}
It converts multiplication into addition:
\[
 (fg)^\dagger=f^\dagger+g^\dagger.
\]
For typical scales,
\[
 (x^a)^\dagger=\frac ax,
 \qquad
 (e^{x^b})^\dagger=bx^{b-1},
 \qquad
 (e^{e^x})^\dagger=e^x.
\]
Thus logarithmic derivatives encode the relation between growth and differentiation. Much of asymptotic differential algebra is organized by how valuation classes transform under $f\mapsto f^\dagger$.

A key phenomenon is that differentiation may reverse naive comparisons. Although $1\precdom\log x\precdom x$, their derivatives are
\[
 0,\quad x^{-1},\quad1.
\]
The correct theory tracks both valuations and logarithmic derivatives, leading to the structure called an \emph{asymptotic couple}.

\section{Small derivation}

An $H$-field has \emph{small derivation} if
\begin{equation}\label{eq:small-derivation}
 f\precdom1\quad\Longrightarrow\quad f'\precdom1.
\end{equation}
For germs this says that the derivative of an infinitesimal is still infinitesimal. It holds in Hardy fields and in the standard transseries field.

The condition depends on the chosen asymptotic variable. Under a change of independent variable, the derivation rescales. In differential-algebra language one often replaces $\partial$ by $\phi\partial$ to reveal the correct dominant part of an equation.

\section{Liouville closure}

An $H$-field is \emph{Liouville closed} if it is real closed and, for every $f$, it contains solutions to
\begin{equation}\label{eq:liouville-closure}
 g'=f
 \qquad\text{and}\qquad
 h'=fh,
 \quad h\ne0.
\end{equation}
The second equation asks for an exponential integral:
\[
 h=c\exp\left(\int f\right).
\]
Thus Liouville closure supplies antiderivatives, exponentials of antiderivatives, and algebraic roots.

The full transseries field has these closure properties. Its formal integration algorithm from \cref{ch:integration} is therefore not merely a convenience; it reflects a structural completeness property.

\section{Newtonianity and differential equations}

Hensel's lemma says that an approximate root of a polynomial over a valued field can often be corrected to an exact root. \emph{Newtonianity} is an analogous property for differential polynomials. After rescaling the derivation and isolating a dominant differential part, a suitable approximate root can be lifted.

Another condition, called $\omega$-freeness, excludes a particular kind of pseudolimit obstruction arising from iterated logarithms. The combination of Liouville closure, newtonianity, and $\omega$-freeness gives a powerful axiomatization of the first-order theory of transseries. At a sophomore level, it is enough to remember the roles:
\begin{description}
 \item[Liouville closed:] elementary integrations and exponential integrations exist;
 \item[newtonian:] dominant differential equations admit corrected roots;
 \item[$\omega$-free:] no hidden logarithmic-gap obstruction remains.
\end{description}

\section{An intermediate-value principle for differential polynomials}

A differential polynomial has the form
\[
 P(Y)=P(Y,Y',\dots,Y^{(r)})
\]
with coefficients in the field. An $H$-field version of the intermediate value property says:
\[
 f<g,\qquad P(f)<0<P(g)
 \quad\Longrightarrow\quad
 \text{some }y\text{ between }f\text{ and }g\text{ satisfies }P(y)=0.
\]
For ordinary polynomials this characterizes real closed fields. For differential polynomials it captures a remarkable degree of closure under asymptotic differential equations.

The theory developed by Aschenbrenner, van den Dries, and van der Hoeven gives complete and model-complete axiomatizations in this setting; \cref{ch:model-theory} explains what those logical words mean \cite{adh-book,adh-numbers}.

\section{From transseries to germs}

A formal transseries $T$ can sometimes be realized by a germ $f$ so that:
\begin{align*}
 T+U&\longmapsto f+g,\\
 TU&\longmapsto fg,\\
 T'&\longmapsto f',\\
 T<U&\longmapsto f<g\text{ eventually}.
\end{align*}
Such an embedding makes formal dominance agree with eventual dominance. For logarithmic-exponential transseries, sophisticated realization results place copies of transseries inside Hardy-field-like structures.

The reverse direction is also important: one asks whether a Hardy germ has a transseries expansion. Many tame germs do, but arbitrary maximal Hardy fields are larger and subtler. Current research aims to understand common extensions and embeddings among Hardy fields, transseries, and surreal numbers.

\section{A three-way dictionary}

\begin{center}
\begin{tabularx}{\textwidth}{>{\bfseries}p{0.20\textwidth}XX}
\toprule
Idea & Formal transseries & Hardy-field germ\\
\midrule
Large variable & symbol $x>\R$ & identity germ $x\mapsto x$\\
Infinitesimal & $x^{-1}\precdom1$ & function tending to $0$\\
Order & sign of leading coefficient & eventual sign\\
Dominance & leading monomial comparison & bounded quotient\\
Derivative & formal strong derivation & derivative of representatives\\
Constant & coefficient in $\R$ & constant germ\\
Composition & formal substitution under support rules & ordinary function composition where defined\\
Integration & formal antiderivative in a closure & germ of an actual antiderivative\\
\bottomrule
\end{tabularx}
\end{center}

\section{Exercises}

\begin{exercise}
Prove that a nonzero germ in a field of germs must be eventually nonzero.
\end{exercise}

\begin{exercise}
Order the germs $\log\log x$, $(\log x)^{10}$, $x^{1/100}$, and $e^{\sqrt{\log x}}$ by strict dominance.
\end{exercise}

\begin{exercise}
Compute the logarithmic derivatives of $x^{\log x}$ and $\exp(x/\log x)$.
\end{exercise}

\begin{exercise}
Give a function tending to $+\infty$ whose derivative changes sign infinitely often, and explain why this behavior conflicts with the $H$-field growth axiom.
\end{exercise}

\begin{summarybox}
Hardy fields are differential fields of nonoscillatory germs at $+\infty$. Their eventual order and dominance mirror leading-term order in transseries. $H$-fields abstract this common structure; Liouville closure, newtonianity, and related axioms describe their ability to solve asymptotic differential equations.
\end{summarybox}


\chapter{Surreal numbers: a numerical home for transseries}\label{ch:surreals}

Transseries treat $x$ as a formal infinitely large quantity. Surreal numbers provide actual canonical numbers that are infinite or infinitesimal. The first infinite ordinal $\omega$ can play the role of $x$, and many transseries can be evaluated at $\omega$ or at other positive infinite surreal numbers.

\section{Numbers born from cuts}

A surreal number is constructed recursively as a cut
\begin{equation}\label{eq:surreal-cut}
 x=\{L\mid R\},
\end{equation}
where every element of the set $L$ is less than every element of the set $R$. The value is the \emph{simplest} number lying strictly between the two sides.

The first examples are
\begin{align*}
 0&=\{\varnothing\mid\varnothing\},\\
 1&=\{0\mid\varnothing\},\\
 -1&=\{\varnothing\mid0\},\\
 \frac12&=\{0\mid1\},\\
 2&=\{1\mid\varnothing\}.
\end{align*}
At each stage, or ``birthday,'' new cuts using earlier numbers create new numbers. Continuing through all ordinal birthdays produces a proper class $\mathbf{No}$ rather than a set.

\section{The first infinite and infinitesimal surreals}

After all finite natural numbers have appeared, the cut
\begin{equation}\label{eq:surreal-omega}
 \omega=\{0,1,2,3,\dots\mid\varnothing\}
\end{equation}
is greater than every natural number. Arithmetic in $\mathbf{No}$ gives
\[
 \omega+1,\qquad \sqrt\omega,\qquad
 \omega^\pi,\qquad e^\omega,
\]
and their reciprocals. In particular,
\[
 \varepsilon=\frac1\omega
\]
is positive but smaller than every positive real number.

Unlike a nonstandard real obtained from an ultrafilter, $\omega$ is canonical inside the surreal construction. No auxiliary choice is used to single it out.

\section{Conway normal form}

Every surreal number has a unique expansion
\begin{equation}\label{eq:conway-normal-form}
 x=\sum_{i<\alpha}r_i\omega^{a_i},
\end{equation}
where:
\begin{itemize}
 \item $\alpha$ is an ordinal;
 \item $r_i\in\R\setminus\{0\}$;
 \item the surreal exponents satisfy $a_0>a_1>a_2>\cdots$.
\end{itemize}
This is a Hahn-series normal form with monomials $\omega^a$. The support is reverse well ordered, exactly as in \cref{ch:hahn}.

For example,
\[
 3\omega^2-\pi\omega^{1/2}+7+\omega^{-1}
\]
is a surreal number whose leading monomial is $\omega^2$. Infinite normal forms can be far longer:
\[
 \omega+1+\omega^{-1}+\omega^{-2}+\cdots
\]
is also meaningful.

\begin{intuitionbox}
A surreal number is simultaneously a number and a generalized series. Conway normal form makes the bridge to transseries visible: replace the formal monomial $x^a$ by the numerical monomial $\omega^a$.
\end{intuitionbox}

\section{Summable families}

Not every arbitrary family of surreal numbers has a sum. A family of normal forms is summable when:
\begin{enumerate}
 \item the union of all monomial supports is reverse well ordered;
 \item each fixed monomial occurs in only finitely many members of the family.
\end{enumerate}
Then coefficients can be added monomial by monomial. This is the same finiteness principle that made Hahn multiplication safe.

Thus the surreal symbol
\[
 \sum_{n\ge0}\omega^{-n}
\]
is legitimate, while an attempted sum placing infinitely many nonzero contributions on the same monomial is not a formal surreal sum unless those real coefficients converge through some separate operation.

\section{Exponentiation and logarithms}

The surreal field admits a canonical exponential function extending the real exponential and satisfying the usual laws:
\[
 \exp(x+y)=\exp(x)\exp(y),
 \qquad
 \exp(0)=1.
\]
It is an order-preserving isomorphism from the additive group of $\mathbf{No}$ onto its positive multiplicative group. Hence every positive surreal has a logarithm.

For an infinitesimal $\varepsilon$, familiar series work:
\begin{align*}
 e^\varepsilon
 &=1+\varepsilon+\frac{\varepsilon^2}{2!}+\cdots,\\
 \log(1+\varepsilon)
 &=\varepsilon-\frac{\varepsilon^2}{2}
 +\frac{\varepsilon^3}{3}-\cdots.
\end{align*}
These are summable Hahn families. For infinite arguments, the full surreal exponential is defined recursively and is compatible with normal forms in a subtler way.

\section{Why \texorpdfstring{$e^\omega$}{exp(omega)} is not merely \texorpdfstring{$\omega^a$}{omega^a}}

Conway monomials have the form $\omega^a$, but exponentiation produces quantities such as $e^\omega$ that are much larger than every fixed power $\omega^a$ with real $a$. They are still surreal monomials, because the exponent $a$ in $\omega^a$ is itself allowed to be a surreal number. There exists a surreal $b$ with
\[
 e^\omega=\omega^b.
\]
The relation between additive exponents and Conway exponents is encoded by an omega-map and related functions. This flexibility lets the surreal monomial group contain towers comparable to transmonomials.

\section{Derivations on the surreals}

A derivative on all surreal numbers cannot mean an ordinary limit of difference quotients: the surreals are not a conventional real manifold. Instead one asks for a derivation
\[
 \partial:\mathbf{No}\to\mathbf{No}
\]
satisfying algebraic and asymptotic requirements:
\begin{align*}
 \partial(x+y)&=\partial x+\partial y,\\
 \partial(xy)&=x\partial y+y\partial x,\\
 \partial(e^x)&=e^x\partial x,\\
 \partial\omega&=1,
\end{align*}
plus compatibility with infinite summation and positivity of derivatives of positive infinite elements.

Berarducci and Mantova constructed a distinguished ``simplest'' derivation with these properties, making $\mathbf{No}$ a Liouville-closed $H$-field \cite{berarducci-mantova-derivation}. Under it,
\[
 \partial(\omega^r)=r\omega^{r-1}
 \qquad(r\in\R),
\]
just as expected when $\omega$ plays the role of an infinite variable.

\section{Embedding ordinary transseries}

There is a natural differential-field embedding
\begin{equation}\label{eq:transseries-surreal-embedding}
 \iota:\T\hookrightarrow\mathbf{No}
\end{equation}
that sends the formal variable $x$ to $\omega$, respects order and exponentiation, and is compatible with the Berarducci--Mantova derivation. At the simplest level,
\begin{align*}
 x^r&\longmapsto\omega^r,\\
 \log x&\longmapsto\log\omega,\\
 e^x&\longmapsto e^\omega,\\
 \sum c_\m\m&\longmapsto\sum c_\m\iota(\m).
\end{align*}
The last line is valid because admissible transseries supports map to summable surreal families.

A major theorem strengthens this: the natural embedding of $\T$ into $\mathbf{No}$ is elementary in the language of ordered differential fields with the relevant structure. In logical terms, every first-order statement about embedded transseries has the same truth value in the ambient surreal field \cite{adh-universal-surreal}.

\section{Omega-series}

Inside $\mathbf{No}$, define the field of \emph{omega-series} as the smallest subfield containing $\R$ and $\omega$ and closed under:
\begin{enumerate}
 \item exponential and logarithm;
 \item sums of all summable families.
\end{enumerate}
This class contains copies of LE-series and EL-series but is substantially larger. It is a proper class, not merely a set-sized field.

Berarducci and Mantova showed that omega-series admit composition with positive infinite surreal arguments and are compatible with differentiation. Each omega-series $f$ therefore acts as a surreal function
\[
 x\longmapsto f\circ x,
 \qquad x>\R,
\]
obtained by substituting $x$ for $\omega$ in a controlled way \cite{berarducci-mantova-germs}.

\section{Evaluation is not naive textual substitution}

Consider a nested expression with infinite sums and exponentials. Replacing every printed $\omega$ by a surreal $x$ is only a mnemonic. One must prove that:
\begin{itemize}
 \item each transformed family remains summable;
 \item different constructions of the same omega-series give the same value;
 \item composition is associative;
 \item differentiation obeys a chain rule;
 \item exponentials and logarithms remain compatible.
\end{itemize}
The formal tree structure of a transseries is used to establish these facts. This is analogous to proving that a formal power series can be composed: the symbolic instruction is simple, but the support theorem is the real work.

\section{Why not every surreal has a satisfactory composition}

The whole surreal field is vastly larger than the omega-series. It can be viewed as an enormous Hahn field of transseries-like normal forms, but there is no known composition operation on all of $\mathbf{No}$ that simultaneously extends the desired substitutions and is compatible with the standard derivation. Indeed, structural obstructions show that the omega-series boundary is meaningful in this context \cite{berarducci-mantova-germs}.

Thus:
\[
 \boxed{\text{every omega-series is surreal}}
 \qquad\text{but}\qquad
 \boxed{\text{not every surreal is an omega-series}.}
\]
Being a generalized number is weaker than being a composable transseries function.

\section{Surreal evaluation versus real asymptotic realization}

If a transseries $T$ maps to a surreal number $T(\omega)$, has it been summed as a real function? Not necessarily.

Surreal evaluation answers an algebraic question: what number results when the infinite variable is interpreted as $\omega$? Borel summation answers an analytic question: what complex or real function has a given divergent asymptotic expansion in a sector? The two constructions preserve different structures and need not coincide term for term.

For a convergent power series in an infinitesimal, both interpretations often agree with familiar algebra. For factorially divergent resurgent sectors, Borel singularities and Stokes choices carry analytic data not supplied merely by a Hahn normal form.

\section{A useful conceptual triangle}

\begin{center}
\begin{tikzpicture}[node distance=4.1cm,>=Latex]
 \node[draw,rounded corners,fill=PaleBlue,align=center] (T) {formal\\transseries};
 \node[draw,rounded corners,fill=PaleTeal,align=center,right=of T] (H) {Hardy-field\\germs};
 \node[draw,rounded corners,fill=PaleGold,align=center,below=2.3cm of $(T)!0.5!(H)$] (N) {surreal\\numbers};
 \draw[->,thick,MediumBlue] (T)--node[above]{realization}(H);
 \draw[->,thick,DeepTeal] (T)--node[left,align=right]{evaluation at\\$\omega$}(N);
 \draw[->,thick,DeepGold] (H)--node[right,align=left]{asymptotic\\embedding}(N);
\end{tikzpicture}
\end{center}

The arrows require hypotheses and choices; the diagram is a research program, not a claim that all three classes are identical.

\section{Exercises}

\begin{exercise}
Using cuts, verify that $\{0\mid1\}$ lies between $0$ and $1$ and is simpler than every other previously available number in that interval.
\end{exercise}

\begin{exercise}
Write the leading Conway normal forms of $\omega+\omega^{-1}$, $(\omega+1)^2$, and $(1+\omega^{-1})^{-1}$.
\end{exercise}

\begin{exercise}
Explain why the family $\{\omega^{-n}:n\in\N\}$ is summable, while a family containing the term $\omega^{-1}$ with nonzero coefficient infinitely many times fails the finite-contribution condition.
\end{exercise}

\begin{exercise}
Distinguish in one paragraph the statements ``$T$ can be evaluated at $\omega$'' and ``$T$ is Borel summable along the positive real direction.''
\end{exercise}

\begin{summarybox}
Surreal numbers form a canonical ordered field containing infinite and infinitesimal numbers. Conway normal form is Hahn-like, and the surreal exponential and Berarducci--Mantova derivation make room for transseries. Omega-series admit controlled infinite sums, composition, and differentiation, while all surreals are too large to carry the same function structure.
\end{summarybox}

\chapter{Beyond the standard LE field}\label{ch:beyond-LE}

There is no single maximal object universally called ``the transseries.'' Different constructions close different starting fields under different kinds of infinite sums, exponentials, logarithms, composition, or transfinite iteration. The standard logarithmic-exponential field $\T$ is extraordinarily useful, but it is one carefully chosen member of a larger ecosystem.

\section{Finite construction depth is a real restriction}

A standard LE-transseries is built by finitely many rounds of operations, even though each round may introduce an infinite Hahn sum. Thus every individual element has finite exponential height and finite logarithmic depth.

Examples such as
\[
 e^{e^{e^x}}
 \quad\text{or}\quad
 \log\log\log x
\]
have large but finite depth and cause no problem. The challenge begins with objects requiring unbounded depth in a single support, for example
\begin{equation}\label{eq:unbounded-depth-example}
 e^{\sqrt{x}+\sqrt{\log x}+\sqrt{\log\log x}+\cdots}.
\end{equation}
No finite construction stage contains all iterated logarithms appearing in the exponent. Larger transseries fields admit such supports while preserving summability, differentiation, and composition.

\section{LE-series and EL-series}

The names reverse the order of two closure ideas.

\begin{description}
 \item[Logarithmic-exponential (LE) series] arise from the construction developed by van den Dries, Macintyre, and Marker, with roots in Dahn--G\"oring. One adjoins exponentials and logarithms through staged generalized-series extensions.
 \item[Exponential-logarithmic (EL) series] arise from prelogarithmic Hahn fields followed by exponential closure, as developed in work of Kuhlmann and collaborators.
\end{description}

The resulting fields satisfy many of the same first-order properties as non-Archimedean models of the real exponential field, but the constructions are intrinsically different. Kuhlmann and Tressl proved that the LE field embeds as an ordered exponential field into any of the relevant EL fields, whereas an EL field does not embed into the LE field in the corresponding way \cite{kuhlmann-tressl}.

This asymmetry is a useful warning: two fields may contain all familiar finite expressions and satisfy the same broad logical theory while differing on which infinite generalized sums they admit.

\section{Grid-based versus well-based transseries}

A grid-based transseries has support inside a finitely generated ratio grid
\[
 \m_0\boldsymbol\mu^{\mathbf k},
 \qquad \mathbf k\in\N^r.
\]
A well-based transseries needs only reverse well-order, with no finite ratio basis.

Grid-based fields are smaller but algorithmically friendly:
\begin{itemize}
 \item monomials can be represented by integer exponent vectors;
 \item truncation regions are finite-dimensional;
 \item multiplication is sparse convolution on $\Z^r$;
 \item fixed-point proofs can use contraction in grid order.
\end{itemize}

Well-based fields capture supports such as
\[
 \{e^{-x},e^{-\sqrt2x},e^{-\pi x},\dots\}
\]
or increasingly complex iterated-log scales that do not lie in one finite grid. They are closer to a maximal formal universe but harder to compute with.

\section{Schmeling's axiomatic transseries fields}

Schmeling's thesis develops broad axioms for fields of transseries, including strong derivations and exponential-logarithmic structure \cite{schmeling}. Rather than defining one field by a single inductive recipe, the axiomatic approach asks which support, summation, exponential, logarithm, and derivation laws a satisfactory transseries field should satisfy.

Typical desiderata include:
\begin{enumerate}
 \item a Hahn-field presentation with a distinguished monomial group;
 \item closure under summable families;
 \item exponential and logarithm extending the real functions;
 \item a strong derivation preserving admissible sums;
 \item compatibility of derivative with exponential and logarithm;
 \item enough monomials to integrate broad classes of series;
 \item a composition law, when the theory is intended to treat functions.
\end{enumerate}
These properties are not automatically compatible on an arbitrarily large Hahn field.

\section{Superexponentials and iteration beyond integers}

The ordinary exponential tower is defined recursively:
\[
 E_0(x)=x,
 \qquad
 E_{n+1}(x)=e^{E_n(x)}.
\]
A \emph{superexponential} seeks a function $E_t(x)$ defined for noninteger $t$ such that
\[
 E_{s+t}=E_s\circ E_t,
 \qquad
 E_1=e^x.
\]
This is an Abel/iteration problem of the kind introduced in \cref{ch:fractional-iteration}, but now the function grows so quickly that ordinary LE height is shifted by composition.

Its inverse family gives hyperlogarithms. Closing a transseries field under such operators requires monomials indexed not only by finite exponential height but by more general ordinal levels.

\section{Logarithmic hyperseries}

Logarithmic hyperseries extend iterated logarithms into transfinite hierarchies. One introduces symbols behaving like
\[
 \ell_0=x,
 \qquad
 \ell_{n+1}=\log\ell_n,
\]
and then continues the indexing beyond finite $n$ in a coherent formal system. At a limit ordinal, no single finite iterate is available, so the construction must use asymptotic relations and generalized supports rather than an ordinary recursive formula.

Van den Dries, van der Hoeven, and Kaplan developed logarithmic hyperseries as a field with composition and derivation, providing a rigorous home for extremely slowly varying scales \cite{log-hyperseries}. These scales arise naturally when repeatedly solving equations whose inverses introduce more logarithms than any fixed finite depth.

\section{Hyperserial fields}

A hyperserial field is designed to support hyperexponentials, hyperlogarithms, Hahn summation, and composition across transfinite levels. Recent work shows that the surreal numbers carry substantial hyperserial structure and can serve as a vast ambient field for these hierarchies \cite{hyperserial-field,surreal-hyperseries}.

The guiding ladder is
\[
 \text{power series}
 \subset
 \text{LE-transseries}
 \subset
 \text{omega-series}
 \subset
 \text{hyperseries-like structures}
 \subset
 \mathbf{No},
\]
but each inclusion suppresses qualifications about the chosen embedding and operations. In particular, a larger ambient field may contain the elements of a smaller one without extending every composition or derivation in the desired way.

\section{Height, level, and comparability}

Two positive infinite elements $f,g$ are logarithmically comparable if repeated logarithms eventually bring them to comparable magnitude. Informally, one studies whether there is an integer $n$ such that
\[
 \log_n f\asymp\log_n g.
\]
The resulting \emph{levels} distinguish growth like
\[
 x,
 \quad e^x,
 \quad e^{e^x},
 \quad \dots.
\]
LE-transseries can move through finitely many such levels. Hyperexponential structures allow systematic movement through transfinite level indices.

These comparisons matter for composition. Composing with $e^x$ raises levels; composing with $\log x$ lowers them. A field closed under arbitrary repeated composition must have room for every level reached.

\section{Oscillatory and complex transseries}

Ordered real transseries deliberately avoid persistent oscillation, but asymptotic analysis often needs phases such as
\[
 e^{\ii x},
 \qquad
 e^{\ii A/z},
 \qquad
 \cos(e^x).
\]
One can work over complex coefficients and monomial groups that are not totally ordered by magnitude, or separate magnitude from phase:
\[
 e^{-(a+\ii b)/z}
 =e^{-a/z}e^{-\ii b/z}.
\]
The first factor controls exponential size and the second controls oscillation.

Multisummable and resurgent transseries routinely use complex actions. Their sector order is often only partial: two exponentials can have the same magnitude but different phase. This analytic tradition overlaps with, but is not identical to, the ordered Hahn-field tradition.

\section{Transseries with fractional powers and logarithmic branching}

Near singular points, natural ansatzes include
\[
 z^\alpha(\log z)^m
 \sum_{n\ge0}a_nz^{n/q}.
\]
Ramification replaces $z$ by $w^q$, turning fractional powers into integer powers on a covering surface. Logarithms record monodromy. Exponential sectors then take the form
\[
 e^{-A/z^p}z^\alpha(\log z)^m\trans\phi(z^{1/q}).
\]
Such objects fit naturally into resurgent analysis even when they are not presented as elements of one canonical ordered LE field.

\section{Why there is no final largest field}

Suppose a field contains a fast-growing $F$. One can ask for a compositional inverse, a fractional iterate, or a solution to
\[
 G'=F,
 \qquad
 H\circ H=F,
 \qquad
 K(x+1)=F(K(x)).
\]
Adjoining these solutions may create new growth classes, which invite new exponentials, sums, and inverses. Any chosen closure principle suggests a stronger one.

Therefore ``the field of all transseries'' is best understood relative to a signature of desired operations. The standard $\T$ is a canonical and remarkably complete answer for finite logarithmic-exponential asymptotics and differential algebra, not an absolute terminal object for every possible asymptotic operation.

\section{A map of common constructions}

\begin{longtable}{>{\raggedright\arraybackslash}p{0.20\textwidth}
                  >{\raggedright\arraybackslash}p{0.32\textwidth}
                  >{\raggedright\arraybackslash}p{0.36\textwidth}}
\toprule
Construction & Main strength & Typical limitation or question\\
\midrule
\endfirsthead
\toprule
Construction & Main strength & Typical limitation or question\\
\midrule
\endhead
Grid-based transseries & finite-dimensional support bookkeeping; algorithms & excludes many well-based supports\\
Well-based LE-series & broad Hahn sums; exp, log, derivative & every element still has finite construction depth\\
EL-series & large exponential closures of prelogarithmic Hahn fields & not isomorphic to the LE construction\\
Resurgent transseries & analytic continuation, Stokes data, large order & requires controlled Borel singularities\\
Omega-series & surreal evaluation, summation, composition & proper class; smaller than all surreals\\
Logarithmic hyperseries & transfinite slowly varying scales & specialized hierarchy and substantial machinery\\
Hyperserial fields & hyperexponentials and hyperlogarithms & active structural development\\
All surreal numbers & universal numerical ambient class & no compatible global composition of the desired kind\\
\bottomrule
\end{longtable}

\section{Exercises}

\begin{exercise}
Explain why every finite tower $\exp^{\circ n}(x)$ can belong to an LE field even though the family of all heights is unbounded.
\end{exercise}

\begin{exercise}
Give a support involving all iterated logarithms that cannot be contained in any fixed finite logarithmic depth.
\end{exercise}

\begin{exercise}
Compare the algorithmic advantages of a two-generator ratio grid with a merely reverse-well-ordered support.
\end{exercise}

\begin{exercise}
Why does containment of two fields as sets not guarantee that composition on the smaller field extends to the larger one?
\end{exercise}

\begin{summarybox}
The standard LE field is one disciplined closure of generalized series under finite rounds of exp, log, and summation. EL fields, omega-series, logarithmic hyperseries, and hyperserial fields admit different larger supports or iteration levels. ``Transseries'' names a family of compatible constructions rather than one uncontested maximal object.
\end{summarybox}


\chapter{Model theory and asymptotic differential algebra}\label{ch:model-theory}

Model theory studies mathematical structures through statements written in a fixed formal language. For transseries, it asks a surprisingly practical question:

\begin{quote}
Which systems of algebraic differential equations and inequalities must have solutions merely because the ambient field has the right asymptotic structure?
\end{quote}

The answer turns the field $\T$ from a collection of formal expressions into a universal domain for nonoscillatory asymptotic differential algebra.

\section{Languages, structures, and sentences}

A first-order language lists allowed symbols. For ordered differential fields, one may use
\[
 0,1,+,-,\cdot,<,\partial,
\]
and perhaps a dominance relation $\preceq$. A structure interprets these symbols: in $\T$, $\partial$ is formal differentiation and $<$ is leading-term order.

A first-order formula can quantify over field elements:
\[
 \exists y\ \bigl(y'>0\ \wedge\ y^2=x\bigr).
\]
It cannot directly quantify over arbitrary subsets, over all sequences, or over all supports unless those are encoded as field elements.

The \emph{theory} of a structure is the collection of all first-order sentences true in it.

\section{Elementary equivalence}

Two structures are elementarily equivalent if they satisfy exactly the same first-order sentences. They need not be isomorphic or have the same cardinality.

For example, different real closed fields satisfy the same first-order theory in the language of ordered fields. One may be Archimedean and another may contain infinitesimals, yet every sentence involving only field operations, order, and quantification over elements has the same truth value.

Likewise, different $H$-closed fields can share the first-order theory of transseries even if their elements have very different concrete representations.

\section{Elementary embeddings}

An embedding $K\hookrightarrow L$ is elementary if every first-order formula with parameters from $K$ has the same truth value before and after embedding. It preserves not only equations and inequalities but every finite logical combination and quantifier pattern.

The natural embedding
\[
 \T\hookrightarrow\mathbf{No}
\]
with the Berarducci--Mantova derivation is elementary in the appropriate $H$-field language \cite{adh-universal-surreal}. Thus $\mathbf{No}$ contains $\T$ without introducing a new first-order solution to a finite differential-algebraic problem over $\T$ that was absent from $\T$ itself.

\section{Existential closure}

An existential statement has the form
\[
 \exists y_1\cdots\exists y_n\;\Phi(\mathbf y,\mathbf a),
\]
where $\Phi$ has no further quantifiers. It can express a finite system of differential equations and inequalities.

A field $K$ is existentially closed in a class if every such system over $K$ that has a solution in an extension from the class already has a solution in $K$.

For transseries, this means roughly:
\begin{quote}
If a finite nonoscillatory asymptotic differential problem is compatible with the $H$-field axioms, then $\T$ already contains a solution.
\end{quote}
This is a structural existence theorem, not an algorithm that immediately prints the solution's coefficients.

\section{Model completeness}

A theory is model complete if every embedding between its models is elementary. Once two $H$-closed fields sit inside one another, no new first-order facts about the smaller field appear in the larger.

The theory of $H$-closed fields is model complete. In one expository axiomatization, the theory of $\T$ is specified by:
\begin{enumerate}
 \item being an $H$-field with small derivation;
 \item being Liouville closed;
 \item satisfying the differential intermediate value property.
\end{enumerate}
Equivalent formulations use the more structural conditions of newtonianity and $\omega$-freeness \cite{adh-book,adh-numbers}.

\section{Quantifier elimination}

Quantifier elimination means that every first-order formula is equivalent, in the theory, to one without quantifiers in an enriched language. For real closed fields, quantifier elimination underlies exact algorithms for polynomial inequalities.

For $H$-closed fields, quantifier elimination requires additional predicates and functions encoding certain asymptotic cuts. The result says that complicated quantified questions ultimately reduce to explicit algebraic, differential, order, and valuation conditions in that richer language.

Consequences include:
\begin{itemize}
 \item completeness: all models satisfy the same sentences;
 \item decidability: there is an algorithm that eventually answers every sentence in the formal language;
 \item tameness of definable sets, including strong restrictions on eventual one-variable behavior.
\end{itemize}
``Decidable'' here concerns exact symbolic sentences, not numerical convergence or the practical complexity of computing a transseries solution.

\section{A simple logical example}

Consider the sentence
\begin{equation}\label{eq:logical-antiderivative}
 \forall f\ \exists g\;(g'=f).
\end{equation}
This is one part of Liouville closure. The sentence
\begin{equation}\label{eq:logical-exp-integral}
 \forall f\ \exists h\;(h\ne0\wedge h'=fh)
\end{equation}
asks for exponential integrals.

A differential intermediate-value statement has more structure:
\[
 \forall f\forall g\,
 \bigl(f<g\wedge P(f)<0<P(g)
 \Rightarrow\exists y(f<y<g\wedge P(y)=0)\bigr).
\]
Once $P$ has fixed finite order and degree, this is first order.

\section{Differential algebraicity}

An element $y$ is differentially algebraic over $K$ if it satisfies a nonzero differential polynomial equation
\begin{equation}\label{eq:differential-algebraic}
 P(y,y',\dots,y^{(r)})=0
\end{equation}
with coefficients in $K$. Otherwise it is differentially transcendental.

Solutions of algebraic ODEs are differentially algebraic by definition. Many familiar special functions are differentially algebraic over rational-function fields; others, such as the gamma function, satisfy difference equations but are differentially transcendental over $\C(x)$.

The transseries field contains both differential-algebraic and differential-transcendental elements. Its model theory describes finite differential-polynomial constraints without reducing every element to a closed form.

\section{Newton degree}

For an ordinary polynomial, the lowest significant terms after a valuation rescaling determine whether Newton's method can lift a root. For a differential polynomial, one rescales both the unknown and the derivation to expose a stable dominant differential polynomial called the Newton polynomial.

The \emph{Newton degree} is its degree in the differential indeterminate. If the Newton degree is $1$, newtonianity guarantees a root in the valuation ring. This abstractly justifies many transseries correction procedures:
\[
 \begin{gathered}
  \text{dominant balance}
  \longrightarrow
  \text{linearized leading equation}\\[-0.2ex]
  \Downarrow\\[-0.2ex]
  \text{smaller correction}
  \longrightarrow
  \text{exact formal root}.
 \end{gathered}
\]
The concrete fixed-point algorithms of \cref{ch:dominant-balance} are computational shadows of this valuation-theoretic principle.

\section{Why oscillation is the boundary}

The order in an $H$-field is total and compatible with eventual growth. Oscillatory solutions can defeat this framework. For example,
\[
 y''+y=0
\]
has nonzero real solutions that change sign forever. No nonzero solution has an eventual sign, so it cannot be an element of a real Hardy field.

One may still study the equation with complex transseries, exponential phases $e^{\pm\ii x}$, and resurgent methods. What fails is not asymptotic analysis itself, but the ordered $H$-field interpretation.

This explains the common phrase that $\T$ is a universal domain for \emph{nonoscillatory} asymptotic differential algebra.

\section{Definability versus explicit construction}

Model theory can prove
\[
 \exists y\;P(y)=0
\]
without giving a support-by-support recipe for $y$. Conversely, a constructive transseries algorithm may compute thousands of coefficients without clarifying whether the solution is unique under every definable condition.

The two viewpoints complement one another:
\begin{description}
 \item[constructive asymptotics] finds leading monomials, recurrences, truncations, and numerical values;
 \item[model theory] proves closure, transfer, uniqueness patterns, and impossibility of certain definable pathologies.
\end{description}

\section{Decidability does not mean easy computation}

A decidable theory has an algorithm that terminates on every sentence. It may be far too slow for practical use. Moreover, questions about a specific infinitely described coefficient sequence may not be first-order sentences in the field language.

For example, the following are different tasks:
\begin{enumerate}
 \item decide whether every solution of a fixed differential equation eventually exceeds $\log x$;
 \item compute the first million coefficients of one formal solution;
 \item locate all Borel singularities of that solution;
 \item prove that a numerical Stokes constant is nonzero.
\end{enumerate}
Only the first is immediately in the style of first-order $H$-field logic. The others require algorithms or analytic information outside bare decidability.

\section{Current bridges among germs, transseries, and surreals}

The common $H$-field language makes it possible to compare:
\begin{itemize}
 \item formal LE-transseries;
 \item Hardy fields of actual germs;
 \item the surreal field with its distinguished derivation.
\end{itemize}
Major embedding and elementary-equivalence results now connect these worlds. Questions about maximal Hardy fields, hyperserial extensions, and canonical analytic realizations remain active. The goal is not to collapse the three worlds into one, but to determine exactly which algebraic, differential, order, and composition properties transfer among them.

\section{Exercises}

\begin{exercise}
Write a first-order sentence expressing that every positive element has a square root.
\end{exercise}

\begin{exercise}
Explain the difference between an embedding of ordered differential fields and an elementary embedding.
\end{exercise}

\begin{exercise}
Why can the property ``every sequence has a least upper bound'' not be expressed directly by one first-order sentence in the language of ordered fields?
\end{exercise}

\begin{exercise}
Classify $e^x$, a solution of $y'=y$, and a generic formal series with algebraically independent coefficients as differentially algebraic or potentially differentially transcendental over $\R(x)$.
\end{exercise}

\begin{summarybox}
Model theory describes $\T$ through finite statements about field operations, order, dominance, and differentiation. $H$-closed fields form a model-complete theory, and an enriched language admits quantifier elimination and decidability. These results make transseries a universal domain for broad nonoscillatory differential-algebraic problems, while remaining distinct from coefficient computation and Borel analysis.
\end{summarybox}

\chapter{How to represent and compute with transseries}\label{ch:computation}

A computer algebra system cannot store a completed infinite transseries. It stores finite descriptions, lazy generators, and truncation rules. The mathematical support conditions tell us which operations are legitimate; the software design must make those conditions explicit enough to enforce.

This chapter describes an implementation architecture rather than a particular programming language. The central rule is:

\begin{quote}
Never represent a transseries as an unstructured syntax tree plus a hope that simplification will sort it out. Separate monomials, coefficients, ordered support, construction level, and truncation policy.
\end{quote}

\section{Three layers of representation}

A practical system benefits from three distinct layers.

\begin{enumerate}
 \item \textbf{Expression layer.} User syntax such as
 \[
  e^{x+1/x}\left(1+\frac2x\right)+\log x.
 \]
 \item \textbf{Canonical transseries layer.} A sparse ordered map
 \[
  \m\longmapsto c_\m
 \]
 with normalized transmonomials and admissible support.
 \item \textbf{Truncated computation layer.} A finite projection determined by a cutoff, ratio grid, requested precision, or observable.
\end{enumerate}

The expression layer preserves convenient structure. The canonical layer decides equality and leading terms. The truncated layer makes algorithms terminate.

\section{A monomial data type}

At one finite logarithmic-exponential level, a monomial may be represented conceptually as
\begin{equation}\label{eq:monomial-record}
 \m=x^a(\log x)^b\exp(L),
\end{equation}
where $a,b\in\R$ and $L$ is a purely large transseries of lower or controlled height. More generally, iterated logarithms can be stored as a finite exponent vector:
\[
 x^{a_0}(\log x)^{a_1}(\log_2x)^{a_2}\cdots
 \exp(L).
\]

A possible immutable record is:

\begin{lstlisting}[style=pseudo]
Monomial:
    logPowers : finite map depth -> exact real exponent
    exponent  : PurelyLargeSeries
    height    : nonnegative integer
    cachedHash, cachedLeadingData
\end{lstlisting}

Multiplication adds logarithmic-power vectors and exponent series. Inversion negates them. Equality is structural after normalization.

\section{Canonical normalization}

Expressions with different syntax can denote the same monomial:
\[
 e^{2\log x}=x^2,
 \qquad
 x^ae^L\,x^be^M=x^{a+b}e^{L+M}.
\]
A normalizer should:
\begin{enumerate}
 \item flatten products and exponentials where formal identities permit;
 \item collect real multiples of iterated logarithms into power coordinates;
 \item normalize $L$ as a purely large series with no constant or small part;
 \item remove zero exponents;
 \item canonicalize exact real coefficients and branch conventions.
\end{enumerate}

Over the real ordered field, $x>0$ and principal real logarithms remove many branch ambiguities. Over complex transseries, the normal form must record branches or work on a chosen logarithmic covering.

\section{Comparing monomials}

To decide whether
\[
 \m=e^Lx^a\ell_1^{a_1}\cdots
 \quad\text{is larger than}\quad
 \n=e^Mx^b\ell_1^{b_1}\cdots,
\]
compare their logarithms:
\begin{equation}\label{eq:monomial-log-compare}
 \log\frac\m\n
 =(L-M)+(a-b)\log x+(a_1-b_1)\log_2x+\cdots.
\end{equation}
The sign of the leading term of \eqref{eq:monomial-log-compare} determines the order.

Thus monomial comparison is recursive in exponential height. Memoizing leading data is essential: sorting a support can otherwise recompute the same lower-level differences many times.

\begin{lstlisting}[style=pseudo]
CompareMonomials(m, n):
    if m is n: return Equal
    D := Normalize(LogMonomial(m) - LogMonomial(n))
    return SignOfLeadingTerm(D)
\end{lstlisting}

A production implementation avoids constructing the full difference when a lexicographic leading signature already distinguishes the monomials.

\section{Series as sparse ordered maps}

A finite truncation can use a balanced search tree keyed by monomial order:

\begin{lstlisting}[style=pseudo]
Series:
    terms : OrderedMap<Monomial, Coefficient>
    ambientLevel : LevelDescriptor
    supportProof : SupportCertificate
\end{lstlisting}

The map should enumerate from largest to smallest monomial. Then
\[
 \lm(T)=\text{first key},
 \qquad
 \lc(T)=\text{its value}.
\]
Zero coefficients are removed immediately.

For lazy infinite series, replace the complete map by a stream of decreasing terms plus a certificate that the stream is well based. Not every generator that emits decreasing terms is safe for multiplication; one also needs local finiteness guarantees.

\section{Ratio-grid representation}

If the support lies in
\[
 \m_0\mu_1^{k_1}\cdots\mu_r^{k_r},
 \qquad\mathbf k\in\N^r,
\]
store each term by its exponent vector $\mathbf k$. Multiplication becomes addition of integer vectors.

\begin{lstlisting}[style=pseudo]
GridSeries:
    base       : Monomial
    generators : array<SmallMonomial>
    coeffs     : SparseMap<IntegerVector, Coefficient>
    cone       : admissible affine submonoid
\end{lstlisting}

The monomial order is induced by the generators, but it need not be ordinary lexicographic order on $\N^r$. One can precompute a comparison oracle from leading logarithms of the $\mu_j$.

Grid representation makes truncation geometric. A cutoff corresponds to a finite set of lattice points in a cone.

\section{Truncation is part of the semantics}

A request such as ``expand to order $x^{-10}$'' is insufficient when exponentials and logarithms coexist. Possible cutoffs include:
\begin{description}
 \item[monomial cutoff:] keep every term $\m\succeqdom\tau$;
 \item[grid box:] keep $0\le k_j\le N_j$;
 \item[sector cutoff:] keep instanton degree $|\mathbf k|\le K$ and perturbative order $n\le N$;
 \item[term budget:] keep the first $M$ terms in dominance order;
 \item[adaptive observable:] keep enough terms to determine a requested leading residual or numerical tolerance.
\end{description}

A mathematically reliable API passes a \texttt{TruncationContext} through every operation. Silent global defaults invite errors, because an intermediate product may require terms below the final cutoff to produce retained terms after cancellation or division.

\section{Addition}

Addition is a merge of two decreasing support streams.

\begin{lstlisting}[style=pseudo]
Add(A, B, cutoff):
    i := A.leadingIterator()
    j := B.leadingIterator()
    while i or j has term above cutoff:
        compare i.monomial and j.monomial
        emit larger term
        if equal: emit summed coefficient if nonzero
\end{lstlisting}

The cost is linear in the number of retained terms, apart from monomial comparisons.

Cancellation can expose a much smaller leading term. Therefore a lazy iterator must be able to request further terms when equal leading coefficients cancel.

\section{Multiplication}

For finite truncations, multiplication is sparse convolution:
\begin{equation}\label{eq:sparse-product}
 [\m](AB)=\sum_{\n\mathfrak p=\m}[\n]A\,[\mathfrak p]B.
\end{equation}
The Neumann support theorem guarantees finite coefficient sums for well-based supports. The algorithm still needs a way to enumerate only products above the cutoff.

A naive double loop over $N$ and $M$ retained terms costs $O(NM)$. A priority-queue algorithm can enumerate products in descending monomial order:

\begin{lstlisting}[style=pseudo]
MultiplyStreams(A, B, cutoff):
    insert candidate (0, 0) with monomial A[0].m * B[0].m
    while queue maximum is above cutoff:
        pop all candidates with the same product monomial
        sum their coefficient products
        emit if nonzero
        insert unseen neighboring index pairs
\end{lstlisting}

For grid series, hash convolution on integer vectors is often simpler and faster. Fast polynomial multiplication can be used when a finite rectangular grid is dense.

\section{Geometric inversion}

Given nonzero
\[
 T=a\m(1+S),
 \qquad S\precdom1,
\]
compute
\begin{equation}\label{eq:computational-inverse}
 T^{-1}=a^{-1}\m^{-1}\sum_{n\ge0}(-S)^n.
\end{equation}

\begin{lstlisting}[style=pseudo]
Inverse(T, cutoff):
    (a, m) := LeadingTerm(T)
    S := T/(a*m) - 1
    result := 1
    power := 1
    n := 1
    while power can still contribute above cutoff:
        power := Truncate(power * S, cutoffRelativeTo(m))
        result += (-1)^n * power
        n += 1
    return Truncate(result/(a*m), cutoff)
\end{lstlisting}

The termination proof is asymptotic: successive powers of $S$ become smaller. A coarse cutoff can often bound in advance how many powers are needed.

\section{Exponentiation}

Use the decomposition
\[
 T=L+c+S
\]
with $L$ purely large, $c$ constant, and $S$ small. Then
\begin{equation}\label{eq:computational-exp}
 e^T=e^Le^c\sum_{n\ge0}\frac{S^n}{n!}.
\end{equation}
The factor $e^L$ is a new monomial at the next height. Only the small part is expanded.

\begin{lstlisting}[style=pseudo]
Exp(T, cutoff):
    (L, c, S) := SplitLargeConstantSmall(T)
    m := InternMonomial(ExpOfPurelyLarge(L))
    P := TruncatedPowerSeriesExp(S, cutoff / m)
    return exp(c) * m * P
\end{lstlisting}

Interning canonical monomials ensures that identical $e^L$ nodes share storage and compare quickly.

\section{Logarithms and real powers}

For positive
\[
 T=a\m(1+S),
 \qquad a>0,
\]
use
\begin{equation}\label{eq:computational-log}
 \log T=\log a+\log\m+
 \sum_{n\ge1}\frac{(-1)^{n+1}}nS^n.
\end{equation}
Then
\[
 T^r=\exp(r\log T).
\]

The positivity test is formal: $T>0$ exactly when its leading coefficient is positive in the real ordered field. For rational powers of a negative series, a separate real-domain policy is needed.

\section{Differentiation}

Differentiate termwise:
\[
 (c\m)'=c\m\,(\log\m)'.
\]
A single monomial may differentiate into a series because $(\log\m)'$ can have several terms.

\begin{lstlisting}[style=pseudo]
Differentiate(T, cutoff):
    result := 0
    for (m, c) in terms of T that may contribute:
        result += c * m * Differentiate(LogMonomial(m))
    return NormalizeAndTruncate(result, cutoff)
\end{lstlisting}

A strong derivation must preserve admissible infinite sums. For a lazy implementation, this requires a support theorem showing that derivative contributions remain locally finite above each cutoff.

\section{Formal integration}

Integration is not simply monomial-by-monomial division. For
\[
 \int e^L\dd x,
\]
start with
\[
 \frac{e^L}{L'}
\]
and correct the residual. The fixed-point algorithm of \cref{ch:integration} can be implemented as:

\begin{lstlisting}[style=pseudo]
Integrate(F, cutoff):
    G := LeadingAntiderivativeGuess(F)
    repeat:
        R := Truncate(F - Differentiate(G), residualCutoff)
        if R has no term above cutoff: return G
        G := Truncate(G + LeadingAntiderivativeGuess(R), cutoff)
\end{lstlisting}

Each correction removes the current leading residual. Termination at finite cutoff follows if residual leading monomials strictly decrease.

Logarithmic resonance must be detected. Integrating $x^{-1}$ creates $\log x$ rather than $x^0/0$.

\section{Composition}

Composition is the most delicate routine. For $T\circ S$:
\begin{enumerate}
 \item verify that $S$ is positive infinite in the domain expected by $T$;
 \item recursively substitute in logarithms and exponentials;
 \item preserve summability of every support family;
 \item use Taylor expansion only when the increment is small in the correct relative sense;
 \item propagate a cutoff backward through the substitution.
\end{enumerate}

For a Taylor step,
\[
 F(S+U)=\sum_{n\ge0}\frac{F^{(n)}(S)}{n!}U^n,
\]
the stopping criterion is not merely $n=N$. Stop when no later $U^nF^{(n)}(S)$ can reach the requested monomial cutoff.

\begin{warningbox}
A syntactically small $U$ need not be small enough for Taylor composition. The relevant comparison depends on the derivative scale of $F$. For $F=e^{e^x}$, an increment $U=e^{-x/2}$ tends to zero, yet $e^xU=e^{x/2}$ is large, so the first exponential response is not a small perturbation.
\end{warningbox}

\section{Fixed points by coefficient stabilization}

Many equations can be written
\[
 Y=\Phi(Y)
\]
with $\Phi$ contracting in the dominance topology. A lazy solver iterates
\[
 Y_{n+1}=\Phi(Y_n)
\]
and records which leading coefficients have stabilized.

\begin{lstlisting}[style=pseudo]
SolveFixedPoint(Phi, cutoff):
    Y := InitialApproximation()
    stablePrefix := empty
    loop:
        Z := Truncate(Phi(Y), cutoff)
        stablePrefix := CommonLeadingPrefix(Y, Z)
        if stablePrefix reaches cutoff: return Z
        Y := Z
\end{lstlisting}

Equality of finite truncations in two consecutive iterations is not by itself a proof if omitted lower terms can feed upward through a noncontractive operation. The solver should carry a contraction certificate or use an equation-specific theorem.

\section{Solving algebraic equations}

A transseries Newton method uses
\begin{equation}\label{eq:computer-newton}
 Y_{n+1}=Y_n-\frac{P(Y_n)}{P'(Y_n)}.
\end{equation}
Dominance analysis chooses the first approximation and verifies that $P'(Y_n)$ is invertible. Truncation must preserve enough residual terms to compute the next correction.

For multiple roots, Newton's method may not improve valuation quadratically. A Newton-polygon phase should first detect balances among monomials of $P$ and split branches.

\section{Sector-indexed resurgent data}

A resurgent transseries is naturally a sparse map from sector multi-indices to perturbative series:

\begin{lstlisting}[style=pseudo]
ResurgentTransseries:
    actions      : array<Complex>
    parameters   : symbolic variables
    sectors      : map<MultiIndex, FormalPowerSeries>
    prefactors   : map<MultiIndex, (power, logPolynomial)>
    stokesData   : directed labeled graph
\end{lstlisting}

Products convolve multi-indices. Differential equations generate sector recurrences. The Borel transform acts inside each sector. Alien derivatives move along labeled edges of the Stokes graph.

For numerical work, coefficient arrays are finite and often high precision. Store exact rationals or symbolic parameters as long as possible, then convert to arbitrary-precision complex numbers for Borel--Pad\'e analysis.

\section{Borel--Pad\'e reconstruction}

Given coefficients $a_0,\dots,a_N$:
\begin{enumerate}
 \item form the truncated Borel series
 \[
  g_N(\xi)=\sum_{n=0}^N\frac{a_n}{n!}\xi^n;
 \]
 \item compute a near-diagonal Pad\'e approximant $P(\xi)/Q(\xi)$;
 \item inspect poles and pole strings as approximations to Borel singularities and cuts;
 \item integrate the rational approximant along a deformed ray;
 \item vary $N$, precision, Pad\'e order, and contour to estimate stability.
\end{enumerate}

Spurious pole-zero pairs, called Froissart doublets, can mimic singularities. True singularity estimates should persist as the order and precision change.

For a singular direction, choose upper and lower contours and compute the difference. Comparing it with a separately summed neighboring sector estimates the Stokes constant.

\section{Exact coefficients and coefficient domains}

Coefficient arithmetic may involve:
\begin{itemize}
 \item rational numbers;
 \item algebraic numbers;
 \item symbolic constants such as $\pi$ or Stokes parameters;
 \item formal power series in auxiliary parameters;
 \item arbitrary-precision complex approximations.
\end{itemize}

Do not identify coefficients merely because their floating-point approximations are close. Canonical equality belongs to an exact coefficient domain; numerical equality belongs to an error model.

A tower of coefficient domains can be useful:
\[
 \Q\subset\overline{\Q}\subset
 \Q(\pi,S_1,\dots)\subset\C_{\mathrm{arb}}.
\]
Coercions upward should be explicit.

\section{Support certificates}

A \emph{support certificate} records why a lazy family is admissible. Examples:
\begin{itemize}
 \item finite support;
 \item support inside a certified ratio grid;
 \item decreasing sequence with a proven order type;
 \item image of a well-based support under a monotone monomial map;
 \item finite union, Minkowski sum, or derivative image of certified supports.
\end{itemize}

The certificate need not be a formal proof object in every implementation, but the architecture should preserve enough provenance to reject obviously unsafe operations.

\section{Testing invariants}

Property-based tests are especially valuable. Generate random finite grid truncations and verify:
\begin{align*}
 A(B+C)&=AB+AC,\\
 (AB)'&=A'B+AB',\\
 \exp(A+B)&=\exp(A)\exp(B)
 \quad\text{when the operation is admitted},\\
 \log(AB)&=\log A+\log B
 \quad(A,B>0),\\
 A\cdot A^{-1}&=1\pmod{\text{cutoff}},\\
 (F\circ G)'&=(F'\circ G)G'
 \quad\text{under composition hypotheses}.
\end{align*}

Also test ordering laws and support invariants:
\begin{itemize}
 \item every term stream is strictly decreasing;
 \item no zero coefficient is stored;
 \item multiplication produces only finitely many contributions per retained monomial;
 \item increasing the cutoff refines rather than changes already stable leading coefficients.
\end{itemize}

\section{A minimal end-to-end example}

To expand
\[
 T=\frac{e^{x+1/x}}{x^2(1-3/x)}
\]
through $e^xx^{-6}$:
\begin{enumerate}
 \item split $e^{x+1/x}=e^xe^{1/x}$;
 \item create the monomial $e^xx^{-2}$;
 \item expand
 \[
 e^{1/x}=1+x^{-1}+\tfrac12x^{-2}+\tfrac16x^{-3}+\tfrac1{24}x^{-4}+O(x^{-5});
 \]
 \item expand
 \[
 (1-3/x)^{-1}=1+3x^{-1}+9x^{-2}+27x^{-3}+81x^{-4}+O(x^{-5});
 \]
 \item convolve the two grid coefficient arrays through relative order $x^{-4}$;
 \item multiply by $e^xx^{-2}$.
\end{enumerate}
The result is
\begin{equation}\label{eq:implementation-example}
 T=e^xx^{-2}\left(
 1+4x^{-1}+\frac{25}{2}x^{-2}
 +\frac{113}{3}x^{-3}
 +\frac{2713}{24}x^{-4}
 +O(x^{-5})\right).
\end{equation}
This simple calculation exercises monomial extraction, two small-series expansions, sparse convolution, and an absolute cutoff.

\section{Exercises}

\begin{exercise}
Design a canonical record for monomials of the form $x^a(\log x)^b e^L$ and specify multiplication, inversion, and comparison.
\end{exercise}

\begin{exercise}
Explain why a global ``keep 20 terms'' policy can fail when leading terms cancel during addition.
\end{exercise}

\begin{exercise}
Implement on paper the geometric inversion of $1+x^{-1}+e^{-x}$ through the cutoff $e^{-x}x^{-2}$, carefully ordering all retained terms.
\end{exercise}

\begin{exercise}
Recompute \eqref{eq:implementation-example} and verify its coefficients.
\end{exercise}

\begin{summarybox}
Reliable transseries software separates syntax, canonical Hahn data, and finite truncation. Monomials require recursive canonical comparison; series require ordered sparse supports and explicit cutoffs. Algebra uses geometric expansions and sparse convolution, while composition and lazy fixed points demand support and contraction certificates. Resurgent numerics add sector arrays, Borel transforms, Pad\'e continuation, and contour integration.
\end{summarybox}


\chapter{A practical workflow and the most common traps}\label{ch:workflow}

Transseries calculations become manageable when performed in a disciplined order. Most mistakes come from doing a correct local manipulation in the wrong ambient scale, or from silently changing between formal and analytic meanings.

\section{The end-to-end workflow}

Given an equation and an asymptotic limit:

\begin{enumerate}
 \item \textbf{Specify the limit and sector.} State whether $x\to+\infty$, $z\to0$ in a complex sector, or $\eps\to0^+$ with other variables fixed.
 \item \textbf{List candidate scales.} Include powers, logarithms, exponentials, and any scales suggested by homogeneous or linearized equations.
 \item \textbf{Find dominant balances.} Compare terms before expanding anything deeply.
 \item \textbf{Choose a leading branch.} Algebraic equations and nonlinear ODEs usually have several.
 \item \textbf{Compute the perturbative sector.} Derive a recurrence and inspect coefficient growth.
 \item \textbf{Linearize about it.} Integrate the linearized exponent to discover actions and power prefactors.
 \item \textbf{Introduce one parameter per independent mode.} Let nonlinear products generate higher sectors.
 \item \textbf{Derive sector equations.} Use index convolution rather than expanding manually sector by sector.
 \item \textbf{Check support admissibility.} Verify a ratio grid or well-based support and identify resonances.
 \item \textbf{Determine analytic status.} Prove or invoke Gevrey bounds, Borel continuation, and summability if actual functions are claimed.
 \item \textbf{Compute Stokes geometry.} Locate action rays and define lateral sums.
 \item \textbf{Match global data.} Use initial conditions, boundary conditions, reality, monodromy, or pole-free requirements to fix parameters.
 \item \textbf{Validate.} Substitute truncated results, estimate residuals, and compare with high-precision numerical solutions.
\end{enumerate}

Each stage answers a different question. Skipping directly from a guessed ansatz to numerical summation often hides a wrong branch or missing scale.

\section{Trap 1: confusing asymptotic equality with equality}

The statement
\[
 f(x)\asym\sum_{n\ge0}a_nx^{-n}
\]
does not mean that the infinite sum converges to $f(x)$. It means a sequence of remainder estimates after finite truncation. Write equality only inside a formal field or after naming a summation operation.

Good notation is:
\begin{align*}
 \trans f&=\sum a_nx^{-n} &&\text{formal object},\\
 f&\asym\sum a_nx^{-n} &&\text{asymptotic relation},\\
 f&=\lateralsum{\theta}\trans f &&\text{summation theorem}.
\end{align*}

\section{Trap 2: believing a zero power series means a zero function}

The function $e^{-x}$ has zero expansion in powers of $x^{-1}$. A boundary condition may therefore be encoded entirely in an exponentially small sector. Before declaring uniqueness, identify the kernel of the linearized operator and the width of the analytic sector.

\section{Trap 3: expanding the large part of an exponential}

For
\[
 e^{x+1/x},
\]
do not write a Taylor series in $x+1/x$. Keep $e^x$ as a monomial and expand only $e^{1/x}$:
\[
 e^{x+1/x}=e^x\left(1+x^{-1}+\frac12x^{-2}+\cdots\right).
\]
The canonical split $T=L+c+S$ exists precisely to enforce this rule.

\section{Trap 4: using a geometric series without a small ratio}

The identity
\[
 (1+S)^{-1}=1-S+S^2-\cdots
\]
is valid formally because $S\precdom1$, not because the symbol visually resembles a correction. If $S=x$, powers grow and the support runs in the wrong direction.

Always factor by the leading monomial first.

\section{Trap 5: treating support as decoration}

A generalized sum with no largest term can make multiplication undefined. For example,
\[
 1+x+x^2+x^3+\cdots
\]
is not a transseries at $x\to+\infty$ in the usual Hahn orientation. Its support climbs forever.

Support conditions are the analogue of convergence conditions for formal algebra. They guarantee that each product coefficient receives finitely many contributions.

\section{Trap 6: assuming every decreasing family is a finite grid}

A reverse-well-ordered support can require infinitely many independent ratios. Algorithms designed for $\N^r$ may silently omit terms. State whether a result is grid-based or merely well-based.

\section{Trap 7: forgetting the power prefactor}

Linearization may give
\[
 \eta'+\left(A-\frac\beta x+\cdots\right)\eta=0.
\]
Integrating yields
\[
 \eta\asym e^{-Ax}x^\beta(1+\cdots),
\]
not just $e^{-Ax}$. The logarithmic integral of the $x^{-1}$ term creates the power.

\section{Trap 8: mistaking a free parameter for a Stokes constant}

The parameter $\sigma$ labels solutions. A Stokes constant $S$ belongs to the equation and chosen normalization. Crossing a ray changes $\sigma$ by a map involving $S$.

Initial data determine $\sigma$; analytic continuation of the equation determines $S$.

\section{Trap 9: counting generated sectors as new constants}

A first-order nonlinear ODE can generate sectors $e^{-kx}$ for every $k\ge1$ while still having only one free parameter. Higher sector coefficients are polynomially generated from that parameter.

Parameter count follows the dimension of the solution space, not the number of sector labels.

\section{Trap 10: missing resonance}

If two action combinations coincide, sector recurrences may become singular. A vanished denominator often signals a required factor of $x$, $\log x$, or a polynomial in parameters. Do not regularize it away numerically before checking resonance.

\section{Trap 11: assuming factorial divergence proves resurgence}

Growth like $n!$ guarantees only local convergence of the first Borel transform. Resurgence additionally requires controlled analytic continuation. A Borel transform may have a natural boundary or uncontrolled singular set.

\section{Trap 12: reading every Pad\'e pole as real}

Finite Pad\'e approximants contain spurious pole-zero pairs. A trustworthy Borel singularity persists under changes of order and precision and fits the predicted action geometry. Branch cuts appear as strings of poles, not necessarily as isolated true poles.

\section{Trap 13: ignoring summation direction}

A Borel sum is directional. The same formal series can have different lateral sums across a singular ray. Always attach $\theta$, $\theta^+$, $\theta^-$, or a median prescription.

\section{Trap 14: using one asymptotic expansion globally}

An expansion valid in one complex sector may fail after rotation because an exponential becomes large. Cover the plane by sectors and use connection formulas. ``As $|z|\to\infty$'' is incomplete without angular restrictions.

\section{Trap 15: interchanging infinite sums without uniform control}

In
\[
 \sum_k\tau^k\sum_na_{k,n}z^n,
\]
Borel summing every $n$-series does not justify summing over $k$. Transasymptotic reorganization, sector convergence, and uniform estimates are separate issues.

\section{Trap 16: Taylor expanding composition under the wrong smallness condition}

The increment $U$ in $F(S+U)$ must be small relative to the variation scale of $F$. For rapidly growing $F$, even $U\to0$ may produce a large exponent change. Examine $F'(S)U/F(S)$ and higher derivative scales.

\section{Trap 17: ignoring branches}

Powers, logarithms, Lambert $W$, and action integrals are multivalued in the complex plane. Record branch cuts and continuation paths. Two formulas with different branches may be different sectors of one global solution rather than contradictions.

\section{Trap 18: trusting a formal residual without ordering it}

After substituting a truncation, do not report merely that the residual is ``small.'' Compute its leading monomial. A successful $N$-term approximation should push the residual below the declared cutoff in the ambient dominance order.

\section{Trap 19: truncating by printed term count}

Ten terms can mean ten power coefficients, ten monomials across several sectors, or ten grid lattice points. Cancellations and mixed scales make printed length a poor precision measure. Use a monomial or sector/order cutoff.

\section{Trap 20: overclaiming analytic meaning}

A formal transseries solution is already valuable. There is no need to call it a convergent function unless a theorem supports that statement. Precise modesty is stronger mathematics:
\begin{quote}
``This is a formal solution in the grid-based LE field; in the stated non-Stokes sector, the rank-one summability theorem provides an analytic realization.''
\end{quote}

\section{A complete miniature example}

Consider
\begin{equation}\label{eq:workflow-example}
 \eps y'+y=x,
 \qquad y(0)=0,
 \qquad \eps\to0^+.
\end{equation}

\subsection*{Step 1: outer balance}
Set $\eps=0$:
\[
 y_0=x.
\]
A regular expansion gives
\[
 y_{\mathrm{out}}=x-\eps,
\]
with no further power corrections because differentiating $x-\eps$ gives $1$ exactly.

\subsection*{Step 2: detect the boundary mismatch}
At $x=0$,
\[
 y_{\mathrm{out}}(0)=-\eps,
\]
so the expansion satisfies the boundary condition only at leading order, not exactly.

\subsection*{Step 3: homogeneous scale}
The homogeneous equation
\[
 \eps h'+h=0
\]
gives
\[
 h=Ce^{-x/\eps}.
\]

\subsection*{Step 4: impose the boundary condition}
The transseries ansatz is
\[
 y=x-\eps+Ce^{-x/\eps}.
\]
At $x=0$, $C=\eps$, so
\begin{equation}\label{eq:workflow-example-solution}
 y=x-\eps+\eps e^{-x/\eps}.
\end{equation}

\subsection*{Step 5: verify}
Direct substitution gives zero residual and $y(0)=0$. For every fixed $x>0$, the exponential is beyond all algebraic orders in $\eps$. In the inner variable $X=x/\eps$ it becomes order $\eps$:
\[
 y=\eps\bigl(X-1+e^{-X}\bigr).
\]
One formula connects the outer and boundary-layer descriptions.

This elementary problem already displays dominant reduction, an invisible homogeneous mode, a transseries parameter fixed by boundary data, and nonuniformity repaired by a stretched variable.

\section{A diagnostic decision tree}

\begin{description}
 \item[Power coefficients grow geometrically:] the series may converge; use ordinary analytic methods first.
 \item[Power coefficients grow like $n!$:] try an order-one Borel transform and inspect its nearest singularities.
 \item[A boundary condition is missing:] solve the homogeneous or linearized problem for exponential scales.
 \item[Nonlinearity is present:] expect powers and combinations of fundamental actions.
 \item[A recurrence denominator vanishes:] test for resonance and logarithmic sectors.
 \item[Many sectors become comparable:] introduce a transasymptotic variable $\tau$.
 \item[Repeated logarithmic depth grows without bound:] move beyond a fixed LE level toward omega-series or hyperseries.
 \item[Persistent oscillation appears:] leave the real ordered $H$-field setting and use complex phases or another asymptotic framework.
\end{description}

\section{Exercises}

\begin{exercise}
For $\eps y'+(1+x)y=1$ with $y(0)=0$, find the outer expansion and the leading boundary-layer exponential.
\end{exercise}

\begin{exercise}
Construct an example in which $U\to0$ but $e^{F(x+U)}/e^{F(x)}$ does not tend to $1$.
\end{exercise}

\begin{exercise}
A coefficient ratio satisfies $a_n/a_{n-1}\sim(n+3/2)/A$. What action and branch exponent does this suggest?
\end{exercise}

\begin{exercise}
Write a checklist distinguishing what must be shown for a formal solution, an asymptotic realization, and a Borel-summed solution.
\end{exercise}

\begin{summarybox}
A successful transseries calculation proceeds from limit and scale selection through dominant balance, perturbative recursion, linearized actions, sector equations, support checks, and only then analytic reconstruction. Most errors are category errors: confusing equality with asymptotics, parameters with Stokes constants, support with notation, or directional summability with a global function.
\end{summarybox}


\chapter{A guided reading map}\label{ch:reading-map}

The literature uses several dialects. This final chapter of the main text explains what to read for each purpose and how the sources fit together.

\section{For a first formal introduction}

Edgar's \emph{Transseries for Beginners} develops grid-based logarithmic-exponential transseries with unusually concrete examples, ratio sets, derivatives, and fixed-point methods \cite{edgar-beginners}. It is the closest source in spirit to the formal middle of this tutorial.

Edgar's paper on composition, recursion, and convergence treats composition and fixed-point phenomena in greater depth \cite{edgar-composition}. His paper on fractional iteration is a natural continuation of \cref{ch:fractional-iteration} \cite{edgar-fractional}.

\section{For formal differential algebra}

Van der Hoeven's \emph{Transseries and Real Differential Algebra} develops a broad computational and differential-algebraic theory, including operators, asymptotic equations, and algorithmic viewpoints \cite{vdhoeven-book}. It is more advanced and denser than an introductory text, but it supplies much of the structural backbone.

The book by Aschenbrenner, van den Dries, and van der Hoeven gives the definitive asymptotic differential algebra and model theory of $\T$ \cite{adh-book}. Their ICM survey \emph{On Numbers, Germs, and Transseries} is a shorter conceptual bridge among Hardy fields, transseries, and surreals \cite{adh-numbers}.

\section{For Borel summation and resurgence}

Costin's \emph{Asymptotics and Borel Summability} develops Borel methods with rigorous differential-equation applications \cite{costin-book}. His rank-one papers establish the correspondence between generalized Borel-summed transseries and actual nonlinear ODE solutions \cite{costin-1995,costin-1998,costin-borel-stokes}.

Dorigoni's review is a compact modern introduction to resurgence, alien calculus, and physical examples \cite{dorigoni}. The long primer by Aniceto, Ba\c{s}ar, and Schiappa develops resurgent transseries, large-order relations, and many worked examples from mathematics and physics \cite{abs-primer}.

\`Ecalle's original works are foundational and far more extensive. They introduce resurgent functions, alien calculus, acceleration, analyzable functions, and the machinery used in the proof of Dulac's conjecture \cite{ecalle-resurgent,ecalle-analyzable}.

\section{For Stokes phenomena and exponential improvement}

Classical asymptotics texts by Dingle, Berry, Howls, Olver, and Olde Daalhuis explain optimal truncation, Stokes multipliers, hyperasymptotics, and terminants. The supplied note by Olde Daalhuis gives a particularly direct route from transseries to Stokes phenomena \cite{olde-daalhuis}. Costin--Kruskal treats optimal truncation and Berry smoothing for nonlinear systems \cite{costin-kruskal}.

\section{For the ordered exponential-field constructions}

The papers of van den Dries, Macintyre, and Marker construct logarithmic-exponential series and connect them with the model theory of the real exponential field \cite{vdDMM-1997,vdDMM-2001}. Dahn--G\"oring is an early source for exponential-logarithmic terms \cite{dahn-goring}. Kuhlmann--Tressl compares LE and EL constructions directly \cite{kuhlmann-tressl}. Schmeling's thesis provides a broad axiomatic framework \cite{schmeling}.

\section{For surreal numbers and omega-series}

Conway's \emph{On Numbers and Games} and Gonshor's monograph develop surreal arithmetic and normal forms \cite{conway,gonshor}. Berarducci--Mantova constructs derivations and transseries structure on the surreals \cite{berarducci-mantova-derivation}. Their later paper develops omega-series, surreal composition, and the interpretation of transseries as surreal functions \cite{berarducci-mantova-germs}.

The survey by Mantova and Matusinski is a gentle bridge among surreal derivations, Hardy fields, and transseries \cite{mantova-matusinski}. More recent work on logarithmic hyperseries and hyperserial fields continues beyond finite LE depth \cite{log-hyperseries,hyperserial-field,surreal-hyperseries}.

\section{For recent Taylor and composition results}

Formal Taylor expansion over generalized power-series fields has received renewed attention. Bagayoko and Mantova give broad existence and convergence results for Taylor expansions in exponential differential fields of generalized series \cite{bagayoko-mantova-taylor}. Mantova proves monotonicity of composition and a Taylor approximation theorem for omega-series with a maximal validity radius \cite{mantova-monotonicity}.

These papers are appropriate after the reader is comfortable with \cref{ch:composition,ch:surreals}; they clarify exactly when a familiar Taylor heuristic is a theorem in large transseries fields.

\section{Suggested routes through this book}

\begin{description}
 \item[Asymptotics route:] \cref{ch:scales,ch:euler,ch:first-transseries,ch:dominant-balance,ch:linear-odes,ch:nonlinear-odes,ch:borel,ch:stokes,ch:resurgence}.
 \item[Formal algebra route:] \cref{ch:monomials,ch:hahn,ch:grids,ch:construction,ch:arithmetic,ch:explog,ch:differentiation,ch:composition}.
 \item[Differential-equation route:] \cref{ch:euler,ch:differentiation,ch:integration,ch:dominant-balance,ch:linear-odes,ch:nonlinear-odes,ch:singular-perturbation,ch:analytic-status}.
 \item[Logic and surreal route:] \cref{ch:hahn,ch:formal-meaning,ch:hfields,ch:surreals,ch:beyond-LE,ch:model-theory}.
 \item[Computational route:] \cref{ch:grids,ch:arithmetic,ch:fixed-points,ch:dominant-balance,ch:computation,ch:workflow}.
\end{description}

\section{What a reader should now be able to do}

After working through the tutorial, a reader should be able to:
\begin{enumerate}
 \item distinguish formal, asymptotic, Borel-summed, and resurgent claims;
 \item order transmonomials and verify simple support conditions;
 \item compute with large/constant/small decompositions;
 \item differentiate, integrate, invert, exponentiate, and compose basic transseries;
 \item derive perturbative and exponential sectors of elementary ODEs;
 \item recognize factorial divergence and perform a basic Borel--Laplace reconstruction;
 \item explain lateral sums, Stokes parameter changes, and large-order relations;
 \item understand how Hardy fields and surreal numbers supply two different meanings for transseries;
 \item identify when a problem requires a larger field, resonance terms, or complex oscillatory scales;
 \item design a finite truncated representation suitable for symbolic computation.
\end{enumerate}

The subject extends far beyond this list, but these skills are enough to read introductory research papers without treating every infinite exponential expression as magic.

\begin{summarybox}
The formal, analytic, logical, and surreal literatures emphasize different structures. Edgar and van der Hoeven lead into formal computation; Costin and \`Ecalle into summation and resurgence; Aschenbrenner--van den Dries--van der Hoeven into $H$-fields and model theory; Berarducci--Mantova into surreal evaluation and omega-series. The best next source depends on which meaning of transseries the reader needs.
\end{summarybox}

\appendix

\chapter{A compact prerequisite toolkit}\label{app:toolkit}

This appendix collects facts used repeatedly in the main text. It is not a replacement for courses in real analysis, complex analysis, or differential equations, but it makes the tutorial more self-contained.

\section{Asymptotic notation}

Let $g(x)>0$ eventually as $x\to\infty$.

\begin{align*}
 f(x)=O(g(x))
 &\quad\Longleftrightarrow\quad
 |f(x)|\le Cg(x)\text{ eventually for some }C,\\
 f(x)=o(g(x))
 &\quad\Longleftrightarrow\quad
 \frac{f(x)}{g(x)}\longrightarrow0,\\
 f(x)\sim g(x)
 &\quad\Longleftrightarrow\quad
 \frac{f(x)}{g(x)}\longrightarrow1.
\end{align*}

The relation $f\sim g$ is much stronger than $f=O(g)$, but neither means $f=g$.

If
\[
 f(x)\asym\sum_{n\ge0}a_n\phi_n(x)
\]
with $\phi_{n+1}=o(\phi_n)$, then for every fixed $N$,
\[
 f(x)-\sum_{n=0}^{N-1}a_n\phi_n(x)
 =o(\phi_{N-1}(x)).
\]
An equivalent common convention uses $O(\phi_N)$ when the scale has additional regularity. Always inspect the stated remainder convention.

\section{Uniqueness of coefficients in an asymptotic scale}

Suppose
\[
 \sum_{n=0}^{N}c_n\phi_n=o(\phi_N)
\]
and $\phi_{n+1}=o(\phi_n)$. Divide by $\phi_0$ and let $x\to\infty$ to obtain $c_0=0$. Repeat after removing the first term. Thus all $c_n=0$.

Consequently, the coefficients of a Poincar\'e asymptotic expansion in a fixed scale are unique, even though the function realizing them need not be.

\section{Geometric and binomial series}

For an ordinary number $|u|<1$,
\begin{align*}
 \frac1{1-u}&=\sum_{n\ge0}u^n,\\
 \log(1+u)&=\sum_{n\ge1}\frac{(-1)^{n+1}}n u^n,\\
 e^u&=\sum_{n\ge0}\frac{u^n}{n!},\\
 (1+u)^\alpha&=\sum_{n\ge0}\binom\alpha n u^n,
\end{align*}
where
\[
 \binom\alpha n
 =\frac{\alpha(\alpha-1)\cdots(\alpha-n+1)}{n!}.
\]

For transseries, replace $|u|<1$ by $u\precdom1$. The support of $u^n$ moves toward smaller monomials, so each retained monomial receives only finitely many contributions.

\section{Formal coefficient extraction}

Write $[z^n]F(z)$ for the coefficient of $z^n$. Then
\begin{align*}
 [z^n](FG)&=\sum_{k=0}^n[z^k]F\,[z^{n-k}]G,\\
 [z^n]F'&=(n+1)[z^{n+1}]F.
\end{align*}
These identities produce most perturbative recurrences after substitution into an ODE.

For a sector variable $\sigma$, use
\[
 [\sigma^k](FG)=\sum_{i+j=k}[\sigma^i]F[\sigma^j]G.
\]
This is why nonlinear sector labels add.

\section{The gamma and beta integrals}

The gamma function is
\begin{equation}\label{eq:gamma-integral-app}
 \Gamma(s)=\int_0^\infty e^{-t}t^{s-1}\dd t,
 \qquad \Re s>0.
\end{equation}
Integration by parts gives
\[
 \Gamma(s+1)=s\Gamma(s),
 \qquad
 \Gamma(n+1)=n!.
\]
Scaling $t=x\xi$ yields
\[
 \int_0^\infty e^{-x\xi}\xi^{s-1}\dd\xi
 =\Gamma(s)x^{-s},
 \qquad\Re x>0.
\]
This is the continuous form of the Borel--Laplace moment identity.

The beta integral is
\begin{equation}\label{eq:beta-integral-app}
 \int_0^1t^{a-1}(1-t)^{b-1}\dd t
 =\frac{\Gamma(a)\Gamma(b)}{\Gamma(a+b)}.
\end{equation}
After $\eta=\xi t$, it proves the convolution rule for Borel transforms.

\section{Stirling's approximation}

As $n\to\infty$,
\begin{equation}\label{eq:stirling-app}
 n!\sim\sqrt{2\pi n}\left(\frac ne\right)^n.
\end{equation}
For a term
\[
 t_n=\frac{n!}{x^{n+1}},
\]
the ratio is $(n+1)/x$, so the least term occurs near $n=x$. Stirling then gives
\[
 t_n\big|_{n\approx x}
 \asymp x^{-1/2}e^{-x}
\]
up to constant factors. Thus the optimally truncated remainder has the scale of the first missing exponential.

\section{First-order linear ODEs}

For
\begin{equation}\label{eq:first-order-app}
 y'+p(x)y=q(x),
\end{equation}
define the integrating factor
\[
 \mu(x)=\exp\left(\int p(x)\dd x\right).
\]
Then
\[
 (\mu y)'=\mu q,
\]
so
\begin{equation}\label{eq:first-order-app-solution}
 y=\mu^{-1}\left(C+\int\mu q\dd x\right).
\end{equation}
The homogeneous scale $\mu^{-1}$ is exactly the beyond-all-orders sector that a power expansion can miss.

\section{Riccati reduction}

A second-order linear equation
\[
 y''=Q(x)y
\]
can be reduced by
\[
 R=\frac{y'}y.
\]
Then
\[
 R'+R^2=Q.
\]
A formal expansion for $R$ can be integrated to obtain
\[
 y=\exp\left(\int R\dd x\right).
\]
The two signs of the leading square root $R\sim\pm\sqrt Q$ produce the two WKB exponentials.

\section{Leading WKB terms}

Take
\[
 y''=Q(x)y
\]
with slowly varying positive $Q$. Write
\[
 y=e^{S},
 \qquad
 S''+(S')^2=Q.
\]
At leading order, $S'\sim\pm\sqrt Q$, hence
\[
 S\sim\pm\int\sqrt Q\dd x.
\]
Keeping the first derivative correction gives
\[
 y\asym Q^{-1/4}
 \exp\left(\pm\int^x\sqrt{Q(t)}\dd t\right).
\]
The factor $Q^{-1/4}$ follows by substituting $S'=\pm\sqrt Q+r$ and solving the first linear correction
\[
 2(\pm\sqrt Q)r\sim-\frac{Q'}{2\sqrt Q}.
\]

\section{Complex sectors and branches}

A sector at the origin is
\[
 S(\theta_0,\alpha,R)
 =\{0<|z|<R:\ |\arg z-\theta_0|<\alpha/2\}.
\]
Its opening is $\alpha$. A branch of $\log z$ on the sector determines
\[
 z^\beta=e^{\beta\log z}.
\]
If the sector winds around the origin, no single-valued logarithm exists without passing to a covering surface.

For asymptotics at infinity, replace $0<|z|<R$ by $|z|>R$.

\section{Cauchy's coefficient formula}

If $g$ is analytic near $0$,
\[
 [\xi^n]g(\xi)
 =\frac1{2\pi\ii}\oint
 \frac{g(\xi)}{\xi^{n+1}}\dd\xi.
\]
Moving the contour outward makes nearby singularities control large $n$. This is the analytic basis of Darboux's principle and resurgent large-order relations.

\section{Residues and lateral jumps}

If a Laplace path crosses a simple pole of $g$ at $A$, the difference between paths bypassing it on opposite sides is a small closed contour. Thus
\[
 \int_{+}e^{-\xi/z}g(\xi)\dd\xi
 -\int_{-}e^{-\xi/z}g(\xi)\dd\xi
 =\pm2\pi\ii e^{-A/z}\Res_{\xi=A}g(\xi).
\]
The sign depends on path orientation. This one-line residue calculation is the simplest Stokes jump.

\section{Well orders and reverse well orders}

A totally ordered set is well ordered if every nonempty subset has a least element. It is reverse well ordered if every nonempty subset has a greatest element.

A support of transmonomials is reverse well ordered because terms are written from largest to smaller. This rules out an infinite strictly increasing sequence
\[
 \m_0\precdom\m_1\precdom\m_2\precdom\cdots
\]
inside the support.

\section{Why Hahn products have finite coefficients}

Let $A$ and $B$ be reverse-well-ordered subsets of an ordered abelian group. Fix a monomial $\m$. Suppose infinitely many pairs $(a_i,b_i)\in A\times B$ satisfy $a_ib_i=\m$. Choose the $a_i$ in strictly decreasing order after passing to a subsequence. Then
\[
 b_i=\m a_i^{-1}
\]
forms a strictly increasing sequence in $B$, contradicting reverse well-order. Hence only finitely many pairs contribute.

This is the core local-finiteness fact behind Hahn multiplication.

\section{Contraction in dominance order}

A map $\Phi$ is contracting near a candidate fixed point if
\[
 U-V\preceqdom\m
 \quad\Longrightarrow\quad
 \Phi(U)-\Phi(V)\precdom\m.
\]
Then each iteration determines at least one additional layer of smaller monomials. The sequence
\[
 U_{n+1}=\Phi(U_n)
\]
need not converge numerically, but its coefficients stabilize from large to small. This is the formal topology used by many transseries fixed-point proofs.

\section{Optimal truncation in one formula}

Suppose
\[
 a_n\sim C A^{-n}\Gamma(n+\beta).
\]
The $n$th term of $\sum a_nz^n$ has ratio approximately
\[
 \left|\frac{a_{n+1}z^{n+1}}{a_nz^n}\right|
 \sim\frac{|z|}{|A|}(n+\beta).
\]
The least term is near
\[
 n_*\approx\frac{|A|}{|z|}-\beta.
\]
Its exponential size is approximately $e^{-|A|/|z|}$, modified by angular dependence and powers. This is why the distance to the nearest Borel singularity controls both coefficient growth and optimal error.

\begin{summarybox}
The prerequisite toolkit consists of asymptotic scales, coefficient extraction, gamma and beta integrals, linear ODEs, WKB reduction, complex sectors, residues, and well-ordered support. Each has a direct transseries role: factorials lead to Borel transforms, kernels lead to exponentials, and reverse well orders make infinite algebra finite coefficient by coefficient.
\end{summarybox}


\chapter{Hints and selected solutions}\label{app:solutions}

The exercises in the main text are intended for active calculation. This appendix gives representative solutions and hints, emphasizing methods that recur elsewhere.

\section{Scales and ordinary asymptotics}

\subsection*{Ordering logarithmic and exponential scales}

For every $a>0$ and every $b>0$,
\[
 \log x=o(x^a),
 \qquad
 x^a=o(e^{x^b}).
\]
For the first, apply l'H\^opital's rule to $\log x/x^a$. For the second, take logarithms:
\[
 \log\frac{x^a}{e^{x^b}}=a\log x-x^b\longrightarrow-\infty.
\]
Repeated logarithms give the extended chain
\[
 1\precdom\log_3x\precdom\log_2x\precdom\log x
 \precdom x^a\precdom e^{x^b}.
\]

\subsection*{An invisible function}

For every $N>0$,
\[
 x^Ne^{-x}\longrightarrow0.
\]
Repeated l'H\^opital or comparison of logarithms proves this. Therefore $e^{-x}=o(x^{-N})$ for every $N$, so its complete inverse-power asymptotic expansion is zero.

\section{The Euler equation}

Insert
\[
 \trans y=\sum_{n\ge0}a_nx^{-n-1}
\]
into $y'+y=x^{-1}$. Since
\[
 y'=-\sum_{n\ge0}(n+1)a_nx^{-n-2},
\]
the $x^{-1}$ coefficient gives $a_0=1$, and the coefficient of $x^{-n-2}$ gives
\[
 a_{n+1}=(n+1)a_n.
\]
Thus $a_n=n!$.

The Borel transform is
\[
 \sum_{n\ge0}\frac{n!}{n!}\xi^n=\frac1{1-\xi}.
\]
The pole at $1$ predicts the missing $e^{-x}$ sector.

\section{Hahn supports}

\subsection*{Why \texorpdfstring{$1+x+x^2+\cdots$}{1+x+x^2+...} is wrong at infinity}

The support $\{1,x,x^2,\dots\}$ has no largest element. Hence it is not reverse well ordered under the dominance order for $x\to+\infty$. In contrast,
\[
 1+x^{-1}+x^{-2}+\cdots
\]
has largest monomial $1$ and every subset has a largest exponent closest to zero.

\subsection*{A product}

For
\[
 A=1+x^{-1}+x^{-2}+\cdots,
\]
formal multiplication gives
\[
 A^2=1+2x^{-1}+3x^{-2}+4x^{-3}+\cdots.
\]
The coefficient of $x^{-n}$ receives exactly $n+1$ contributions. Local finiteness holds despite both supports being infinite.

\section{Grids}

Let $\mu_1=x^{-1}$ and $\mu_2=e^{-x}$. The monomial
\[
 x^{-3}e^{-2x}
\]
corresponds to $(3,2)\in\N^2$. Multiplying by $x^{-4}e^{-x}$ adds vectors:
\[
 (3,2)+(4,1)=(7,3),
\]
which represents $x^{-7}e^{-3x}$.

A cutoff keeping sector degree $k\le2$ and power order $n\le4$ is a rectangle in this grid. A dominance cutoff is generally a different shape because every $e^{-x}$ is smaller than every fixed power $x^{-N}$.

\section{Exponentials and logarithms}

To expand
\[
 \log\left(3e^{x^2}x^{-5}(1+2/x+e^{-x})\right),
\]
use the canonical factorization:
\begin{align*}
 \log T
 &=\log3+x^2-5\log x
 +\log(1+2/x+e^{-x})\\
 &=x^2-5\log x+\log3
 +(2/x+e^{-x})
 -\frac12(2/x+e^{-x})^2+\cdots.
\end{align*}
The first mixed terms are
\[
 \log T=x^2-5\log x+\log3
 +2x^{-1}-2x^{-2}+\frac83x^{-3}+\cdots
 +e^{-x}(1-2x^{-1}+4x^{-2}-\cdots)
 -\frac12e^{-2x}+\cdots.
\]
One must state a cutoff before deciding how far each sector is expanded.

\section{Formal integration}

Seek
\[
 \int e^{x^2}\dd x
 \asymp e^{x^2}\sum_{n\ge0}c_nx^{-2n-1}.
\]
Differentiate one term:
\[
 \frac{\dd}{\dd x}
 \left(e^{x^2}c_nx^{-2n-1}\right)
 =e^{x^2}\left(2c_nx^{-2n}- (2n+1)c_nx^{-2n-2}\right).
\]
Matching the constant term gives $2c_0=1$, so $c_0=1/2$. Cancellation at $x^{-2n-2}$ gives
\[
 2c_{n+1}=(2n+1)c_n.
\]
Hence
\[
 c_n=\frac{(2n-1)!!}{2^{n+1}},
\]
and
\[
 \int e^{x^2}\dd x
 \asymp\frac{e^{x^2}}{2x}
 \left(1+\frac1{2x^2}+\frac3{(2x^2)^2}
 +\frac{15}{(2x^2)^3}+\cdots\right).
\]

\section{Lambert \texorpdfstring{$W$}{W}}

For large positive $x$, write
\[
 W(x)=L_1-L_2+u,
 \qquad
 L_1=\log x,
 \quad
 L_2=\log\log x.
\]
The equation $W+\log W=L_1$ becomes
\[
 u+\log\left(1-\frac{L_2-u}{L_1}\right)=0.
\]
Expanding the logarithm and solving recursively gives
\[
 u=\frac{L_2}{L_1}
 +\frac{L_2(-2+L_2)}{2L_1^2}
 +\frac{L_2(6-9L_2+2L_2^2)}{6L_1^3}
 +\cdots.
\]
The large term is $L_1$, the next correction is $-L_2$, and powers of $L_1^{-1}$ carry polynomial coefficients in $L_2$.

\section{Fixed points}

For
\[
 Y=1+tY^2,
\]
write $Y=\sum_{n\ge0}c_nt^n$. Then $c_0=1$ and
\[
 c_n=\sum_{i+j=n-1}c_ic_j.
\]
This gives
\[
 1,1,2,5,14,42,\dots,
\]
the Catalan numbers
\[
 c_n=\frac1{n+1}\binom{2n}{n}.
\]
The exact branch satisfying $Y(0)=1$ is
\[
 Y=\frac{1-\sqrt{1-4t}}{2t}.
\]
Formal iteration stabilizes coefficients because multiplication by $t$ pushes every new error to higher order.

\section{Dominant balance in an exponential cubic}

For
\[
 Y^3-xY-e^x=0,
\]
the exponentially large branch balances $Y^3$ with $e^x$:
\[
 Y\sim e^{x/3}.
\]
Set $Y=e^{x/3}(1+u)$. Then
\[
 (1+u)^3-1-xe^{-2x/3}(1+u)=0.
\]
The small parameter is $q=xe^{-2x/3}$. Solving
\[
 3u+3u^2+u^3=q(1+u)
\]
recursively gives $u=q/3+O(q^2)$.

The other two complex branches use $c=\omega$ and $c=\omega^2$, where
$\omega^3=1$ and $\omega\ne1$. Setting
\[
 Y=c\,e^{x/3}(1+u)
\]
gives
\[
 (1+u)^3-1=cq(1+u),
 \qquad q=xe^{-2x/3},
\]
so that $u=cq/3+O(q^2)$. Thus all three branches have magnitude
$e^{x/3}$; only the branch with $c=1$ is real for large positive $x$.
The lesson is to test every neglected term after proposing a balance.

\section{Airy WKB}

For $y''-xy=0$, take
\[
 y=e^S.
\]
Then
\[
 S''+(S')^2=x.
\]
The leading equation gives $S'\sim\pm x^{1/2}$, hence
\[
 S\sim\pm\frac23x^{3/2}.
\]
The first correction is $S'\sim\pm x^{1/2}-1/(4x)$, so
\[
 y\asym x^{-1/4}e^{\pm2x^{3/2}/3}.
\]
Successive inverse powers of $x^{3/2}$ produce the standard Airy coefficient series.

\section{The nonlinear Riccati recurrence}

From
\[
 a_m=(m-1)a_{m-1}+\sum_{i=1}^{m-1}a_ia_{m-i}
\]
and $a_1=1$:
\begin{align*}
 a_2&=1+1=2,\\
 a_3&=2\cdot2+(1\cdot2+2\cdot1)=8,\\
 a_4&=3\cdot8+(1\cdot8+2\cdot2+8\cdot1)=44,\\
 a_5&=4\cdot44+(1\cdot44+2\cdot8+8\cdot2+44\cdot1)=296.
\end{align*}
The convolution term is the coefficient of $x^{-m}$ in $\phi_0^2$.

\section{A boundary layer}

For
\[
 \eps y'+(1+x)y=1,
 \qquad y(0)=0,
\]
set $y\asym y_0+\eps y_1+\cdots$. Then
\[
 y_0=\frac1{1+x},
 \qquad
 y_1=-\frac{y_0'}{1+x}=\frac1{(1+x)^3}.
\]
The leading outer value at $0$ is $1$, so the boundary mismatch is order one. The homogeneous equation gives
\[
 h=C\exp\left(-\frac1\eps\int_0^x(1+t)\dd t\right)
 =C\exp\left(-\frac{x+x^2/2}{\eps}\right).
\]
At leading order $C=-1$. Near $x=0$, the inner variable $X=x/\eps$ recovers $e^{-X}$.

\section{Borel transforms}

For
\[
 \trans f(x)=\sum_{n\ge0}(-1)^nn!x^{-n-1},
\]
we have
\[
 \Borel\trans f(\xi)=\frac1{1+\xi}.
\]
Its nearest singularity is at $\xi=-1$, so positive-direction Borel summation is unobstructed while the negative direction is singular.

For
\[
 \sum_{n\ge0}A^{-n-1}n!x^{-n-1},
\]
we obtain
\[
 \frac1{A-\xi},
\]
with a singularity at $\xi=A$.

\section{A simple Stokes jump}

If
\[
 g(\xi)=\frac{r}{\xi-A},
\]
then upper-minus-lower lateral contours differ by an oriented loop around $A$. In the small-$z$ convention,
\[
 \operatorname{Disc}\int e^{-\xi/z}g(\xi)\frac{\dd\xi}{z}
 =\pm\frac{2\pi\ii r}{z}e^{-A/z}.
\]
If the Borel convention assigns the pole with an extra factor of $z$ or uses $r/(A-\xi)$, the displayed prefactor and sign change. This is why conventions must be fixed before quoting a Stokes constant.

\section{Branch singularities and large order}

The binomial expansion gives
\[
 (1-w)^{-\beta}
 =\sum_{n\ge0}
 \frac{\Gamma(n+\beta)}{\Gamma(\beta)\Gamma(n+1)}w^n.
\]
With $w=\xi/A$ and $a_n=n![\xi^n]g$, one gets
\[
 a_n=\frac{\Gamma(n+\beta)}{\Gamma(\beta)A^n}.
\]
This exactly realizes the factorial-over-power law rather than merely approximating it.

\section{Hardy-field ordering}

Order
\[
 \log\log x,
 \quad(\log x)^{10},
 \quad e^{\sqrt{\log x}},
 \quad x^{1/100}.
\]
Taking logarithms where helpful:
\[
 \log\log x
 \precdom(\log x)^{10}
 \precdom e^{\sqrt{\log x}}
 \precdom x^{1/100}.
\]
For the last comparison,
\[
 \log\left(\frac{e^{\sqrt{\log x}}}{x^{1/100}}\right)
 =\sqrt{\log x}-\frac1{100}\log x\to-\infty.
\]

The logarithmic derivatives are
\begin{align*}
 (x^{\log x})^\dagger
 &=\frac{\dd}{\dd x}\bigl((\log x)^2\bigr)
 =\frac{2\log x}{x},\\
 \left(e^{x/\log x}\right)^\dagger
 &=\left(\frac{x}{\log x}\right)'
 =\frac{\log x-1}{(\log x)^2}.
\end{align*}

\section{Surreal normal forms}

Using formal Hahn arithmetic,
\begin{align*}
 (\omega+1)^2&=\omega^2+2\omega+1,\\
 (1+\omega^{-1})^{-1}
 &=1-\omega^{-1}+\omega^{-2}-\omega^{-3}+\cdots.
\end{align*}
The latter family is summable because its monomials strictly decrease.

\section{Computational inversion with mixed scales}

Let
\[
 T=1+x^{-1}+e^{-x}.
\]
Set $S=x^{-1}+e^{-x}$. Then
\[
 T^{-1}=1-S+S^2-S^3+\cdots.
\]
Ordering by dominance gives, through algebraic order $x^{-2}$ and exponential order $e^{-x}x^{-2}$,
\begin{align*}
 T^{-1}
 &=1-x^{-1}+x^{-2}+O(x^{-3})\\
 &\quad-e^{-x}\bigl(1-2x^{-1}+3x^{-2}+O(x^{-3})\bigr)
 +e^{-2x}\bigl(1+O(x^{-1})\bigr)+\cdots.
\end{align*}
Whether the $e^{-2x}$ term is retained depends on the precise cutoff: it is smaller than every $e^{-x}x^{-N}$, so a cutoff at $e^{-x}x^{-2}$ normally discards it.

\section{A composition counterexample}

Let
\[
 F(x)=e^{e^x},
 \qquad U(x)=e^{-x/2}.
\]
Then $U\to0$, but
\[
 \frac{F(x+U)}{F(x)}
 =\exp\left(e^x(e^U-1)\right).
\]
Since $e^U-1\sim U$,
\[
 e^x(e^U-1)\sim e^{x/2}\to\infty.
\]
Therefore the ratio tends to infinity, not to $1$. Absolute smallness of the increment is insufficient for Taylor-small composition.

\begin{summarybox}
The selected solutions repeatedly use the same moves: compare logarithms, factor the leading monomial, derive coefficient convolutions, linearize for exponential scales, and inspect neglected terms. These are more important than memorizing any one coefficient list.
\end{summarybox}


\chapter{Glossary of transseries language}\label{app:glossary}

The same word can carry slightly different emphases in formal differential
algebra, exponential asymptotics, resurgence, model theory, and surreal
analysis. This glossary records the meaning intended in this tutorial and,
where useful, points out common variations.

\renewcommand{\arraystretch}{1.16}
\begin{longtable}{@{}p{0.235\textwidth}p{0.715\textwidth}@{}}
\toprule
\textbf{Term} & \textbf{Meaning in this tutorial}\\
\midrule
\endfirsthead
\toprule
\textbf{Term} & \textbf{Meaning in this tutorial}\\
\midrule
\endhead
\midrule
\multicolumn{2}{r}{\small Continued on the next page}\\
\endfoot
\bottomrule
\endlastfoot

\textbf{Action}\index{action}
& A quantity $A$ appearing in an exponential scale $e^{-A/z}$ or $e^{-A x}$.
In differential equations it is normally obtained by integrating an
appropriate eigenvalue of the linearized or WKB problem. Borel singularities
frequently occur at the same values $A$.\\

\textbf{Alien derivative}\index{alien derivative}
& An operator $\alien_A$ that extracts, in a normalized way, the resurgent
singularity located at Borel position $A$. It does not mean an ordinary
derivative with respect to the independent variable.\\

\textbf{Analyzable function}\index{analyzable function}
& In \`Ecalle's theory, roughly, a function represented by a transseries that
can be reconstructed through generalized Borel--Laplace procedures, possibly
including acceleration. Precise definitions require more structure than
formal LE transseries alone.\\

\textbf{Anti-Stokes direction}\index{anti-Stokes direction}
& A direction on which two relevant exponentials have equal magnitude. Some
authors interchange the names ``Stokes'' and ``anti-Stokes,'' so the defining
equation should always be stated.\\

\textbf{Asymptotic expansion}\index{asymptotic expansion}
& A sequence of finite approximation statements relative to a chosen ordered
scale. It is not an assertion that the corresponding infinite sum converges.\\

\textbf{Asymptotic scale}\index{asymptotic scale}
& A sequence $(\phi_n)$ with $\phi_{n+1}=o(\phi_n)$ in the relevant limit.
The order may combine powers, logarithms, exponentials, and iterated
exponentials rather than only powers of one variable.\\

\textbf{Beyond all orders}\index{beyond all orders}
& Smaller than every term in a specified asymptotic scale. For example,
$e^{-x}$ is beyond all orders of $x^{-1}$ as $x\to+\infty$. The phrase is
relative to a scale, not absolute.\\

\textbf{Borel plane}\index{Borel plane}
& The complex $\xi$-plane in which the Borel transform lives. Singularities
there encode factorial growth, Stokes jumps, and often relations among
transseries sectors.\\

\textbf{Borel sum}\index{Borel sum}
& A function obtained by Borel transforming a divergent formal series,
analytically continuing the transform, and applying a Laplace integral along
a specified direction. Direction and normalization are part of the data.\\

\textbf{Borel transform}\index{Borel transform}
& A transform that removes factorial growth, typically sending
$\sum_{n\ge0}a_n z^{n+1}$ to $\sum_{n\ge0}a_n\xi^n/n!$. Several shifted
normalizations are standard.\\

\textbf{Bridge equation}\index{bridge equation}
& A relation connecting alien derivatives to ordinary variation of
transseries parameters, schematically
$\alien_A\Phi=C_A(\sigma)\partial_\sigma\Phi$. It turns Borel singularity data
into transformations of the full transseries.\\

\textbf{Coefficient stabilization}\index{coefficient stabilization}
& The phenomenon that successive formal iterates agree through progressively
smaller monomials. This is the correct replacement for numerical convergence
in many transseries fixed-point constructions.\\

\textbf{Constant part}\index{constant part}
& In the decomposition $T=L+c+S$, the real constant $c$. The terms $L$ and
$S$ are respectively purely large and purely small.\\

\textbf{Depth}\index{logarithmic depth}
& A measure of how many iterated logarithms are needed to express a
transseries after an appropriate change of variable. Conventions vary among
constructions; it should not be confused with exponential height.\\

\textbf{Dominance}\index{dominance}
& The order of asymptotic size. We write $f\precdom g$ when $f/g\to0$ in a
functional realization, or when the corresponding ordered monomials have that
formal relation.\\

\textbf{Dominant balance}\index{dominant balance}
& A leading-order cancellation among the largest candidate terms of an
equation. Every proposed balance must be checked by substituting its implied
scale into all neglected terms.\\

\textbf{Écalle acceleration}\index{acceleration}
& A generalized summation step used when ordinary Borel--Laplace summation has
insufficient growth control. It changes the effective time scale in the
Borel plane and permits treatment of more rapidly divergent levels.\\

\textbf{EL series}\index{EL series}
& ``Exponential-logarithmic'' series constructed by a particular ordering of
field, exponential, logarithmic, and Hahn-series extensions. EL and LE
constructions are closely related but are not identical by definition.\\

\textbf{Exponential height}\index{exponential height}
& The number of genuinely large exponential-extension stages needed to build
a monomial or transseries. Expanding a small exponential in an already
constructed field need not increase height.\\

\textbf{Formal equality}\index{formal equality}
& Equality of coefficients attached to every monomial in the chosen
transseries field. It does not by itself assert convergence, summability, or
equality of analytic functions.\\

\textbf{Gevrey order}\index{Gevrey order}
& A quantitative factorial-growth class. In a common convention, Gevrey--1
coefficients satisfy $|a_n|\le C K^n n!$. Different authors index Gevrey
classes differently, so the bound is safer than the label alone.\\

\textbf{Grid}\index{grid}
& A finitely generated multiplicative family of monomials, usually generated
by small ratio monomials. Grid-based transseries are especially convenient
for algorithms but form a proper subclass of all well-based transseries.\\

\textbf{Hahn field}\index{Hahn field}
& A field of generalized series $\sum a_\m\m$ with coefficients in a field,
monomials in an ordered abelian group, and suitably well-ordered support.
Transseries fields are built from iterated Hahn-field ideas plus exponential,
logarithmic, and differential structure.\\

\textbf{Hardy field}\index{Hardy field}
& A field of germs at $+\infty$ closed under differentiation. Ordering is by
eventual sign. It provides a functional interpretation of asymptotic
comparison but does not naturally contain persistent oscillation such as
$\sin x$.\\

\textbf{$H$-field}\index{H-field@$H$-field}
& An ordered differential field satisfying axioms abstracting the behavior of
Hardy fields and transseries, including compatibility between order and
derivation for infinitely large elements.\\

\textbf{Hyperasymptotics}\index{hyperasymptotics}
& A hierarchy of exponentially improved approximations obtained by expanding
the remainder after optimal truncation, then repeating the process. It is an
analytic approximation method rather than merely a larger formal support.\\

\textbf{Hyperseries}\index{hyperseries}
& Generalized transseries allowing levels beyond the standard finite
logarithmic-exponential hierarchy, designed to accommodate transfinite
iteration of exponential and logarithmic behavior. Exact definitions depend
on the chosen hyperserial framework.\\

\textbf{Instanton sector}\index{instanton sector}
& Physics terminology for a component multiplied by an exponential such as
$e^{-kA/z}$. In this tutorial it is synonymous with an exponential
transseries sector; no physical instanton interpretation is required.\\

\textbf{Lateral sum}\index{lateral Borel sum}
& A Borel--Laplace sum taken just above or just below a singular direction.
The two lateral sums may differ by exponentially small terms.\\

\textbf{Large part}\index{large part}
& The sum $L$ of monomials larger than $1$ in the canonical decomposition
$T=L+c+S$. ``Purely large'' means that all nonzero monomials are larger than
$1$.\\

\textbf{LE series}\index{LE series}
& ``Logarithmic-exponential'' generalized power series obtained by iterating
appropriate exponential and logarithmic extensions. The standard field
$\T$ is one canonical LE-type construction.\\

\textbf{Logarithmic derivative}\index{logarithmic derivative}
& The quotient $f^\dagger=f'/f$. It converts products into sums and is often
the fastest way to compare derivatives of transmonomials or to derive WKB
corrections.\\

\textbf{Median summation}\index{median summation}
& A balanced prescription, built from lateral sums and part of the Stokes
automorphism, that can preserve reality on a Stokes direction when the
problem has the appropriate conjugation symmetry.\\

\textbf{Monomial}\index{monomial}
& A multiplicative unit of asymptotic size, such as $x^a$, $e^{-x}x^b$, or
$\exp(-e^x+x^2)$. A transseries is a well-organized formal linear combination
of such monomials.\\

\textbf{Newton polygon}\index{Newton polygon}
& A geometric device for detecting balances from exponents or valuations of
terms. In transseries problems it may be generalized to ordered monomial
values rather than only integer powers.\\

\textbf{Normal form}\index{normal form}
& A canonical representation in which like monomials are collected and the
support is ordered. Normalization is essential for reliable equality testing,
truncation, and coefficient extraction.\\

\textbf{Omega-series}\index{omega-series}
& A large field of surreal expressions generated from real numbers and the
simplest infinite surreal $\omega$ using summation, exponential, logarithm,
and composition subject to summability conditions. It extends the ordinary
LE transseries viewpoint.\\

\textbf{Optimal truncation}\index{optimal truncation}
& Stopping a divergent asymptotic series near its least term. The resulting
error is often exponentially small and reveals the scale of the next
transseries sector.\\

\textbf{Perturbative sector}\index{perturbative sector}
& The sector with no exponentially small factor, usually denoted $\Phi_0$.
It is often an ordinary formal power series and may be factorially divergent.\\

\textbf{Ratio set}\index{ratio set}
& A finite set of small monomials whose nonnegative products generate a grid.
Ratio sets provide a concrete calculus for witnessing and manipulating
relative smallness.\\

\textbf{Resonance}\index{resonance}
& An integer relation among actions, exponents, or linearized eigenvalues that
causes sectors to collide. Resonance may force powers of logarithms, mixed
parameters, or modified recurrences.\\

\textbf{Resurgence}\index{resurgence}
& The principle, made precise through endlessly continuable Borel transforms
and singularity calculus, that different asymptotic sectors encode one
another. Factorial large-order growth is one visible consequence.\\

\textbf{Sector}\index{transseries sector}
& Either a region of the complex independent-variable plane or a component
of a transseries carrying a fixed exponential weight. Context normally makes
the two meanings clear; ``angular sector'' and ``instanton sector'' remove the
ambiguity.\\

\textbf{Small part}\index{small part}
& The sum $S$ of monomials smaller than $1$ in $T=L+c+S$. Formal Taylor,
geometric, logarithmic, and binomial expansions are valid around $S=0$ when
the resulting support is summable.\\

\textbf{Stokes automorphism}\index{Stokes automorphism}
& The formal transformation relating the two lateral resummations across a
singular direction. In resurgent notation it is built from exponentials of
alien derivatives, with convention-dependent normalization.\\

\textbf{Stokes constant}\index{Stokes constant}
& A numerical invariant measuring a normalized Borel singularity or the
associated jump between sectors. It is not the same thing as an arbitrary
integration constant or transseries parameter.\\

\textbf{Stokes direction}\index{Stokes direction}
& A direction along which a Borel singularity lies on the Laplace ray, or,
in WKB language, a curve/direction where the relevant exponential phases have
a specified alignment. Naming conventions vary, so the defining condition is
primary.\\

\textbf{Support}\index{support}
& The set of monomials with nonzero coefficients in a generalized series.
Its order-theoretic properties determine whether infinite addition,
multiplication, composition, and differentiation are legitimate.\\

\textbf{Summable family}\index{summable family}
& A family of generalized series whose union of supports is suitably
well-based and for which each monomial occurs in only finitely many members.
This is the formal condition that allows the family to be added.\\

\textbf{Surreal number}\index{surreal number}
& An element of Conway's proper class $\mathbf{No}$. Every surreal has a Hahn-
like normal form $\sum r_i\omega^{a_i}$. Suitable transseries can be evaluated
at infinite surreal arguments, but this is different from Borel summation.\\

\textbf{Transasymptotics}\index{transasymptotics}
& Reorganization of a transseries when infinitely many ordinary sectors
become comparable, often by keeping a composite variable such as
$\tau=\sigma e^{-A/z}z^\beta$ unexpanded and summing its dependence first.\\

\textbf{Transmonomial}\index{transmonomial}
& Another name for a transseries monomial, emphasizing that its scale may
involve nested exponentials and logarithms rather than a single power.\\

\textbf{Transseries}\index{transseries}
& A formal generalized series on an ordered family of transmonomials, equipped
with enough support control to permit algebra and often calculus. In analytic
asymptotics the word also denotes a sectorial formal object intended for
Borel--Laplace reconstruction. The surrounding structure must be specified.\\

\textbf{Transseries parameter}\index{transseries parameter}
& A formal constant such as $\sigma$ multiplying a homogeneous exponential
sector. Boundary or initial data select its value; Stokes continuation may
change its coordinate description.\\

\textbf{Well-based series}\index{well-based series}
& A generalized series whose support is reverse well ordered when monomials
are listed from largest to smaller. Well-based supports are more general than
finitely generated grids.\\

\textbf{Well ordering}\index{well ordering}
& The order condition guaranteeing an extremal element in every nonempty
subset. Hahn-series authors often use well-ordered support with the valuation
order; this tutorial uses reverse well ordering with the dominance order.
The two conventions encode the same idea after reversing the order.\\

\textbf{WKB transseries}\index{WKB transseries}
& A transseries built from exponential phases $\exp(\pm\int\sqrt Q)$,
algebraic prefactors, and divergent correction series for a differential
equation with a large parameter or asymptotic independent variable.\\

\end{longtable}
\renewcommand{\arraystretch}{1}

\begin{summarybox}
When terminology is uncertain, ask four questions: What is the independent
limit? What is the ordered monomial group? What support condition is imposed?
What analytic reconstruction, if any, is claimed? Those questions usually
resolve apparent contradictions between different transseries dialects.
\end{summarybox}


\chapter{Notation and a cross-dialect concordance}\label{app:notation}

\section{Core notation used in this book}

\renewcommand{\arraystretch}{1.16}
\begin{longtable}{@{}p{0.23\textwidth}p{0.70\textwidth}@{}}
\toprule
\textbf{Notation} & \textbf{Meaning}\\
\midrule
\endfirsthead
\toprule
\textbf{Notation} & \textbf{Meaning}\\
\midrule
\endhead
\midrule
\multicolumn{2}{r}{\small Continued on the next page}\\
\endfoot
\bottomrule
\endlastfoot

$x\to+\infty$ & Default large-variable limit in the formal LE chapters.\\
$z\to0$ & Default small-parameter limit in the Borel and resurgence chapters.\\
$f\precdom g$ & $f$ is asymptotically smaller than $g$; functionally,
$f/g\to0$.\\
$f\succdom g$ & $f$ is asymptotically larger than $g$.\\
$f\equivdom g$ & $f$ and $g$ have comparable magnitude; neither ratio is
forced to tend to zero.\\
$f\sim g$ & Usually asymptotic equivalence, $f/g\to1$. In a few cited sources
$\sim$ may denote comparability, which is why this tutorial uses
$\equivdom$ for the latter.\\
$f\asym\sum a_n\phi_n$ & $f$ has the displayed asymptotic expansion. The
symbol does not mean convergence of the infinite sum.\\
$\M$ or $\G$ & An ordered multiplicative group of monomials.\\
$\m,\n$ & Individual monomials.\\
$\supp T$ & The support of $T$: all monomials with nonzero coefficient.\\
$\lm(T)$ & Leading, hence largest, monomial of a nonzero transseries $T$.\\
$\lc(T)$ & Coefficient of the leading monomial.\\
$\T$ & The standard logarithmic-exponential transseries field in the formal
chapters. The exact construction is specified in \cref{ch:construction}.\\
$T=L+c+S$ & Canonical decomposition into purely large part, real constant,
and purely small part.\\
$T=a\m(1+S)$ & Leading-monomial factorization, with $a\ne0$ and
$S\precdom1$.\\
$f^\dagger$ & Logarithmic derivative $f'/f$.\\
$\Borel$ & A Borel transform, with normalization stated locally.\\
$\Laplace_\theta$ & Laplace transform along the ray of angle $\theta$.\\
$\lateralsum{\theta^\pm}$ & Upper or lower lateral Borel--Laplace sum along a
singular direction.\\
$\operatorname{Disc}_\theta$ & Difference of the two lateral values across
direction $\theta$, with order stated in context.\\
$\Stokes_\theta$ & Stokes automorphism in direction $\theta$.\\
$\alien_A$ & Alien derivative associated with Borel singularity/action $A$.\\
$A$ & An action, so that an exponential sector contains $e^{-A/z}$ or an
equivalent large-variable scale.\\
$\sigma$ & A transseries parameter, normally indexing a homogeneous sector.\\
$\Phi_k$ & The coefficient series in the $k$th exponential sector.\\
$\Gamma$ & Gamma function; $\Gamma(n+1)=n!$.\\
$\Res$ & Complex residue.\\
$\mathbf{No}$ & The proper class of surreal numbers.\\
$\omega$ & The simplest infinite positive surreal number; also occasionally a
root of unity in local algebraic examples, where this is stated explicitly.\\

\end{longtable}
\renewcommand{\arraystretch}{1}

\section{Large-variable and small-parameter forms}

Many formulas occur in two equivalent-looking conventions. Setting
\[
 x=\frac1z
\]
translates the elementary scales as follows:
\[
 \begin{array}{c@{\qquad\longleftrightarrow\qquad}c}
 x^{-n} & z^n,\\
 e^{-Ax} & e^{-A/z},\\
 x^\beta e^{-Ax} & z^{-\beta}e^{-A/z},\\
 \log x & -\log z.
 \end{array}
\]
The translation is simple algebraically but not harmless analytically:
angular sectors at $z=0$ become sectors at $x=\infty$ with reversed angles,
and branch choices for $\log z$ must be tracked.

\section{Three common Borel normalizations}

Suppose
\[
 \widehat f(z)=\sum_{n\ge0}a_nz^{n+1}.
\]
A convenient normalization in this book is
\[
 \Borel\widehat f(\xi)=\sum_{n\ge0}\frac{a_n}{n!}\xi^n,
 \qquad
 \mathcal S_\theta\widehat f(z)
 =\int_0^{e^{i\theta}\infty}e^{-\xi/z}
   \Borel\widehat f(\xi)\,\dd\xi.
\]
Equally common are:
\[
 \widehat g(z)=\sum_{n\ge0}b_nz^n,
 \qquad
 \Borel_1\widehat g(\xi)
 =b_0\delta(\xi)+\sum_{n\ge1}\frac{b_n}{(n-1)!}\xi^{n-1},
\]
or a Laplace integral with $\dd\xi/z$. These conventions relocate factors of
$z$, signs, and residues. Consequently, a bare formula for a Stokes jump or
alien derivative is incomplete unless its Borel and Laplace normalizations
are known.

\section{Formal, asymptotic, and analytic assertions}

The following statements are logically different:
\begin{align*}
 \mathcal E(\widehat F)&=0
 &&\text{in a formal transseries field},\\
 f(z)&\asym\widehat F(z)
 &&\text{in a specified sector},\\
 f(z)&=\mathcal S_\theta\widehat F(z)
 &&\text{after a specified summation procedure},\\
 f_+(z)&=f_-(z)\circ\Stokes_\theta
 &&\text{as a resurgent continuation relation}.
\end{align*}
A proof of the first line does not automatically prove any later line. A
rigorous theorem must supply hypotheses that connect them.

\section{A dictionary among five viewpoints}

\begin{center}
\small
\begin{tabularx}{\textwidth}{@{}p{0.17\textwidth}X X@{}}
\toprule
\textbf{Viewpoint} & \textbf{Primary object} & \textbf{Question it answers}\\
\midrule
Formal Hahn/LE
& A well-based sum of ordered monomials
& Are algebra, differentiation, composition, and recursive construction
well defined coefficient by coefficient?\\
Asymptotic analysis
& A hierarchy of finite remainder estimates
& How accurately does a finite truncation approximate a function in a limit
or sector?\\
Borel--resurgent
& Analytic continuation and singularities of Borel transforms
& Can a divergent transseries be reconstructed, and how do its sectors jump
and determine large-order growth?\\
Hardy-field
& A germ of a real function at $+\infty$
& Which formal order and differential identities have an eventual functional
realization?\\
Surreal
& A value or function at infinite surreal arguments
& Can the transseries be interpreted inside a vast ordered exponential field
with canonical infinite and infinitesimal numbers?\\
\bottomrule
\end{tabularx}
\end{center}

\section{A minimum data sheet for a transseries claim}

Before trusting or communicating a transseries formula, record:
\begin{enumerate}
 \item the limit and angular domain;
 \item the coefficient field and monomial group;
 \item the support condition (grid-based, well-based, multisummable, and so on);
 \item the normalization of logarithms, powers, Borel transforms, and Laplace
       transforms;
 \item the meaning of equality: formal, asymptotic, summed, germ equality, or
       surreal equality;
 \item the status of free parameters and how boundary data choose them;
 \item singular directions and the convention for Stokes jumps;
 \item the theorem, if any, that realizes the formal object analytically.
\end{enumerate}
This data sheet is deliberately repetitive: most avoidable mistakes in the
subject arise from silently changing one of these items.

\begin{summarybox}
Notation in transseries theory is not fully standardized. The safest practice
is to define dominance, support orientation, Borel normalization, Stokes-jump
orientation, and the intended notion of equality before using compact
formulas. Once those choices are fixed, the different dialects fit together
cleanly.
\end{summarybox}


\backmatter

\begin{thebibliography}{99}
\phantomsection
\addcontentsline{toc}{chapter}{Bibliography}

\bibitem{abs-primer}
I.~Aniceto, G.~Ba\c{s}ar, and R.~Schiappa,
\emph{A primer on resurgent transseries and their asymptotics},
Physics Reports \textbf{809} (2019), 1--135.

\bibitem{adh-book}
M.~Aschenbrenner, L.~van den Dries, and J.~van der Hoeven,
\emph{Asymptotic Differential Algebra and Model Theory of Transseries},
Annals of Mathematics Studies, vol.~195, Princeton University Press, 2017.

\bibitem{adh-numbers}
M.~Aschenbrenner, L.~van den Dries, and J.~van der Hoeven,
\emph{On numbers, germs, and transseries},
in \emph{Proceedings of the International Congress of Mathematicians 2018},
vol.~2, 19--42.

\bibitem{adh-universal-surreal}
M.~Aschenbrenner, L.~van den Dries, and J.~van der Hoeven,
\emph{The surreal numbers as a universal $H$-field},
Journal of the European Mathematical Society \textbf{21} (2019), no.~4,
1179--1199.

\bibitem{bagayoko-mantova-taylor}
V.~Bagayoko and V.~Mantova,
\emph{Taylor expansions over generalised power series},
arXiv:2509.08473 (2025).

\bibitem{hyperserial-field}
V.~Bagayoko and J.~van der Hoeven,
\emph{The hyperserial field of surreal numbers},
arXiv:2310.14873 (2023).

\bibitem{surreal-hyperseries}
V.~Bagayoko and J.~van der Hoeven,
\emph{Surreal numbers as hyperseries},
arXiv:2310.14879 (2023).

\bibitem{balser}
W.~Balser,
\emph{From Divergent Power Series to Analytic Functions: Theory and
Application of Multisummable Power Series},
Lecture Notes in Mathematics, vol.~1582, Springer, 1994.

\bibitem{berarducci-mantova-derivation}
A.~Berarducci and V.~Mantova,
\emph{Surreal numbers, derivations and transseries},
Journal of the European Mathematical Society \textbf{20} (2018), no.~2,
339--390.

\bibitem{berarducci-mantova-germs}
A.~Berarducci and V.~Mantova,
\emph{Transseries as germs of surreal functions},
Transactions of the American Mathematical Society \textbf{371} (2019),
no.~5, 3549--3592.

\bibitem{conway}
J.~H. Conway,
\emph{On Numbers and Games}, 2nd ed., A~K Peters, 2001.

\bibitem{costin-1995}
O.~Costin,
\emph{Exponential asymptotics, transseries, and generalized Borel summation
for analytic nonlinear rank-one systems of ordinary differential equations},
International Mathematics Research Notices (1995), no.~8, 377--417.

\bibitem{costin-1998}
O.~Costin,
\emph{On Borel summation and Stokes phenomena for rank-one nonlinear systems
of ordinary differential equations},
Duke Mathematical Journal \textbf{93} (1998), no.~2, 289--344.

\bibitem{costin-borel-stokes}
O.~Costin,
\emph{On Borel summation and Stokes phenomena of nonlinear differential
systems}, arXiv:math/0608408 (2006).

\bibitem{costin-book}
O.~Costin,
\emph{Asymptotics and Borel Summability},
Chapman \& Hall/CRC Monographs and Surveys in Pure and Applied Mathematics,
vol.~141, 2009.

\bibitem{costin-kruskal}
O.~Costin and M.~D. Kruskal,
\emph{On optimal truncation of divergent series solutions of nonlinear
differential systems; Berry smoothing},
Proceedings of the Royal Society of London A \textbf{455} (1999),
1931--1956.

\bibitem{dahn-goring}
B.~Dahn and P.~G\"oring,
\emph{Notes on exponential-logarithmic terms},
Fundamenta Mathematicae \textbf{127} (1987), no.~1, 45--50.

\bibitem{dingle}
R.~B. Dingle,
\emph{Asymptotic Expansions: Their Derivation and Interpretation},
Academic Press, 1973.

\bibitem{dorigoni}
D.~Dorigoni,
\emph{An introduction to resurgence, trans-series and alien calculus},
Annals of Physics \textbf{409} (2019), article 167914.

\bibitem{ecalle-resurgent}
J.~\`Ecalle,
\emph{Les Fonctions R\'esurgentes}, 3 vols.,
Publications Math\'ematiques d'Orsay, 1981.

\bibitem{ecalle-analyzable}
J.~\`Ecalle,
\emph{Introduction aux fonctions analysables et preuve constructive de la
conjecture de Dulac}, Hermann, 1992.

\bibitem{edgar-beginners}
G.~A. Edgar,
\emph{Transseries for beginners},
Real Analysis Exchange \textbf{35} (2009/2010), no.~2, 253--310;
arXiv:0801.4877.

\bibitem{edgar-composition}
G.~A. Edgar,
\emph{Transseries: composition, recursion, and convergence},
arXiv:0909.1259 (2009).

\bibitem{edgar-fractional}
G.~A. Edgar,
\emph{Fractional iteration of series and transseries},
arXiv:1002.2378 (2010).

\bibitem{gonshor}
H.~Gonshor,
\emph{An Introduction to the Theory of Surreal Numbers},
London Mathematical Society Lecture Note Series, vol.~110,
Cambridge University Press, 1986.

\bibitem{hahn}
H.~Hahn,
\emph{\"Uber die nichtarchimedischen Gr\"o\ss{}ensysteme},
Sitzungsberichte der Kaiserlichen Akademie der Wissenschaften,
Mathematisch-Naturwissenschaftliche Klasse, Abteilung IIa,
\textbf{116} (1907), 601--655.

\bibitem{hardy-orders}
G.~H. Hardy,
\emph{Orders of Infinity: The ``Infinit\"arcalc\"ul'' of Paul du Bois-Reymond},
Cambridge Tracts in Mathematics and Mathematical Physics, no.~12,
Cambridge University Press, 1910.

\bibitem{kuhlmann-tressl}
S.~Kuhlmann and M.~Tressl,
\emph{Comparison of exponential-logarithmic and logarithmic-exponential
series},
Mathematical Logic Quarterly \textbf{58} (2012), no.~6, 434--448.

\bibitem{log-hyperseries}
L.~van den Dries, J.~van der Hoeven, and E.~Kaplan,
\emph{Logarithmic hyperseries},
Transactions of the American Mathematical Society \textbf{372} (2019),
no.~7, 5199--5241.

\bibitem{mantova-matusinski}
V.~Mantova and M.~Matusinski,
\emph{Surreal numbers with derivation, Hardy fields and transseries: a
survey},
in \emph{New Paths in Algebraic Geometry}, Contemporary Mathematics,
vol.~697, American Mathematical Society, 2017, 265--290.

\bibitem{mantova-monotonicity}
V.~Mantova,
\emph{Monotonicity and a Taylor approximation theorem for transseries},
arXiv:2601.07747v2 (8 May 2026).

\bibitem{mitschisauzin}
C.~Mitschi and D.~Sauzin,
\emph{Divergent Series, Summability and Resurgence I: Monodromy and
Resurgence}, Lecture Notes in Mathematics, vol.~2153, Springer, 2016.

\bibitem{olde-daalhuis}
A.~B. Olde Daalhuis,
\emph{Transseries analysis and Stokes phenomena},
unpublished lecture/project note, 2018.

\bibitem{olver}
F.~W. J. Olver,
\emph{Asymptotics and Special Functions},
Academic Press, 1974; reprinted by A~K Peters, 1997.

\bibitem{schmeling}
M.~C. Schmeling,
\emph{Corps de transs\'eries},
Ph.D. thesis, Universit\'e Paris~VII, 2001.

\bibitem{vdDMM-1997}
L.~van den Dries, A.~Macintyre, and D.~Marker,
\emph{Logarithmic-exponential power series},
Journal of the London Mathematical Society \textbf{56} (1997), no.~3,
417--434.

\bibitem{vdDMM-2001}
L.~van den Dries, A.~Macintyre, and D.~Marker,
\emph{Logarithmic-exponential series},
Annals of Pure and Applied Logic \textbf{111} (2001), 61--113.

\bibitem{vdhoeven-book}
J.~van der Hoeven,
\emph{Transseries and Real Differential Algebra},
Lecture Notes in Mathematics, vol.~1888, Springer, 2006.

\end{thebibliography}

\printindex

\end{document}
